% Generated mechanically from frozen input inputs/p08/ja/spaces-properties.tex.
% Do not edit this generated file; edit the bound locale source instead.

% OK, start here.
%

\setcounter{chapter}{65}
\chapter{代数空間の諸性質}



\phantomsection
\label{spaces-properties-section-phantom}


\section{はじめに}
\label{spaces-properties-section-introduction}

\noindent
代数空間の簡単な導入については
Spaces, Section \ref{spaces-section-introduction}
を参照されたい。また、代数空間に関する基本的な定義と規約について、
同章の一部を読まれたい。本章では、代数空間の基本的な概念と性質を
いくつか導入し始める。準分離的な代数空間の場合の基本的な文献は
\cite{Kn} である。

\medskip\noindent
議論がやや不自然になる箇所がある。これは、まず代数空間そのものの
性質を扱い、その後で初めて代数空間の射の性質を扱う、という設計上の
判断をしたためである。この原則について、代数空間の
{\it エタール射}だけは例外とし、Section \ref{spaces-properties-section-etale-morphisms}
で導入する。しかしその節までは、射がある性質をもつと言うとき、
その射の始域がスキームである（あるいは、その射が表現可能である）
ことを自動的に意味する。

\medskip\noindent
本章の内容の一部（とりわけ点に関する内容）は、良好な代数空間を扱う章で
改良される。


\section{規約}
\label{spaces-properties-section-conventions}

\noindent
すべてのスキームは大 fppf サイト $\Sch_{fppf}$ に含まれるものと
常に仮定する。また、考察するすべての環 $A$ は、
$\Spec(A)$ がこの大サイトの対象と（同型で）あるという性質をもつものとする。

\medskip\noindent
$S$ をスキーム、$X$ を $S$ 上の代数空間とする。本章および次章では、
$X$ とそれ自身との積（$S$ 上の代数空間の圏における積）を
$X \times X$ ではなく $X \times_S X$ と書く。これは、
Spaces, Section \ref{spaces-section-change-base-scheme}
のように基底スキームを変更するときの混同を避けるためである。


\section{分離公理}
\label{spaces-properties-section-separation}

\noindent
この節では、代数空間の「絶対的」な分離条件をすべてまとめる。
本書の用語では、任意の代数空間はある定まった基底スキーム上の
代数空間であるから、$S$ 上の $X$ の任意の絶対的性質は、
$X$ を $\Spec(\mathbf{Z})$ 上の代数空間とみなしたときに課される
条件に対応する。正確な定式化は次のとおりである。

\begin{definition}
\label{spaces-properties-definition-separated}
（Spaces, Definition \ref{spaces-definition-separated} と比較されたい。）
大 fppf サイト
$\Sch_{fppf} = (\Sch/\Spec(\mathbf{Z}))_{fppf}$
を考える。$X$ を
$\Spec(\mathbf{Z})$ 上の代数空間とする。
$\Delta : X \to X \times X$ を対角射とする。
\begin{enumerate}
\item $\Delta$ が閉埋め込みであるとき、$X$ は {\it 分離的}であるという。
\item $\Delta$ が埋め込みであるとき、$X$ は
{\it 局所分離的}\footnote{文献では、これはしばしば準分離的かつ
局所分離的な代数空間を指す。}であるという。
\item $\Delta$ が準コンパクトであるとき、$X$ は {\it 準分離的}であるという。
\item 各 $X_i$ が準分離的であるような Zariski 被覆
$X = \bigcup_{i \in I} X_i$（Spaces,
Definition \ref{spaces-definition-Zariski-open-covering} を参照）が存在するとき、
$X$ は {\it Zariski 局所準分離的}\footnote{
この概念は B.\ Conrad によって提案された。}であるという。
\end{enumerate}
$S$ を $\Sch_{fppf}$ に含まれるスキームとし、
$X$ を $S$ 上の代数空間とする。このとき、$X$ を
$\Spec(\mathbf{Z})$ 上の代数空間とみなしたもの（Spaces,
Definition \ref{spaces-definition-base-change} を参照）が対応する性質を
もつ場合に、$X$ は {\it 分離的}、{\it 局所分離的}、{\it 準分離的}、
または {\it Zariski 局所準分離的}であるという。
\end{definition}

\noindent
$S$ 上の代数空間 $X$ が（上の絶対的な意味で）分離的ならば、
$S$ 上分離的でもある（上記の他の絶対的分離性質についても同様である）。
これは
Morphisms of Spaces, Section \ref{spaces-morphisms-section-separation-axioms}
で詳しく論じる。
Lemma \ref{spaces-properties-lemma-quasi-separated-quasi-compact-pieces}
において、Zariski 局部分離的であることは基底スキームに依存しない
（したがって絶対的概念と同値である）ことを見る。

\begin{lemma}
\label{spaces-properties-lemma-trivial-implications}
$S$ をスキームとし、$X$ を $S$ 上の代数空間とする。
Definition \ref{spaces-properties-definition-separated} の分離公理の間には
次の含意がある。
\begin{enumerate}
\item 分離的ならば、他のすべての性質をもつ。
\item 準分離的ならば、Zariski 局所準分離的である。
\end{enumerate}
\end{lemma}

\begin{proof}
省略する。
\end{proof}

\begin{lemma}
\label{spaces-properties-lemma-characterize-quasi-separated}
$S$ をスキームとし、$X$ を $S$ 上の代数空間とする。
次は同値である。
\begin{enumerate}
\item $X$ は準分離的な代数空間である。
\item $U$, $V$ が準コンパクトなスキームであるような
$U \to X$, $V \to X$ に対して、ファイバー積
$U \times_X V$ は準コンパクトである。
\item $U$, $V$ がアフィンであるような
$U \to X$, $V \to X$ に対して、ファイバー積
$U \times_X V$ は準コンパクトである。
\end{enumerate}
\end{lemma}

\begin{proof}
Spaces, Lemma
\ref{spaces-lemma-category-of-spaces-over-smaller-base-scheme}
を用いれば、$S = \Spec(\mathbf{Z})$ と仮定してよい。
$U \times_X V = X \times_{X \times X} (U \times V)$ であり、
$U$, $V$ が準コンパクトならば $U \times V$ も準コンパクトであるから、
(1) は (2) を含意する。(2) が (3) を含意することは明らかである。
(3) を仮定する。スキーム $W$ と全射エタール射
$W \to X$ を選ぶ。このとき $W \times W \to X \times X$ は
全射エタールである。したがって、Spaces, Lemma
\ref{spaces-lemma-descent-representable-transformations-property}
により、
$$
j : W \times_X W = X \times_{(X \times X)} (W \times W) \to W \times W
$$
が準コンパクトであることを示せば十分である。
$U \subset W$ と $V \subset W$ がアフィン開部分ならば、
仮定により $j^{-1}(U \times V) = U \times_X V$ は準コンパクトである。
アフィン開部分 $U \times V$ は $W \times W$ のアフィン開被覆をなすので
（Schemes, Lemma \ref{schemes-lemma-affine-covering-fibre-product}）、
Schemes, Lemma \ref{schemes-lemma-quasi-compact-affine}
から結論が従う。
\end{proof}

\begin{lemma}
\label{spaces-properties-lemma-characterize-separated}
$S$ をスキームとし、$X$ を $S$ 上の代数空間とする。
次は同値である。
\begin{enumerate}
\item $X$ は分離的な代数空間である。
\item $U$, $V$ がアフィンであるような $U \to X$, $V \to X$ に対して、
ファイバー積 $U \times_X V$ はアフィンであり、
$$
\mathcal{O}(U) \otimes_\mathbf{Z} \mathcal{O}(V)
\longrightarrow
\mathcal{O}(U \times_X V)
$$
は全射である。
\end{enumerate}
\end{lemma}

\begin{proof}
Spaces, Lemma
\ref{spaces-lemma-category-of-spaces-over-smaller-base-scheme}
を用いれば、$S = \Spec(\mathbf{Z})$ と仮定してよい。
$U \times_X V = X \times_{X \times X} (U \times V)$ であり、
$U$, $V$ がアフィンならば $U \times V$ もアフィンであるから、
(1) は (2) を含意する。(2) を仮定する。スキーム $W$ と
全射エタール射 $W \to X$ を選ぶ。このとき
$W \times W \to X \times X$ は全射エタールである。
したがって、Spaces, Lemma
\ref{spaces-lemma-descent-representable-transformations-property}
により、
$$
j : W \times_X W = X \times_{(X \times X)} (W \times W) \to W \times W
$$
が閉埋め込みであることを示せば十分である。
$U \subset W$ と $V \subset W$ がアフィン開部分ならば、
仮定により $j^{-1}(U \times V) = U \times_X V$ はアフィンであり、
対応する環準同型が全射であるから、
$U \times_X V \to U \times V$ は閉埋め込みである。
アフィン開部分 $U \times V$ は $W \times W$ のアフィン開被覆をなすので
（Schemes, Lemma \ref{schemes-lemma-affine-covering-fibre-product}）、
Morphisms, Lemma \ref{morphisms-lemma-closed-immersion}
から結論が従う。
\end{proof}







\section{代数空間の点}
\label{spaces-properties-section-points}

\noindent
Spaces, Example \ref{spaces-example-affine-line-translation}
から明らかなように、代数空間の点を体のスペクトルからの単射として
定義すべきではない。代わりに、
Schemes, Section \ref{schemes-section-points}
で説明されているのとまったく同様に、体のスペクトルの射の同値類として
定義する。

\medskip\noindent
$S$ をスキームとする。$F$ を $(\Sch/S)_{fppf}$ 上の前層とする。
$K$ を体とする。射
$$
\Spec(K) \longrightarrow F.
$$
を考える。Yoneda の補題により、これは元
$p \in F(\Spec(K))$ によって与えられる。このような二つの組
$(\Spec(K), p)$ と $(\Spec(L), q)$ が {\it 同値}であるとは、
第三の体 $\Omega$ と可換図式
$$
\xymatrix{
\Spec(\Omega) \ar[r] \ar[d] &
\Spec(L) \ar[d]^q \\
\Spec(K) \ar[r]^p &
F.
}
$$
が存在することをいう。言い換えれば、体拡大
$K \to \Omega$ と $L \to \Omega$ が存在し、$p$ と $q$ が
$F(\Spec(\Omega))$ の同じ元に写ることをいう。これが同値関係を定めることの
検証は省略する。

\begin{definition}
\label{spaces-properties-definition-points}
$S$ をスキームとし、$X$ を $S$ 上の代数空間とする。
$X$ の {\it 点}とは、体のスペクトルから $X$ への射の同値類である。
$X$ の点の集合を $|X|$ と表す。
\end{definition}

\noindent
$f : X \to Y$ が $S$ 上の代数空間の射ならば、誘導される写像
$|f| : |X| \to |Y|$ があり、代表元
$x : \Spec(K) \to X$ を代表元
$f \circ x : \Spec(K) \to Y$ に写すことに注意する。

\begin{lemma}
\label{spaces-properties-lemma-scheme-points}
$S$ をスキームとし、$X$ を $S$ 上のスキームとする。
スキームとしての $X$ の点と、代数空間としての $X$ の点との間には
標準的な一対一対応がある。
\end{lemma}

\begin{proof}
これは Schemes, Lemma \ref{schemes-lemma-characterize-points} である。
\end{proof}

\begin{lemma}
\label{spaces-properties-lemma-points-cartesian}
$S$ をスキームとする。
$$
\xymatrix{
Z \times_Y X \ar[r] \ar[d] & X \ar[d] \\
Z \ar[r] & Y
}
$$
を $S$ 上の代数空間の Cartesian 図式とする。このとき、点集合の写像
$$
|Z \times_Y X|
\longrightarrow
|Z| \times_{|Y|} |X|
$$
は全射である。
\end{lemma}

\begin{proof}
実際、体 $K$, $L$ と射
$\Spec(K) \to X$, $\Spec(L) \to Z$ が与えられたとする。
これらが $|Y|$ の元として一致するという仮定は、共通拡大
$M/K$ と $M/L$ が存在して、
$\Spec(M) \to \Spec(K) \to X \to Y$ と
$\Spec(M) \to \Spec(L) \to Z \to Y$ が一致することを意味する。
これはちょうど、射 $\Spec(M) \to Z \times_Y X$ が得られるための条件である。
\end{proof}

\begin{lemma}
\label{spaces-properties-lemma-characterize-surjective}
$S$ をスキームとする。$X$ を $S$ 上の代数空間とする。
$f : T \to X$ をスキームから $X$ への射とする。次は同値である。
\begin{enumerate}
\item $f : T \to X$ は全射である（Spaces, Definition
\ref{spaces-definition-relative-representable-property} の意味で）。
\item $|f| : |T| \to |X|$ は全射である。
\end{enumerate}
\end{lemma}

\begin{proof}
(1) を仮定する。$x : \Spec(K) \to X$ を体のスペクトルから
$X$ への射とする。仮定により、スキームの射
$\Spec(K) \times_X T \to \Spec(K)$ は全射である。
したがって、体拡大 $K'/K$ と射
$\Spec(K') \to \Spec(K) \times_X T$ が存在して、図式
$$
\xymatrix{
\Spec(K') \ar[r] \ar[d] &
\Spec(K) \times_X T \ar[d] \ar[r] &
T \ar[d] \\
\Spec(K) \ar@{=}[r] &
\Spec(K) \ar[r]^-x & X
}
$$
の左の正方形は可換である。これは
$|f| : |T| \to |X|$ が全射であることを示す。

\medskip\noindent
(2) を仮定する。$Z$ をスキームとし、$Z \to X$ を射とする。
スキームの射 $Z \times_X T \to T$ が全射であること、
すなわち $|Z \times_X T| \to |Z|$ が全射であることを示さなければならない。
これは (2) と Lemma \ref{spaces-properties-lemma-points-cartesian} から従う。
\end{proof}

\begin{lemma}
\label{spaces-properties-lemma-points-presentation}
$S$ をスキームとする。$X$ を $S$ 上の代数空間とする。
$X = U/R$ を $X$ の表示とする。Spaces, Definition
\ref{spaces-definition-presentation} を参照されたい。
このとき $|R| \to |U| \times |U|$ の像は同値関係であり、
$|X|$ はこの同値関係による $|U|$ の商である。
\end{lemma}

\begin{proof}
仮定は、$U$ がスキームであり、$p : U \to X$ が全射エタール射であり、
$R = U \times_X U$ がスキームであって $U$ 上のエタール同値関係を定め、
層として $X = U/R$ であることを意味する。
Lemma \ref{spaces-properties-lemma-characterize-surjective} により、
$|U| \to |X|$ は全射である。Lemma \ref{spaces-properties-lemma-points-cartesian}
により、写像
$$
|R| \longrightarrow |U| \times_{|X|} |U|
$$
は全射である。したがって、$|R| \to |U| \times |U|$ の像は、
$u_1$ と $u_2$ の $|X|$ における像が一致する組
$(u_1, u_2) \in |U| \times |U|$ の集合にちょうど等しい。
これら二つの主張を合わせれば補題の結論を得る。
\end{proof}

\begin{lemma}
\label{spaces-properties-lemma-topology-points}
$S$ をスキームとする。$S$ 上の代数空間の点集合には、
次の性質をもつ位相が一意に存在する。
\begin{enumerate}
\item $X$ が $S$ 上のスキームならば、$|X|$ の位相は通常の位相である
（Lemma \ref{spaces-properties-lemma-scheme-points} の同一視を介して）。
\item $S$ 上の代数空間の任意の射 $X \to Y$ に対して、
写像 $|X| \to |Y|$ は連続である。
\item $U$ がスキームである任意のエタール射 $U \to X$ に対して、
位相空間の写像 $|U| \to |X|$ は連続かつ開である。
\end{enumerate}
\end{lemma}

\begin{proof}
$X$ を $S$ 上の代数空間とする。$U$ を $S$ 上のスキームとし、
$p : U \to X$ を全射エタール射とする。
$W \subset |X|$ が開であることを、$|p|^{-1}(W)$ が $|U|$ の
開部分集合であることとして定義する。これは $|X|$ 上の位相である
（$|X|$ 上の商位相である。Topology, Lemma
\ref{topology-lemma-quotient} を参照）。

\medskip\noindent
この位相が表示の選択に依存しないことを示そう。そのためには、
$U'$ がスキームで、$U' \to X$ がエタール射ならば、写像
$|U'| \to |X|$（ここで $|X|$ の位相は上の $U \to X$ を用いて
定義したもの）が開かつ連続であることを示せば十分である。
これは同時に (3) も証明する。
$U'' = U \times_X U'$ とおくと、可換図式
$$
\xymatrix{
U'' \ar[r] \ar[d] & U' \ar[d] \\
U \ar[r] & X
}
$$
を得る。$U \to X$ と $U' \to X$ はエタールであるから、
$U'' \to U$ と $U'' \to U'$ はともにスキームのエタール射である。
さらに、$U'' \to U'$ は全射である。したがって集合の写像の可換図式
$$
\xymatrix{
|U''| \ar[r] \ar[d] & |U'| \ar[d] \\
|U| \ar[r] & |X|
}
$$
を得る。下の水平矢印は全射であり
（Lemma \ref{spaces-properties-lemma-characterize-surjective} または
Lemma \ref{spaces-properties-lemma-points-presentation} を参照）、
$|X|$ の位相の定義により連続である。上の水平矢印は
Morphisms, Lemma \ref{morphisms-lemma-etale-open}
により全射、連続、かつ開である。左の垂直矢印も連続かつ開である
（再び Morphisms, Lemma \ref{morphisms-lemma-etale-open}）。
したがって形式的に、右の垂直矢印は連続かつ開である。

\medskip\noindent
証明を終えるため (2) を示す。
$a : X \to Y$ を代数空間の射とする。Spaces, Lemma
\ref{spaces-lemma-lift-morphism-presentations}
により、図式
$$
\xymatrix{
U \ar[d]_p \ar[r]_\alpha & V \ar[d]^q \\
X \ar[r]^a & Y
}
$$
で、$U$, $V$ はスキーム、$p$, $q$ は全射エタールであるものを取れる。
これから図式
$$
\xymatrix{
|U| \ar[d]_p \ar[r]_\alpha & |V| \ar[d]^q \\
|X| \ar[r]^a & |Y|
}
$$
が得られる。下の水平矢印以外はすべて連続であることが分かっており、
二つの垂直矢印は全射かつ開である。したがって、望みどおり下の水平矢印も
連続である。
\end{proof}

\begin{definition}
\label{spaces-properties-definition-topological-space}
$S$ をスキームとし、$X$ を $S$ 上の代数空間とする。
$X$ の台 {\it 位相空間}とは、Lemma \ref{spaces-properties-lemma-topology-points}
で構成した位相を備えた点集合 $|X|$ のことである。
\end{definition}

\noindent
この位相空間は、Spaces, Definition
\ref{spaces-definition-small-Zariski-site} の小 Zariski サイト
$X_{Zar}$ と同じ情報をもつことが分かる。

\begin{lemma}
\label{spaces-properties-lemma-open-subspaces}
$S$ をスキームとし、$X$ を $S$ 上の代数空間とする。
\begin{enumerate}
\item 対応 $X' \mapsto |X'|$ は、$X$ の開部分空間 $X'$
（Spaces, Definition \ref{spaces-definition-immersion} を参照）と
位相空間 $|X|$ の開部分との間の、包含関係を保つ全単射を定める。
\item $X$ の開部分空間の族
$\{X_i \subset X\}_{i \in I}$ が Zariski 被覆である
（Spaces, Definition \ref{spaces-definition-Zariski-open-covering} を参照）
ための必要十分条件は、$|X| = \bigcup |X_i|$ であることである。
\end{enumerate}
言い換えれば、$X$ の小 Zariski サイト $X_{Zar}$ は、位相空間
$|X|$ に付随するサイトと標準的に同一視される
（Sites, Example \ref{sites-example-site-topological} を参照）。
\end{lemma}

\begin{proof}
(1) を証明するため、この対応の逆を構成しよう。
$W \subset |X|$ を開とする。全射エタール射
$p : U \to X$ とエタール射 $s, t : R \to U$ に対応する表示
$X = U/R$ を選ぶ。構成により $|p|^{-1}(W)$ は $U$ の開部分である。
対応する開部分スキームを $W' \subset U$ と書く。
$R' = s^{-1}(W') = t^{-1}(W')$ は明らかに $R$ の Zariski 開部分であり、
$W'$ 上のエタール同値関係を定める。Spaces, Lemma
\ref{spaces-lemma-finding-opens} により、射
$X' = W'/R' \to X$ は開埋め込みである。したがって Spaces, Lemma
\ref{spaces-lemma-representable-over-space} により $X'$ は代数空間である。
構成により $|X'| = W$、すなわち $X'$ は $W$ に対応する $X$ の部分空間である。
これで (1) が証明された。

\medskip\noindent
(2) を証明するため、$\{X_i \subset X\}_{i \in I}$ が
開部分空間の族ならば、それが Zariski 被覆であるための必要十分条件は
$U = \bigcup U \times_X X_i$ が開被覆であることに注意する。
これは Zariski 被覆の定義と、射 $U \to X$ が
$(\Sch/S)_{fppf}$ 上の前層の写像として全射であることから従う。
一方、Lemma \ref{spaces-properties-lemma-points-presentation} により、
$|X| = \bigcup |X_i|$ であるための必要十分条件も
$U = \bigcup U \times_X X_i$ であることである
（射影 $U \times_X X_i \to X_i$ が全射エタールであることも用いる）。
したがって (2) の同値性が従う。
\end{proof}

\begin{lemma}
\label{spaces-properties-lemma-factor-through-open-subspace}
$S$ をスキームとする。$X$, $Y$ を $S$ 上の代数空間とする。
$X' \subset X$ を開部分空間とする。
$f : Y \to X$ を $S$ 上の代数空間の射とする。
このとき、$f$ が $X'$ を経由するための必要十分条件は、
$|f| : |Y| \to |X|$ が $|X'| \subset |X|$ を経由することである。
\end{lemma}

\begin{proof}
Spaces, Lemma \ref{spaces-lemma-base-change-immersions}
により、$Y' = Y \times_X X' \to Y$ は開埋め込みである。
$|f|(|Y|) \subset |X'|$ ならば、明らかに $|Y'| = |Y|$ である。
したがって Lemma \ref{spaces-properties-lemma-open-subspaces} により $Y' = Y$ である。
\end{proof}

\begin{lemma}
\label{spaces-properties-lemma-etale-image-open}
$S$ をスキームとし、$X$ を $S$ 上の代数空間とする。
$U$ をスキームとし、$f : U \to X$ をエタール射とする。
Lemma \ref{spaces-properties-lemma-open-subspaces} によって開部分
$|f|(|U|) \subset |X|$ に対応する開部分空間を $X' \subset X$ とする。
このとき $f$ は、全射エタール射 $f' : U \to X'$ を経由する。
さらに、$R = U \times_X U$ ならば $R = U \times_{X'} U$ であり、
$X'$ は表示 $X' = U/R$ をもつ。
\end{lemma}

\begin{proof}
分解の存在は Lemma \ref{spaces-properties-lemma-factor-through-open-subspace} から従う。
Lemma \ref{spaces-properties-lemma-characterize-surjective} により $f'$ は全射である。
$f'$ がエタールであることを見るため、$T$ をスキームとし、
$T \to X'$ を射とする。$X' \to X$ は層の単射であるから、
$T \times_X U = T \times_{X'} U$ である。したがって、
$f$ がエタールであるという仮定から、射影
$T \times_{X'} U \to T$ はエタールである。
$X' \to X$ は単射なので
$U \times_X U = U \times_{X'} U$ でもある。
Spaces, Lemma \ref{spaces-lemma-space-presentation}
から $X' = U/R$ が従う。
\end{proof}

\begin{lemma}
\label{spaces-properties-lemma-equivalence-class-point-monomorphism}
$S$ をスキームとし、$X$ を $S$ 上の代数空間とする。
$K$, $L$ を体とし、
$p : \Spec(K) \to X$ と $q : \Spec(L) \to X$ を射とする。
$p$ と $q$ が $|X|$ の同じ点を定め、$p$ が単射であると仮定する。
このとき $q$ は $p$ を一意に経由する。
\end{lemma}

\begin{proof}
$p$ と $q$ は $|X|$ の同じ点を定めるので、スキーム
$$
Y = \Spec(K) \times_{p, X, q} \Spec(L)
$$
は空でない。単射の基底変換は単射であるから、射影
$Y \to \Spec(L)$ は単射である。したがって
$Y = \Spec(L)$ である。Schemes, Lemma
\ref{schemes-lemma-mono-towards-spec-field} を参照されたい。
よって $q$ は $p$ を経由する。一意性は $p$ が単射であることから従う。
\end{proof}

\begin{lemma}
\label{spaces-properties-lemma-points-monomorphism}
$S$ をスキームとし、$X$ を $S$ 上の代数空間とする。写像
$$
\{\Spec(k) \to X \text{ は単射であり }k\text{ は体}\}
\longrightarrow
|X|
$$
を考える。この写像は単射である。
\end{lemma}

\begin{proof}
これは Lemma \ref{spaces-properties-lemma-equivalence-class-point-monomorphism} から従う。
\end{proof}

\noindent
Decent Spaces,
Lemma \ref{decent-spaces-lemma-decent-points-monomorphism}
において、$X$ が良好であるとき
Lemma \ref{spaces-properties-lemma-points-monomorphism} の写像は全単射であることを見る。
















\section{準コンパクトな空間}
\label{spaces-properties-section-quasi-compact}

\begin{definition}
\label{spaces-properties-definition-quasi-compact}
$S$ をスキームとし、$X$ を $S$ 上の代数空間とする。
$U$ が準コンパクトなスキームであるような全射エタール射
$U \to X$ が存在するとき、$X$ は {\it 準コンパクト}であるという。
\end{definition}

\begin{lemma}
\label{spaces-properties-lemma-quasi-compact-space}
$S$ をスキームとし、$X$ を $S$ 上の代数空間とする。
$X$ が準コンパクトであるための必要十分条件は、
$|X|$ が準コンパクトであることである。
\end{lemma}

\begin{proof}
スキーム $U$ と全射エタール射 $U \to X$ を選ぶ。
Lemma \ref{spaces-properties-lemma-characterize-surjective} を用いる。
$U$ が準コンパクトならば、$|U| \to |X|$ は全射であるから、
$|X|$ は準コンパクトである。
$|X|$ が準コンパクトならば、$|U| \to |X|$ は開であるから、
$|U'| \to |X|$ が全射であるような準コンパクトな開部分
$U' \subset U$ が存在する（この射は依然としてエタールである）。
これで結論を得る。
\end{proof}

\begin{lemma}
\label{spaces-properties-lemma-finite-disjoint-quasi-compact}
準コンパクトな代数空間の有限非交和は、準コンパクトな代数空間である。
\end{lemma}

\begin{proof}
これは Lemma \ref{spaces-properties-lemma-quasi-compact-space} と対応する位相的事実から明らかである。
\end{proof}

\begin{example}
\label{spaces-properties-example-quasi-compact-not-very-reasonable}
空間 $\mathbf{A}^1_{\mathbf{Q}}/\mathbf{Z}$ は準コンパクトな代数空間である。
\end{example}

\begin{lemma}
\label{spaces-properties-lemma-space-locally-quasi-compact}
$S$ をスキームとし、$X$ を $S$ 上の代数空間とする。
$|X|$ の各点は、準コンパクトな開近傍からなる基本近傍系をもつ。
とくに $|X|$ は、Topology, Definition
\ref{topology-definition-locally-quasi-compact}
の意味で局所準コンパクトである。
\end{lemma}

\begin{proof}
これは、スキーム $U$ と位相空間の全射、開、連続な写像
$U \to |X|$ が存在するという事実から形式的に従う。
もう少し正確に言えば、$u \in U$ が $x \in |X|$ に写るならば、
$u$ のアフィン近傍の像が $x$ の準コンパクトな開近傍の基本近傍系を与える。
\end{proof}







\section{特殊な被覆}
\label{spaces-properties-section-special-coverings}

\noindent
この節では、代数空間の全射エタール被覆の存在に関する
いくつかの直接的な補題をまとめる。

\begin{lemma}
\label{spaces-properties-lemma-cover-by-union-affines}
$S$ をスキームとし、$X$ を $S$ 上の代数空間とする。
$U$ がアフィンスキームの非交和であるような全射エタール射
$U \to X$ が存在する。さらに、これらの各アフィンが
$S$ のアフィン開部分へ写ると仮定してよい。
\end{lemma}

\begin{proof}
$V \to X$ を全射エタール射とする。各 $V_i$ が $S$ の
アフィン開部分へ写るような Zariski 開被覆
$V = \bigcup_{i \in I} V_i$ を取る。
$U = \coprod_{i \in I} V_i$ とおき、誘導される射
$U \to V \to X$ を考える。これは、関手の表現可能な
エタール全射変換の合成としてエタールかつ全射である
（一般原理 Spaces, Lemma
\ref{spaces-lemma-composition-representable-transformations-property}
ならびに Morphisms, Lemmas
\ref{morphisms-lemma-composition-surjective} および
\ref{morphisms-lemma-composition-etale} による）。
\end{proof}

\begin{lemma}
\label{spaces-properties-lemma-union-of-quasi-compact}
$S$ をスキームとし、$X$ を $S$ 上の代数空間とする。
各代数空間 $X_i$ がアフィンスキームによる全射エタール被覆をもつような
Zariski 被覆 $X = \bigcup X_i$ が存在する。
さらに、各 $X_i$ が $S$ のアフィン開部分へ写ると仮定してよい。
\end{lemma}

\begin{proof}
Lemma \ref{spaces-properties-lemma-cover-by-union-affines} により、
$U_i$ はアフィンで $S$ のアフィン開部分へ写るような全射エタール射
$U = \coprod U_i \to X$ を取れる。$U_i \to X$ が全射エタール射
$U_i \to X_i$ を経由するような $X$ の開部分空間を
$X_i \subset X$ とする。Lemma \ref{spaces-properties-lemma-etale-image-open}
を参照されたい。$U = \bigcup U_i$ なので $X = \bigcup X_i$ である。
$U_i \to X_i$ は全射であるから、$X_i \to S$ は $S$ の
アフィン開部分へ写る。
\end{proof}

\begin{lemma}
\label{spaces-properties-lemma-quasi-compact-affine-cover}
$S$ をスキームとし、$X$ を $S$ 上の代数空間とする。
$X$ が準コンパクトであるための必要十分条件は、
$U$ がアフィンスキームであるような全射エタール射
$U \to X$ が存在することである。
\end{lemma}

\begin{proof}
$U$ がアフィンであるような全射エタール射 $U \to X$ が存在すれば、
Definition \ref{spaces-properties-definition-quasi-compact} により $X$ は準コンパクトである。
逆に、$X$ が準コンパクトならば $|X|$ は準コンパクトである。
アフィンスキームの非交和
$U = \coprod_{i \in I} U_i$ とエタール全射
$\varphi : U \to X$ を取る
（Lemma \ref{spaces-properties-lemma-cover-by-union-affines}）。
このとき $|X| = \bigcup \varphi(|U_i|)$ であり、
準コンパクト性により、有限部分集合 $i_1, \ldots, i_n$ で
$|X| = \bigcup \varphi(|U_{i_j}|)$ となるものが存在する。
したがって $U_{i_1} \cup \ldots \cup U_{i_n}$ はアフィンスキームであり、
$X$ への有限全射をもつ。
\end{proof}

\noindent
次の補題は、代数空間の射を扱う章における分離射の議論によって
不要になる。

\begin{lemma}
\label{spaces-properties-lemma-separated-cover}
$S$ をスキームとし、$X$ を $S$ 上の代数空間とする。
$U$ を分離スキームとし、$U \to X$ をエタールとする。
このとき $U \to X$ は分離的であり、
$R = U \times_X U$ は分離スキームである。
\end{lemma}

\begin{proof}
$U \to X$ が全射エタール射 $U \to X'$ を経由するような
開部分スキームを $X' \subset X$ とする。
Lemma \ref{spaces-properties-lemma-etale-image-open} を参照されたい。
$U \to X'$ が分離的ならば $U \to X$ も分離的である。Spaces, Lemma
\ref{spaces-lemma-composition-representable-transformations-property}
を参照されたい（開埋め込み $X' \to X$ は Spaces, Lemma
\ref{spaces-lemma-representable-transformations-property-implication}
および Schemes, Lemma
\ref{schemes-lemma-immersions-monomorphisms}
により分離的である）。
さらに $U \times_{X'} U = U \times_X U$ なので、
$X$ を $X'$ で置き換えた後に結論を証明すれば十分である。
すなわち、$U \to X$ は全射であると仮定してよい。
可換図式
$$
\xymatrix{
R = U \times_X U \ar[r] \ar[d] & U \ar[d] \\
U \ar[r] & X
}
$$
を考える。Spaces, Lemma
\ref{spaces-lemma-properties-diagonal}
の証明で、$j : R \to U \times_S U$ が分離的であることを見た。
$U$ は分離スキームなので、スキームの射 $U \to S$ は分離的である。
Schemes, Lemma \ref{schemes-lemma-compose-after-separated}
を参照されたい。したがって、基底変換である
$U \times_S U \to U$ は分離的である。Schemes, Lemma
\ref{schemes-lemma-separated-permanence} を参照されたい。
よって同じ補題により、スキーム $U \times_S U$ は分離的である。
$j$ は分離的なので、同様に $R$ は分離的である。
したがって $R \to U$ は分離射である
（再び Schemes, Lemma \ref{schemes-lemma-compose-after-separated}）。
以上より、Spaces, Lemma
\ref{spaces-lemma-representable-morphisms-spaces-property}
と上の図式から、$U \to X$ が分離的であると結論する。
\end{proof}

\begin{lemma}
\label{spaces-properties-lemma-quasi-separated}
$S$ をスキームとし、$X$ を $S$ 上の代数空間とする。
準分離スキーム $U$ と全射エタール射 $U \to X$ が存在し、
射影 $U \times_X U \to U$ のいずれか一方が準コンパクトならば、
$X$ は準分離的である。
\end{lemma}

\begin{proof}
$X$ を $\mathbf{Z}$ 上の代数空間と考えてよい。
Cartesian 図式
$$
\xymatrix{
U \times_X U \ar[r] \ar[d]_j & X \ar[d]^\Delta \\
U \times U \ar[r] & X \times X
}
$$
を考える。$U$ は準分離的なので、射影 $U \times U \to U$ は
準分離的である（スキームの準分離射の基底変換である。
Schemes, Lemma \ref{schemes-lemma-separated-permanence} を参照）。
したがって補題の仮定と Schemes, Lemma
\ref{schemes-lemma-quasi-compact-permanence}
により $j$ は準コンパクトである。Spaces, Lemma
\ref{spaces-lemma-representable-morphisms-spaces-property}
により、望みどおり $\Delta$ は準コンパクトである。
\end{proof}

\begin{lemma}
\label{spaces-properties-lemma-quasi-separated-quasi-compact-pieces}
$S$ をスキームとし、$X$ を $S$ 上の代数空間とする。
次は同値である。
\begin{enumerate}
\item $X$ は $S$ 上 Zariski 局所準分離的である。
\item $X$ は Zariski 局所準分離的である。
\item Zariski 開被覆 $X = \bigcup X_i$ が存在し、各 $i$ に対して
アフィンスキーム $U_i$ と準コンパクトな全射エタール射
$U_i \to X_i$ が存在する。
\item Zariski 開被覆 $X = \bigcup X_i$ が存在し、各 $i$ に対して
$S$ のアフィン開部分へ写るアフィンスキーム $U_i$ と
準コンパクトな全射エタール射 $U_i \to X_i$ が存在する。
\end{enumerate}
\end{lemma}

\begin{proof}
(3) のとおりに $U_i \to X_i \subset X$ が与えられたとする。
(4) を証明するため、各 $i$ に対して、各 $U_{ij}$ が $S$ の
アフィン開部分へ写るような有限アフィン開被覆
$U_i = U_{i1} \cup \ldots \cup U_{in_i}$ を選ぶ。
合成 $U_{ij} \to U_i \to X_i$ はエタールかつ準コンパクトである
（Spaces, Lemma
\ref{spaces-lemma-composition-representable-transformations-property}）。
$|U_{ij}| \to |X_i|$ の像に対応する開部分空間を
$X_{ij} \subset X_i$ とする。Lemma
\ref{spaces-properties-lemma-etale-image-open} を参照されたい。
$X_{ij} \subset X_i$ は単射であり、$U_{ij} \to X$ は
準コンパクトなので、$U_{ij} \to X_{ij}$ は準コンパクトであることに注意する。
すると $X = \bigcup X_{ij}$ は (4) のような被覆である。
(4) $\Rightarrow$ (3) は直ちに従う。

\medskip\noindent
(4) を仮定する。$X$ が $S$ 上 Zariski 局所準分離的であることを
示すには、$X_i$ が $S$ 上準分離的であることを示せば十分である。
したがって、$S$ のアフィン開部分へ写るアフィンスキーム $U$ と
準コンパクトな全射エタール射 $U \to X$ が存在すると仮定してよい。
ファイバー積図式
$$
\xymatrix{
U \times_X U \ar[r] \ar[d] & U \times_S U \ar[d] \\
X \ar[r]^-{\Delta_{X/S}} & X \times_S X
}
$$
を考える。右の垂直矢印は全射エタールであり
（Spaces, Lemma
\ref{spaces-lemma-product-representable-transformations-property}）、
$U$ は $S$ のアフィン開部分へ写るので $U \times_S U$ はアフィンである
（Schemes, Section \ref{schemes-section-fibre-products}）。
また、射影 $U \times_X U \to U$ は $U \to X$ の基底変換として
準コンパクトなので、$U \times_X U$ は準コンパクトである。
Spaces, Lemma
\ref{spaces-lemma-representable-morphisms-spaces-property}
により、望みどおり $\Delta_{X/S}$ は準コンパクトである。

\medskip\noindent
(1) を仮定する。(3) を証明するには、$X$ が $S$ 上準分離的である場合へ
直ちに帰着できる。Lemma \ref{spaces-properties-lemma-union-of-quasi-compact}
により、各 $X_i$ が $S$ のアフィン開部分へ写り、アフィンスキーム
$U_i$ と全射エタール射 $U_i \to X_i$ が存在するような
Zariski 開被覆 $X = \bigcup X_i$ を取れる。
$U_i \to S$ は $S$ のアフィン開部分へ写るので、
$U_i \times_S U_i$ はアフィンである。Schemes, Section
\ref{schemes-section-fibre-products} を参照されたい。
$X$ は $S$ 上準分離的なので、
$$
R_i = U_i \times_{X_i} U_i = U_i \times_X U_i
\longrightarrow
U_i \times_S U_i
$$
は $\Delta_{X/S}$ の基底変換として準コンパクトである。
したがって $R_i$ は準コンパクトなスキームである。
すると各射影 $R_i \to U_i$ は準コンパクトである。
よって、被覆 $U_i \to X_i$ と射 $U_i \to X_i$ に
Spaces, Lemma
\ref{spaces-lemma-representable-morphisms-spaces-property}
を適用して、望みどおり射 $U_i \to X_i$ が準コンパクトであると結論する。

\medskip\noindent
以上で (1), (3), (4) が同値であることが分かった。(3) は
基底スキームに言及しないので、これらは (2) とも同値である。
\end{proof}

\noindent
次の補題は非常に有用である。

\begin{lemma}
\label{spaces-properties-lemma-finite-fibres-presentation}
$S$ をスキームとし、$X$ を $S$ 上の代数空間とする。
$U$ をスキームとする。$\varphi : U \to X$ をエタール射とし、
射影 $R = U \times_X U \to U$ が準コンパクトであるとする。
たとえば $\varphi$ が準コンパクトならばよい。このとき
$$
|U| \to |X|
\quad\text{および}\quad
|R| \to |X|
$$
のファイバーは有限である。
\end{lemma}

\begin{proof}
$R = U \times_X U$ と書き、$s, t : R \to U$ を射影とする。
$u \in U$ を点とし、その像を $x \in |X|$ とする。
Lemma \ref{spaces-properties-lemma-points-cartesian} により、
$x$ 上の $|U| \to |X|$ のファイバーは
$s(t^{-1}(\{u\}))$ に等しく、
$x$ 上の $|R| \to |X|$ のファイバーは
$t^{-1}(s(t^{-1}(\{u\})))$ である。
$t : R \to U$ はエタールかつ準コンパクトなので有限なファイバーをもつ
（Morphisms, Lemma \ref{morphisms-lemma-etale-over-field}
によりそのファイバーは体のスペクトルの非交和であり、かつ準コンパクトである）。
したがって結論を得る。
\end{proof}








\section{スキームの性質によって定義される空間の性質}
\label{spaces-properties-section-types-properties}

\noindent
スキームの任意のエタール局所的性質は、次の補題によって
代数空間の対応する性質を与える。

\begin{lemma}
\label{spaces-properties-lemma-type-property}
$S$ をスキームとし、$X$ を $S$ 上の代数空間とする。
$\mathcal{P}$ をエタール位相について局所的なスキームの性質とする。
Descent, Definition \ref{descent-definition-property-local}
を参照されたい。次は同値である。
\begin{enumerate}
\item あるスキーム $U$ と全射エタール射 $U \to X$ に対して、
$U$ は性質 $\mathcal{P}$ をもつ。
\item 任意のスキーム $U$ と任意のエタール射 $U \to X$ に対して、
$U$ は性質 $\mathcal{P}$ をもつ。
\end{enumerate}
$X$ が表現可能ならば、これは $\mathcal{P}(X)$ と同値である。
\end{lemma}

\begin{proof}
(2) $\Rightarrow$ (1) は直ちに従う。逆向きについて、
$\mathcal{P}$ をもつスキーム $U$ からの全射エタール射
$U \to X$ を選び、$V$ をエタールな $X$-スキームとする。
このとき $U \times_X V \rightarrow V$ はスキームの全射エタール射なので、
$V$ は $U \times_X V$ から $\mathcal{P}$ を受け継ぎ、
後者は $U$ から $\mathcal{P}$ を受け継ぐ
（Descent, Definition \ref{descent-definition-property-local}
に続く議論を参照）。
最後の主張は (1) と Descent, Definition
\ref{descent-definition-property-local} から明らかである。
\end{proof}

\begin{definition}
\label{spaces-properties-definition-type-property}
$\mathcal{P}$ をエタール位相について局所的なスキームの性質とする。
$S$ をスキームとし、$X$ を $S$ 上の代数空間とする。
Lemma \ref{spaces-properties-lemma-type-property} の同値な条件のいずれかが成り立つとき、
$X$ は {\it 性質 $\mathcal{P}$ をもつ}という。
\end{definition}

\begin{remark}
\label{spaces-properties-remark-list-properties-local-etale-topology}
エタール位相について局所的な性質を列挙する
（fpqc, fppf, syntomic, smooth 位相はエタール位相より強いことに注意する）。
\begin{enumerate}
\item 局所 Noether、Descent, Lemma
\ref{descent-lemma-Noetherian-local-fppf} を参照。
\item Jacobson、Descent, Lemma
\ref{descent-lemma-Jacobson-local-fppf} を参照。
\item 局所 Noether かつ $(S_k)$、Descent, Lemma
\ref{descent-lemma-Sk-local-syntomic} を参照。
\item Cohen--Macaulay、Descent, Lemma
\ref{descent-lemma-CM-local-syntomic} を参照。
\item Gorenstein、Duality for Schemes, Lemma
\ref{duality-lemma-gorenstein-local-syntomic} を参照。
\item 被約、Descent, Lemma
\ref{descent-lemma-reduced-local-smooth} を参照。
\item 正規、Descent, Lemma
\ref{descent-lemma-normal-local-smooth} を参照。
\item 局所 Noether かつ $(R_k)$、Descent, Lemma
\ref{descent-lemma-Rk-local-smooth} を参照。
\item 正則、Descent, Lemma
\ref{descent-lemma-regular-local-smooth} を参照。
\item Nagata、Descent, Lemma
\ref{descent-lemma-Nagata-local-smooth} を参照。
\end{enumerate}
\end{remark}

\noindent
スキームの芽の任意のエタール局所的性質は、代数空間の対応する性質を与える。
必要となる補題は次のとおりである。

\begin{lemma}
\label{spaces-properties-lemma-local-source-target-at-point}
$\mathcal{P}$ をエタール局所的なスキームの芽の性質とする。
Descent, Definition \ref{descent-definition-local-at-point}
を参照されたい。$S$ をスキームとし、$X$ を $S$ 上の代数空間とする。
$x \in |X|$ を $X$ の点とする。$U$ がスキームであるような
エタール射 $a : U \to X$ を考える。次は同値である。
\begin{enumerate}
\item 上のような任意の $U \to X$ と $a(u) = x$ を満たす
$u \in U$ に対して $\mathcal{P}(U, u)$ が成り立つ。
\item 上のようなある $U \to X$ と $a(u) = x$ を満たす
$u \in U$ に対して $\mathcal{P}(U, u)$ が成り立つ。
\end{enumerate}
$X$ が表現可能ならば、これは $\mathcal{P}(X, x)$ と同値である。
\end{lemma}

\begin{proof}
省略する。
\end{proof}

\begin{definition}
\label{spaces-properties-definition-property-at-point}
$S$ をスキームとし、$X$ を $S$ 上の代数空間とする。
$x \in |X|$ とする。$\mathcal{P}$ をエタール局所的なスキームの芽の性質とする。
Lemma \ref{spaces-properties-lemma-local-source-target-at-point}
の同値な条件のいずれかが成り立つとき、
$X$ は {\it $x$ において性質 $\mathcal{P}$ をもつ}という。
\end{definition}

\begin{remark}
\label{spaces-properties-remark-list-properties-local-ring-local-etale-topology}
$P$ を局所環の性質とする。任意のエタール環準同型 $A \to B$ と、
$A$ の素イデアル $\mathfrak p$ の上にある $B$ の素イデアル
$\mathfrak q$ に対して、
$P(A_\mathfrak p) \Leftrightarrow P(B_\mathfrak q)$
が成り立つと仮定する。このとき
$\mathcal{P}(U, u) = P(\mathcal{O}_{U, u})$
とおくことで、スキームの芽 $(U, u)$ のエタール局所的性質を得る。
この状況では、「$x$ における $X$ の局所環が $P$ をもつ」という用語を、
$X$ が $x$ において性質 $\mathcal{P}$ をもつという意味で用いる。
このような性質 $P$ を列挙する。
\begin{enumerate}
\item Noether、More on Algebra, Lemma
\ref{more-algebra-lemma-Noetherian-etale-extension} を参照。
\item 次元 $d$、More on Algebra, Lemma
\ref{more-algebra-lemma-dimension-etale-extension} を参照。
\item 正則、More on Algebra, Lemma
\ref{more-algebra-lemma-regular-etale-extension} を参照。
\item 離散付値環、(2), (3) および Algebra, Lemma
\ref{algebra-lemma-characterize-dvr} から従う。
\item 被約、More on Algebra, Lemma
\ref{more-algebra-lemma-henselization-reduced} を参照。
\item 正規、More on Algebra, Lemma
\ref{more-algebra-lemma-henselization-normal} を参照。
\item Noether かつ深さ $k$、More on Algebra, Lemma
\ref{more-algebra-lemma-henselization-depth} を参照。
\item Noether かつ Cohen--Macaulay、More on Algebra, Lemma
\ref{more-algebra-lemma-henselization-CM} を参照。
\item Noether かつ Gorenstein、Dualizing Complexes, Lemma
\ref{dualizing-lemma-flat-under-gorenstein} を参照。
\end{enumerate}
この条件を満たす性質は他にもあり、たとえば G-環や Nagata がある。
これらが必要になれば、対応する代数的事実の詳しい証明への参照とともに
ここへ追加する。
\end{remark}







\section{構成可能集合}
\label{spaces-properties-section-constructible}

\begin{lemma}
\label{spaces-properties-lemma-locally-constructible}
$S$ をスキームとし、$X$ を $S$ 上の代数空間とする。
$E \subset |X|$ を部分集合とする。次は同値である。
\begin{enumerate}
\item $U$ がスキームである任意のエタール射 $U \to X$ に対して、
$U$ における $E$ の逆像は $U$ の局所構成可能部分集合である。
\item $U$ がアフィンスキームである任意のエタール射 $U \to X$ に対して、
$U$ における $E$ の逆像は $U$ の構成可能部分集合である。
\item $U$ がスキームであるある全射エタール射 $U \to X$ に対して、
$U$ における $E$ の逆像は $U$ の局所構成可能部分集合である。
\end{enumerate}
\end{lemma}

\begin{proof}
Properties, Lemma \ref{properties-lemma-locally-constructible}
により (1) と (2) は同値である。(1) が (3) を含意することは直ちに従う。
そこで、$\varphi^{-1}(E)$ が局所構成可能であるような、スキーム $U$ からの
全射エタール射 $\varphi : U \to X$ が与えられていると仮定する。
$U'$ をスキームとし、$\varphi' : U' \to X$ を別のエタール射とする。
このとき
$$
E'' = \text{pr}_1^{-1}(\varphi^{-1}(E)) = \text{pr}_2^{-1}((\varphi')^{-1}(E))
$$
である。ここで $\text{pr}_1 : U \times_X U' \to U$ と
$\text{pr}_2 : U \times_X U' \to U'$ は射影である。
Morphisms, Lemma
\ref{morphisms-lemma-inverse-image-constructible}
により、$E''$ は $U \times_X U'$ において局所構成可能である。
$W' \subset U'$ をアフィン開部分とする。$\text{pr}_2$ は
エタールであり、したがって開であるから、
$\text{pr}_2(W'') = W'$ となる準コンパクトな開部分
$W'' \subset U \times_X U'$ を選べる。このとき
$\text{pr}_2|_{W''} : W'' \to W'$ は準コンパクトである。
$\varphi$ は全射なので、Lemma \ref{spaces-properties-lemma-points-cartesian}
により
$W' \cap (\varphi')^{-1}(E) = \text{pr}_2(E'' \cap W'')$
である。したがって望みどおり、
Morphisms, Theorem \ref{morphisms-theorem-chevalley}
により $W' \cap (\varphi')^{-1}(E) = \text{pr}_2(E'' \cap W'')$
は局所構成可能である。
\end{proof}

\begin{definition}
\label{spaces-properties-definition-locally-constructible}
$S$ をスキームとし、$X$ を $S$ 上の代数空間とする。
$E \subset |X|$ を部分集合とする。
Lemma \ref{spaces-properties-lemma-locally-constructible} の同値な条件を満たすとき、
$E$ は {\it エタール局所構成可能}であるという。
\end{definition}

\noindent
もちろん $X$ が表現可能、すなわち $X$ がスキームである場合、
これは単に $E$ が台位相空間の局所構成可能部分集合であることを意味する。







\section{点における次元}
\label{spaces-properties-section-dimension}

\noindent
Descent, Lemma \ref{descent-lemma-dimension-at-point-local}
を用いて、点 $x$ における代数空間 $X$ の次元を定義できる。
これは位相的な概念（すなわち $x$ における $|X|$ の次元）とは
異なる概念を与える。

\begin{definition}
\label{spaces-properties-definition-dimension-at-point}
$S$ をスキームとし、$X$ を $S$ 上の代数空間とする。
$x \in |X|$ を $X$ の点とする。
{\it $x$ における $X$ の次元}を、次を満たす元
$\dim_x(X) \in \{0, 1, 2, \ldots, \infty\}$ として定義する。
すなわち、スキームから $X$ へのエタール射 $a : U \to X$ と
$a(u) = x$ を満たす点 $u \in U$ からなる任意の（同値に、ある）組
$(a : U \to X, u)$ に対して、$\dim_x(X) = \dim_u(U)$ とする。
Definition \ref{spaces-properties-definition-property-at-point},
Lemma \ref{spaces-properties-lemma-local-source-target-at-point}, および
Descent, Lemma \ref{descent-lemma-dimension-at-point-local}
を参照されたい。
\end{definition}

\noindent
注意：一般には $\dim_x(X) = \dim_x(|X|)$ は成り立た{\bf ない}。
反例は Spaces, Example \ref{spaces-example-infinite-product}
の代数空間 $X$ である。実際、$x \in |X|$ を $|X|$ の一般点
$x_0$ と異なる点とすると、
$\dim_x(X) = 0$ である一方、$\dim_x(|X|) = 1$ である。
とくに、以下で定義する $X$ の次元は $|X|$ の次元と異なる。

\begin{definition}
\label{spaces-properties-definition-dimension}
$S$ をスキームとし、$X$ を $S$ 上の代数空間とする。
$X$ の {\it 次元} $\dim(X)$ を
$$
\dim(X) = \sup\nolimits_{x \in |X|} \dim_x(X)
$$
によって定義する。
\end{definition}

\noindent
Properties, Lemma \ref{properties-lemma-dimension}
により、$X$ がスキームならばこれは通常の概念である。
点におけるスキームの次元を測る別の整数として、局所環の次元がある。
この不変量もエタール射と両立する。Section
\ref{spaces-properties-section-dimension-local-ring} を参照されたい。






\section{局所環の次元}
\label{spaces-properties-section-dimension-local-ring}

\noindent
代数空間の局所環の次元は、明確に定義される概念である。

\begin{lemma}
\label{spaces-properties-lemma-pre-dimension-local-ring}
$S$ をスキームとし、$X$ を $S$ 上の代数空間とする。
$x \in |X|$ を点とし、$d \in \{0, 1, 2, \ldots, \infty\}$ とする。
次は同値である。
\begin{enumerate}
\item あるスキーム $U$、エタール射 $a : U \to X$、および
$a(u) = x$ を満たす点 $u \in U$ に対して
$\dim(\mathcal{O}_{U, u}) = d$ である。
\item 任意のスキーム $U$、任意のエタール射 $a : U \to X$、および
$a(u) = x$ を満たす任意の点 $u \in U$ に対して
$\dim(\mathcal{O}_{U, u}) = d$ である。
\end{enumerate}
$X$ がスキームならば、これは $\dim(\mathcal{O}_{X, x}) = d$ と同値である。
\end{lemma}

\begin{proof}
Lemma \ref{spaces-properties-lemma-local-source-target-at-point} と
Descent, Lemma \ref{descent-lemma-dimension-local-ring-local}
を組み合わせればよい。
\end{proof}

\begin{definition}
\label{spaces-properties-definition-dimension-local-ring}
$S$ をスキームとし、$X$ を $S$ 上の代数空間とする。
$x \in |X|$ を点とする。{\it $x$ における $X$ の局所環の次元}とは、
Lemma \ref{spaces-properties-lemma-pre-dimension-local-ring} の同値な条件を満たす元
$d \in \{0, 1, 2, \ldots, \infty\}$ のことである。
この場合、{\it $x$ は $X$ 上の余次元 $d$ の点である}ともいう。
\end{definition}

\noindent
次の補題に加えて、読者には Lemmas
\ref{spaces-properties-lemma-dimension-local-ring} および
\ref{spaces-properties-lemma-dimension-decent-invariant-under-etale}
も参照されたい。

\begin{lemma}
\label{spaces-properties-lemma-dimension}
$S$ をスキームとし、$X$ を $S$ 上の代数空間とする。
次の量は等しい。
\begin{enumerate}
\item $X$ の次元。
\item $X$ の局所環の次元の上限。
\item $x \in |X|$ に対する $\dim_x(X)$ の上限。
\end{enumerate}
\end{lemma}

\begin{proof}
(1) と (3) の数は Definition \ref{spaces-properties-definition-dimension}
により等しい。$U$ をスキームとし、$U \to X$ を全射エタール射とする。
定義により、$x \in |X|$ に対する $\dim_x(X)$ の上限は、
$U$ の点 $u$ に対する $\dim_u(U)$ の上限と同じである。
Properties, Lemma \ref{properties-lemma-dimension}
により、これは $\dim(\mathcal{O}_{U, u})$ の上限と同じである。
さらに定義により、これは (2) と同じである。
\end{proof}





\section{一般点}
\label{spaces-properties-section-generic-points}

\noindent
$T$ を位相空間とする。EGA I 第2版によれば、
{\it $T$ の極大点}とは $T$ の既約成分の一般点である。
$T = |X|$ が代数空間 $X$ に付随する位相空間である場合、
極大点には少なくとも二つの概念がある。$T$ を位相空間とみなしたときの
極大点を考えることもできるし、$U$ がスキームで
$U \to X$ がエタール射であるときの $U$ の極大点の像を考えることもできる。
後者の概念は余次元 $0$ の点の集合に対応する
（Lemma \ref{spaces-properties-lemma-codimension-0-points}）。
一般の代数空間については余次元 $0$ の点の方が扱いやすい。
二つの概念は準分離的な代数空間、より一般には良好な代数空間について一致する
（Decent Spaces, Lemma
\ref{decent-spaces-lemma-decent-generic-points}）。

\begin{lemma}
\label{spaces-properties-lemma-codimension-0-points}
$S$ をスキームとし、$X$ を $S$ 上の代数空間とする。
$x \in |X|$ とする。$U$ がスキームであるようなエタール射
$a : U \to X$ を考える。次は同値である。
\begin{enumerate}
\item $x$ は $X$ 上の余次元 $0$ の点である。
\item 上のようなある $U \to X$ と $a(u) = x$ を満たす
$u \in U$ に対して、$u$ は $U$ の既約成分の一般点である。
\item 上のような任意の $U \to X$ と $x$ に写る任意の
$u \in U$ に対して、$u$ は $U$ の既約成分の一般点である。
\end{enumerate}
$X$ が表現可能ならば、これは $x$ が $|X|$ の既約成分の
一般点であることと同値である。
\end{lemma}

\begin{proof}
スキーム $U$ の点 $u$ が $U$ の既約成分の一般点であるための
必要十分条件は $\dim(\mathcal{O}_{U, u}) = 0$ である
（Properties, Lemma \ref{properties-lemma-generic-point}）。
したがって、これは $X$ 上の点の余次元の定義
（Definition \ref{spaces-properties-definition-dimension-local-ring}）から従う。
\end{proof}

\begin{lemma}
\label{spaces-properties-lemma-codimension-0-points-dense}
$S$ をスキームとし、$X$ を $S$ 上の代数空間とする。
$X$ の余次元 $0$ の点の集合は $|X|$ において稠密である。
\end{lemma}

\begin{proof}
$U$ がスキームならば、既約成分の一般点の集合は $U$ において稠密である
（これは任意の準 sober 位相空間について成り立つ）。
したがって $U \to X$ が全射エタール射ならば、
$X$ の余次元 $0$ の点の集合は $|U|$ の稠密部分集合の像である
（Lemma \ref{spaces-properties-lemma-codimension-0-points}）。
$|X|$ は $|U| \to |X|$ に関する商位相をもつので、結論が従う。
\end{proof}





\section{被約空間}
\label{spaces-properties-section-reduced}

\noindent
被約代数空間はすでに Section \ref{spaces-properties-section-types-properties}
で定義した。ここでは被約代数空間に関するいくつかの簡単な補題だけを証明する。

\begin{lemma}
\label{spaces-properties-lemma-subspace-induced-topology}
$S$ をスキームとし、$Z \to X$ を代数空間の埋め込みとする。
このとき $|Z| \to |X|$ は、$|Z|$ から $|X|$ の局所閉部分集合への
同相写像である。
\end{lemma}

\begin{proof}
$U$ をスキームとし、$U \to X$ を全射エタール射とする。
このとき $Z \times_X U \to U$ はスキームの埋め込みなので、
$|Z \times_X U|$ から $|U|$ の局所閉部分集合 $T'$ への同相写像を与える。
Lemma \ref{spaces-properties-lemma-points-cartesian} により、$T'$ は
$|Z| \to |X|$ の像 $T$ の逆像である。関手の変換 $Z \to X$ は単射なので、
$|Z| \to |X|$ は単射である。Spaces, Section
\ref{spaces-section-Zariski} を参照されたい。
Topology, Lemma
\ref{topology-lemma-open-morphism-quotient-topology}
により、$T$ は $|X|$ において局所閉である。
さらに、$|Z \times_X U| \to T'$ は同相写像であり、
$|Z \times_Y U| \to |Z|$ は商写像なので、
連続写像 $|Z| \to T$ は同相写像である。
\end{proof}

\noindent
次の補題は（局所）閉部分空間を構成するために役立つ。

\begin{lemma}
\label{spaces-properties-lemma-subspaces-presentation}
$S$ をスキームとする。$j : R \to U \times_S U$ を
エタール同値関係とする。$X = U/R$ を付随する代数空間とする
（Spaces, Theorem \ref{spaces-theorem-presentation}）。
標準的な全単射
$$
R\text{-不変な局所閉部分スキーム }Z'\text{、その母空間 }U
\leftrightarrow
\text{局所閉部分空間 }Z\text{、その母空間 }X
$$
がある。さらに、$Z \to X$ が閉（それぞれ開）であるための
必要十分条件は、$Z' \to U$ が閉（それぞれ開）であることである。
\end{lemma}

\begin{proof}
$\varphi : U \to X$ を標準写像とする。この全単射は
$Z \to X$ を $Z' = Z \times_X U \to U$ に写す。
定義から直ちに、$Z \to X$ が埋め込み、閉埋め込み、または開埋め込みならば、
$Z' \to U$ もそれぞれ埋め込み、閉埋め込み、または開埋め込みである。
$Z'$ が $R$-不変であることも明らかである
（Groupoids, Definition \ref{groupoids-definition-invariant-open} を参照）。

\medskip\noindent
逆に、$Z' \to U$ が $R$-不変な埋め込みであると仮定する。
$R'$ を $Z'$ への $R$ の制限とする。Groupoids, Definition
\ref{groupoids-definition-restrict-groupoid} を参照されたい。
この場合
$R' = R \times_{s, U} Z' = Z' \times_{U, t} R$
なので、$R'$ は $Z'$ 上のエタール同値関係である。
Spaces, Theorem \ref{spaces-theorem-presentation}
により $Z = Z'/R'$ は代数空間である。構成により
$U \times_X Z = Z'$ なので、$U \times_X Z \to Z$ は埋め込みである。
性質「埋め込み」は基底変換で保たれ、基底上 fppf 局所的であることに注意する
（Spaces, Section \ref{spaces-section-lists}）。
さらに、埋め込みは分離的かつ局所準有限である
（Schemes, Lemma \ref{schemes-lemma-immersions-monomorphisms}
および Morphisms, Lemma
\ref{morphisms-lemma-immersion-locally-quasi-finite}）。
したがって More on Morphisms, Lemma
\ref{more-morphisms-lemma-separated-locally-quasi-finite-morphisms-fppf-descend}
により、埋め込みは fppf 被覆に関する降下を満たす。
これは、$Z \to X$、$\mathcal{P}=$「埋め込み」および
全射エタール射 $U \to X$ について、Spaces, Lemma
\ref{spaces-lemma-morphism-sheaves-with-P-effective-descent-etale}
のすべての仮定が満たされることを意味する。
したがって $Z \to X$ は表現可能な埋め込みである。
これは部分空間の定義である
（Spaces, Definition \ref{spaces-definition-immersion}）。

\medskip\noindent
これらの構成が互いに逆であることは明らかであり、結論を得る。
\end{proof}

\begin{lemma}
\label{spaces-properties-lemma-reduced-closed-subspace}
$S$ をスキームとし、$X$ を $S$ 上の代数空間とする。
$T \subset |X|$ を閉部分集合とする。次の性質をもつ閉部分空間
$Z \subset X$ が一意に存在する：(a) $|Z| = T$ であり、
(b) $Z$ は被約である。
\end{lemma}

\begin{proof}
$U$ をスキームとし、$U \to X$ を全射エタール射とする。
$R = U \times_X U$ とおくと $X = U/R$ である。
Spaces, Lemma \ref{spaces-lemma-space-presentation}
を参照されたい。通常どおり $s, t : R \to U$ を二つの射影と書く。
Lemma \ref{spaces-properties-lemma-points-presentation} により、$T$ は
$s^{-1}(T') = t^{-1}(T')$ を満たす閉部分集合
$T' \subset |U|$ に対応する。$Z' \subset U$ を $T'$ 上の
被約誘導スキーム構造とする。この場合、ファイバー積
$Z' \times_{U, t} R$ と $Z' \times_{U, s} R$ は $R$ の閉部分スキームであり
（Schemes, Lemma \ref{schemes-lemma-base-change-immersion}）、
$Z'$ 上エタールである
（Morphisms, Lemma \ref{morphisms-lemma-base-change-etale}）。
したがってこれらは被約である
（被約性はエタール位相について局所的である。
Remark \ref{spaces-properties-remark-list-properties-local-etale-topology} を参照）。
これらの台位相空間は同じなので、上の議論により
$Z' \times_{U, t} R = Z' \times_{U, s} R$
と結論する。よって Lemma \ref{spaces-properties-lemma-subspaces-presentation}
を適用して、$U$ への引き戻しが $Z'$ である閉部分空間
$Z \subset X$ を得る。構成により $|Z| = T$ であり、$Z$ は被約である。
これで存在が証明された。一意性の証明は省略する。
\end{proof}

\begin{lemma}
\label{spaces-properties-lemma-map-into-reduction}
$S$ をスキームとする。$X$, $Y$ を $S$ 上の代数空間とする。
$Z \subset X$ を閉部分空間とし、$Y$ は被約であると仮定する。
射 $f : Y \to X$ が $Z$ を経由するための必要十分条件は
$f(|Y|) \subset |Z|$ である。
\end{lemma}

\begin{proof}
$f(|Y|) \subset |Z|$ と仮定する。図式
$$
\xymatrix{
V \ar[d]_b \ar[r]_h & U \ar[d]^a \\
Y \ar[r]^f & X
}
$$
で、$U$, $V$ はスキーム、垂直矢印は全射エタールであるものを選ぶ。
Lemma \ref{spaces-properties-lemma-type-property} によりスキーム $V$ は被約である。
したがって Schemes, Lemma
\ref{schemes-lemma-map-into-reduction}
により $h$ は $a^{-1}(Z)$ を経由する。
よって $a \circ h$ は $Z$ を経由する。
$Z \subset X$ は部分層であり、$V \to Y$ は
$(\Sch/S)_{fppf}$ 上の層の全射なので、
$X \to Y$ が $Z$ を経由すると結論する。
\end{proof}

\begin{definition}
\label{spaces-properties-definition-reduced-induced-space}
$S$ をスキームとし、$X$ を $S$ 上の代数空間とする。
$Z \subset |X|$ を閉部分集合とする。
{\it $Z$ 上の代数空間構造}とは、$|Z'|$ が $Z$ に等しい
$X$ の閉部分空間 $Z'$ によって与えられる構造である。
$Z$ 上の {\it 被約誘導代数空間構造}とは、
Lemma \ref{spaces-properties-lemma-reduced-closed-subspace}
で構成したものである。
$X$ の {\it 被約化 $X_{red}$}とは、$|X|$ 上の
被約誘導代数空間構造である。
\end{definition}












\section{スキーム軌跡}
\label{spaces-properties-section-schematic}

\noindent
任意の代数空間は、スキームである最大の開部分空間をもつ。
これはほぼ明らかだが、以下に証明も記す。
もちろんこの部分空間は空でありうる。たとえば
$X = \mathbf{A}^1_{\mathbf{Q}}/\mathbf{Z}$
（普遍的な反例）がそうである。一方、たとえば $X$ が準分離的ならば、
この最大開部分スキームは実際に $X$ で稠密である。

\begin{lemma}
\label{spaces-properties-lemma-subscheme}
$S$ をスキームとし、$X$ を $S$ 上の代数空間とする。
スキームである最大の開部分空間 $X' \subset X$ が存在する。
\end{lemma}

\begin{proof}
$U$ をスキームとし、$U \to X$ を全射エタール射とする。
$R = U \times_X U$ とおく。$X$ の開部分空間は、
$R$-不変な $U$ の開部分スキームと $1 - 1$ に対応する。
したがって、それらは集合をなす。$X$ の開部分空間であって
スキーム、すなわち表現可能であるものの集合を
$X_i$, $i \in I$ とする。点の台集合が $|X|$ の開部分
$\bigcup |X_i|$ である開部分空間を $X' \subset X$ とする。
Lemma \ref{spaces-properties-lemma-characterize-surjective} により
$$
\coprod X_i \longrightarrow X'
$$
は $(\Sch/S)_{fppf}$ 上の層の全射である。
しかし各 $X_i \to X'$ は開埋め込みによって表現可能なので、
実際にはこの写像は Zariski 位相で全射である。
実際、スキームから $X'$ への射 $T \to X'$ を取ると、
$X_i \times_{X'} T$ は $T$ の開部分スキームである。
したがって Schemes, Lemma \ref{schemes-lemma-glue-functors}
を適用して、$X'$ がスキームであることが分かる。
\end{proof}

\noindent
この節の残りでは、代数空間 $X$ の開部分空間 $X'$ が
{\it 稠密}であるとは、対応する開部分集合
$|X'| \subset |X|$ が稠密であることをいう。

\begin{lemma}
\label{spaces-properties-lemma-quasi-separated-finite-etale-cover-dense-open-scheme}
$S$ をスキームとし、$X$ を $S$ 上の代数空間とする。
$U$ が準分離スキームであるような有限エタール全射
$U \to X$ が存在するならば、$X$ にはスキームである
稠密な開部分空間 $X'$ が存在する。
より正確には、$X$ の余次元 $0$ の任意の点
$x \in |X|$ は $X'$ に含まれる。
\end{lemma}

\begin{proof}
$X' \subset X$ をスキームである最大の開部分空間とする
（Lemma \ref{spaces-properties-lemma-subscheme}）。
$x \in |X|$ を $X$ 上の余次元 $0$ の点とする。
Lemma \ref{spaces-properties-lemma-codimension-0-points-dense}
により、$x \in X'$ を示せば十分である。
$U \to X$ を補題の主張のとおりとする。
$R = U \times_X U$ と書き、通常どおり
$s, t : R \to U$ を射影とする。$s, t$ は全射、有限、かつエタールである。
Lemma \ref{spaces-properties-lemma-finite-fibres-presentation}
により、$x$ 上の $|U| \to |X|$ のファイバーは有限であり、
$\{\eta_1, \ldots, \eta_n\}$ と書ける。
Lemma \ref{spaces-properties-lemma-codimension-0-points}
により、各 $\eta_i$ は $U$ の既約成分の一般点である。
Properties, Lemma \ref{properties-lemma-maximal-points-affine}
により、$\{\eta_1, \ldots, \eta_n\}$ を含むアフィン開部分
$W \subset U$ を取れる（ここで $U$ の準分離性を用いる）。
Groupoids, Lemma \ref{groupoids-lemma-find-invariant-affine}
により、$W$ は $R$-不変であると仮定してよい。
$W \subset U$ は $R$-不変なアフィン開部分なので、
$R$ の $W$ への制限 $R_W$ は
$R_W = s^{-1}(W) = t^{-1}(W)$ に等しい
（Groupoids, Definition \ref{groupoids-definition-invariant-open}
およびその後の議論を参照）。とくに写像 $R_W \to W$ も有限エタールである。
したがって $R_W$ はアフィンである。
よって Groupoids, Proposition
\ref{groupoids-proposition-finite-flat-equivalence}
により $W/R_W$ はスキームである。一方、
Spaces, Lemma \ref{spaces-lemma-finding-opens}
により $W/R_W$ は $X$ の開部分空間であり、構成により $x$ を含む。
\end{proof}

\noindent
次の命題は Decent Spaces, Theorem
\ref{decent-spaces-theorem-decent-open-dense-scheme}
において、良好な代数空間の場合へ改良される。

\begin{proposition}
\label{spaces-properties-proposition-locally-quasi-separated-open-dense-scheme}
$S$ をスキームとし、$X$ を $S$ 上の代数空間とする。
$X$ が Zariski 局所準分離的ならば
（たとえば $X$ が準分離的ならば）、
$X$ にはスキームである稠密な開部分空間 $X'$ が存在する。
より正確には、$X$ 上の余次元 $0$ の任意の点
$x \in |X|$ は $X'$ に含まれる。
\end{proposition}

\begin{proof}
Lemma \ref{spaces-properties-lemma-subscheme} により、問題は $X$ 上局所的である。
したがって Lemma
\ref{spaces-properties-lemma-quasi-separated-quasi-compact-pieces}
により、アフィンスキーム $U$ と全射、準コンパクト、エタールな射
$U \to X$ が存在すると仮定してよい。
さらに $U \to X$ は分離的である
（Lemma \ref{spaces-properties-lemma-separated-cover}）。
$R = U \times_X U$ とおき、通常どおり
$s, t : R \to U$ を射影とする。
このとき $s, t$ は全射、準コンパクト、分離的、かつエタールである。
したがって $s, t$ は局所準有限でもあり、有限なファイバーをもつ
（Morphisms, Lemmas
\ref{morphisms-lemma-etale-locally-quasi-finite},
\ref{morphisms-lemma-quasi-finite-locally-quasi-compact}, および
\ref{morphisms-lemma-quasi-finite}）。
Morphisms, Lemma \ref{morphisms-lemma-generically-finite}
により、$U$ の既約成分の一般点である各 $\eta \in U$ に対して、
$s^{-1}(V) \to V$ が有限となる $\eta$ の開近傍
$V \subset U$ が存在する。
Descent, Lemma \ref{descent-lemma-descending-property-finite}
により、有限性は基底上 fpqc（とくにエタール）局所的である。
したがって More on Groupoids, Lemma
\ref{more-groupoids-lemma-property-invariant}
を適用でき、$s$ が有限となる最大の開部分 $W \subset U$ は
$R$-不変である。この構成により、$W$ は $U$ の既約成分の
すべての一般点を含む。
$R$ の $W$ への制限 $R_W$ は
$R_W = s^{-1}(W) = t^{-1}(W)$ に等しい
（Groupoids, Definition \ref{groupoids-definition-invariant-open}
およびその後の議論を参照）。
構成により $s_W, t_W : R_W \to W$ は有限エタールである。
開部分空間 $X' = W/R_W \subset X$ を考える
（Spaces, Lemma \ref{spaces-lemma-finding-opens}）。
構成により包含写像 $X' \to X$ は余次元 $0$ の点の上で全単射を誘導する。
これで Lemma
\ref{spaces-properties-lemma-quasi-separated-finite-etale-cover-dense-open-scheme}
に帰着した。
\end{proof}




\section{スキームを得ること}
\label{spaces-properties-section-getting-a-scheme}

\noindent
前節では、アフィンスキーム $U$ を同値関係 $R$ で割った商 $U/R$ が、
射 $s, t : R \to U$ が有限エタールならばスキームであることを用いた。
これは次の結果の特別な場合である。

\begin{proposition}
\label{spaces-properties-proposition-finite-flat-equivalence-global}
$S$ をスキームとし、$(U, R, s, t, c)$ を $S$ 上の群oidスキームとする。
次を仮定する。
\begin{enumerate}
\item $s, t : R \to U$ は有限局所自由である。
\item $j = (t, s)$ は同値関係である。
\item $U$ の任意の空でない閉部分集合 $Z$ は、
$R$-同値類 $t(s^{-1}(\{u\}))$ が $U$ のアフィン開部分に含まれるような
点 $u$ を含む\footnote{$t(s^{-1}(\{u\}))$ が $U$ のアフィン開部分に
含まれるような点 $u$ の集合を $E \subset U$ とする。
$E = U$ ならば、または $U$ のすべての有限型点が $E$ に属するならば、
または各 $u \in U$ が $E$ の点へ特殊化するならば、条件 (3) は成り立つ。}。
\end{enumerate}
このとき、$S$ 上のスキームの有限局所自由射 $U \to M$ が存在し、
$R = U \times_M U$ であり、かつ $M$ は fppf 位相における
商層 $U/R$ を表現する。
\end{proposition}

\begin{proof}
仮定 (3) と Groupoids, Lemma
\ref{groupoids-lemma-find-invariant-affine}
により、各 $U_i$ が $U$ の $R$-不変なアフィン開部分であるような
開被覆 $U = \bigcup U_i$ を取れる。$R_i = R|_{U_i}$ とおく。
fppf 層 $F = U/R$ と $F_i = U_i/R_i$ を考える。
Spaces, Lemma \ref{spaces-lemma-finding-opens}
により射 $F_i \to F$ は表現可能な開埋め込みである。
Groupoids, Proposition
\ref{groupoids-proposition-finite-flat-equivalence}
により層 $F_i$ はアフィンスキームによって表現可能である。
$T$ がスキームで $T \to F$ が射ならば、
$V_i = F_i \times_F T$ は $T$ で開であり、
$T = \bigcup V_i$ であると主張する。
実際、$T$ 上 fppf 局所的に $T \to F$ を射
$f : T \to U$ へ持ち上げられ、その場合
$f^{-1}(U_i) \subset V_i$ である。
したがって Schemes, Lemma
\ref{schemes-lemma-glue-functors}
により $F$ はスキームによって表現可能である。
\end{proof}

\noindent
たとえば $U$ がアフィンスキームの局所閉部分スキームと同型であるか、
ある次数付き環 $A$ に対する $\text{Proj}(A)$ の局所閉部分スキームと
同型であるならば、Properties, Lemma
\ref{properties-lemma-ample-finite-set-in-affine}
により第三の仮定が成り立つ。
とくに、準アフィンまたは準射影スキーム上の有限群および有限群スキームの
自由作用にこれを適用できる。たとえば、準射影多様体 $X$ 上の
有限群 $G$ の自由作用による商 $X/G$ はスキームである。
詳しい主張は次のとおりである。

\begin{lemma}
\label{spaces-properties-lemma-quotient-scheme}
$S$ をスキームとする。$G \to S$ を群スキームとし、
$X \to S$ をスキームの射とする。
$a : G \times_S X \to X$ を作用とする。次を仮定する。
\begin{enumerate}
\item $G \to S$ は有限局所自由である。
\item 作用 $a$ は自由である。
\item $X \to S$ はアフィン、準アフィン、射影的、または準射影的であるか、
$X$ はアフィンスキームの開部分スキームと同型であるか、
$X$ はある次数付き環 $A$ に対する $\text{Proj}(A)$ の
開部分スキームと同型であるか、または $G \to S$ は純非分離である。
\end{enumerate}
このとき fppf 商層 $X/G$ はスキームであり、
$X \to X/G$ は fppf $G$-トーサーである。
\end{lemma}

\begin{proof}
まず $X/G$ がスキームであることを示す。
作用は自由なので、射
$j = (a, \text{pr}) : G \times_S X \to X \times_S X$
は単射であり、したがって同値関係である。
Groupoids, Lemma \ref{groupoids-lemma-free-action}
を参照されたい。$G \to S$ は有限局所自由であると仮定したので、
写像 $s, t : G \times_S X \to X$ は有限局所自由である。
Proposition \ref{spaces-properties-proposition-finite-flat-equivalence-global}
の最後の仮定が成り立つことを示せば十分である。
$G$ の作用は $S$ 上なので、$X \to S$ の一つのファイバーにある
任意の有限個の点が $X$ のアフィン開部分に含まれることを示せばよい。
$X$ がアフィンスキームの開部分スキームと同型であるか、
ある次数付き環 $A$ に対する $\text{Proj}(A)$ の
開部分スキームと同型である場合、これは Properties, Lemma
\ref{properties-lemma-ample-finite-set-in-affine}
から従う。$X \to S$ がアフィン、準アフィン、射影的、または
準射影的ならば、$S$ をアフィン開部分で置き換えることにより、
いま扱った場合に戻る。$G \to S$ が純非分離ならば、
$G$ の作用による $X$ の点の軌道は一点集合であり、条件は自明に成り立つ。
細部は省略する。

\medskip\noindent
$X \to X/G$ が fppf $G$-トーサーであることを見るには
（Groupoids, Definition
\ref{groupoids-definition-principal-homogeneous-space}）、
$G \times_S X \to X \times_{X/G} X$ が同型であることと、
$X \to X/G$ が fppf 局所的に切断をもつことを示さなければならない。
第二の部分は、構成により $X \to X/G$ が fppf 層の写像として
全射であることから明らかである。第一の部分は Proposition
\ref{spaces-properties-proposition-finite-flat-equivalence-global}
の結論にある同型 $R = U \times_M U$ から従う
（本件では $R = G \times_S X$ であることに注意する）。
\end{proof}


\begin{lemma}
\label{spaces-properties-lemma-quotient-separated}
記法と仮定は Proposition
\ref{spaces-properties-proposition-finite-flat-equivalence-global}
のとおりとする。このとき
\begin{enumerate}
\item $U$ が $S$ 上準分離的ならば、$U/R$ は $S$ 上準分離的である。
\item $U$ が準分離的ならば、$U/R$ は準分離的である。
\item $U$ が $S$ 上分離的ならば、$U/R$ は $S$ 上分離的である。
\item $U$ が分離的ならば、$U/R$ は分離的である。
\item ここにさらに追加する。
\end{enumerate}
Lemma \ref{spaces-properties-lemma-quotient-scheme} の設定でも同様の結果が成り立つ。
\end{lemma}

\begin{proof}
$M$ は商層を表現するので、スキームの Cartesian 図式
$$
\xymatrix{
R \ar[r]_-j \ar[d] & U \times_S U \ar[d] \\
M \ar[r] & M \times_S M
}
$$
がある。$U \times_S U \to M \times_S M$ は全射かつ有限局所自由なので、
$M \to M \times_S M$ が準コンパクト、あるいは閉埋め込みであることを
示すには、それぞれ $j : R \to U \times_S U$ が準コンパクト、
あるいは閉埋め込みであることを示せば十分である。
Descent, Lemmas
\ref{descent-lemma-descending-property-quasi-compact} および
\ref{descent-lemma-descending-property-closed-immersion}
を参照されたい。
$j : R \to U \times_S U$ は $U$ 上の射であり、$R$ は $U$ 上有限なので、
射影 $U \times_S U \to U$ が準分離的であれば $j$ は準コンパクトである
（Schemes, Lemma \ref{schemes-lemma-quasi-compact-permanence}）。
$j$ は単射かつ局所有限型なので、$j$ が proper であれば閉埋め込みである
（\'Etale Morphisms, Lemma
\ref{etale-lemma-characterize-closed-immersion}）。
そして射影 $U \times_S U \to U$ が分離的ならば、この条件は成り立つ
（Morphisms, Lemma \ref{morphisms-lemma-image-proper-scheme-closed}）。
これで (1) と (3) が証明された。(2) と (4) を証明するには、
$S$ を $\Spec(\mathbf{Z})$ で置き換える。
Definition \ref{spaces-properties-definition-separated} を参照されたい。
Lemma \ref{spaces-properties-lemma-quotient-scheme} は Proposition
\ref{spaces-properties-proposition-finite-flat-equivalence-global}
を適用して証明されるので、最後の主張も明らかである。
\end{proof}




\section{準分離空間の点}
\label{spaces-properties-section-points-quasi-separated}

\noindent
Stacks project で導入される一般性のもとでは、代数空間の点は
非常に悪く振る舞いうる。しかし準分離空間については、その振る舞いは
ほとんどスキームの点の振る舞いと同様である。この節ではこれについて
いくつかの結果を証明する。良好な空間を扱う章にはさらに多くの結果がある。
たとえば Decent Spaces, Section
\ref{decent-spaces-section-points} を参照されたい。

\begin{lemma}
\label{spaces-properties-lemma-quasi-separated-sober}
$S$ をスキームとする。$X$ を $S$ 上の Zariski 局所準分離的な
代数空間とする。このとき位相空間 $|X|$ は sober である
（Topology, Definition \ref{topology-definition-generic-point}）。
\end{lemma}

\begin{proof}
Topology, Lemma \ref{topology-lemma-sober-local}
と Lemma \ref{spaces-properties-lemma-quasi-separated-quasi-compact-pieces}
を組み合わせれば、アフィンスキーム $U$ と全射、準コンパクト、
エタールな射 $U \to X$ が存在すると仮定してよい。
$R = U \times_X U$ とおき、射影を $s, t : R \to U$ とする。
Lemma \ref{spaces-properties-lemma-finite-fibres-presentation}
を適用すると、$s, t$ のファイバーは有限である。
したがって Topology, Lemma
\ref{topology-lemma-quotient-kolmogorov}
の仮定がすべて満たされ、$|X|$ は Kolmogorov であると結論する
\footnote{実際にはここで、Schemes, Lemma
\ref{schemes-lemma-scheme-sober}（スキームの soberness）、
Morphisms, Lemmas \ref{morphisms-lemma-etale-flat} および
\ref{morphisms-lemma-generalizations-lift-flat}
（一般化はエタール射に沿って持ち上がる）、
Lemma \ref{spaces-properties-lemma-points-presentation}
（表示を用いた代数空間の点）、ならびに
Lemma \ref{spaces-properties-lemma-topology-points}
（商写像の開性）も用いる。}。

\medskip\noindent
すべての既約閉部分集合 $T \subset |X|$ が一般点をもつことを
示せば残りの証明は終わる。
Lemma \ref{spaces-properties-lemma-reduced-closed-subspace}
により、$|Z| = |T|$ を満たす閉部分空間 $Z \subset X$ が存在する。
$U \times_X Z \to Z$ はアフィンスキームから $Z$ への準コンパクトな
全射エタール射なので、Lemma
\ref{spaces-properties-lemma-quasi-separated-quasi-compact-pieces}
により $Z$ は Zariski 局所準分離的であることに注意する。
Proposition
\ref{spaces-properties-proposition-locally-quasi-separated-open-dense-scheme}
により、スキームである稠密な開部分空間 $Z' \subset Z$ が存在する。
これは $|Z'| \subset T$ が稠密開であることを意味する。
したがって位相空間 $|Z'|$ は既約であり、すなわち $Z'$ は既約スキームである。
Schemes, Lemma \ref{schemes-lemma-scheme-sober}
により、$|Z'|$ は一つの点
$\eta \in |Z'| \subset T$ の閉包である。
したがって $T = \overline{\{\eta\}}$ でもあり、結論を得る。
\end{proof}

\begin{lemma}
\label{spaces-properties-lemma-quasi-compact-quasi-separated-spectral}
$S$ をスキームとし、$X$ を $S$ 上の準コンパクトかつ
準分離的な代数空間とする。位相空間 $|X|$ は spectral 空間である。
\end{lemma}

\begin{proof}
Topology, Definition \ref{topology-definition-spectral-space}
により、$|X|$ が sober、準コンパクトであり、準コンパクト開部分からなる
基底をもち、任意の二つの準コンパクト開部分の共通部分が
準コンパクトであることを確かめなければならない。
Lemma \ref{spaces-properties-lemma-quasi-separated-sober} により $|X|$ は sober である。
Lemma \ref{spaces-properties-lemma-quasi-compact-space} により $|X|$ は準コンパクトである。
Lemma \ref{spaces-properties-lemma-quasi-compact-affine-cover}
により、アフィンスキーム $U$ と全射エタール射
$f : U \to X$ が存在する。
$|f| : |U| \to |X|$ は開かつ連続であり、$|U|$ は準コンパクト開部分からなる
基底をもつので、$|X|$ も準コンパクト開部分からなる基底をもつ。
最後に $A, B \subset |X|$ を準コンパクト開部分とする。
すると、ある開部分空間 $X', X'' \subset X$ に対して
$A = |X'|$ および $B = |X''|$ である
（Lemma \ref{spaces-properties-lemma-open-subspaces}）。
また、アフィンスキーム $V$, $W$ と全射エタール射
$V \to X'$, $W \to X''$ を選べる
（Lemma \ref{spaces-properties-lemma-quasi-compact-affine-cover}）。
このとき $A \cap B$ は $|V \times_X W| \to |X|$ の像である
（Lemma \ref{spaces-properties-lemma-points-cartesian}）。
$X$ は準分離的なので $V \times_X W$ は準コンパクトであり
（Lemma \ref{spaces-properties-lemma-characterize-quasi-separated}）、
$A \cap B$ は準コンパクトである。これで証明が終わる。
\end{proof}

\noindent
次の補題を用いて、代数空間が体のスペクトルと同型であることを証明できる。

\begin{lemma}
\label{spaces-properties-lemma-point-like-spaces}
$S$ をスキームとし、$k$ を体とする。
$X$ を $S$ 上の代数空間とし、全射エタール射
$\Spec(k) \to X$ が存在すると仮定する。
$X$ が準分離的ならば $X \cong \Spec(k')$ であり、
$k/k'$ は有限分離拡大である。
\end{lemma}

\begin{proof}
$R = \Spec(k) \times_X \Spec(k)$ とおくと、ファイバー積図式
$$
\xymatrix{
R \ar[r]_-s \ar[d]_-t & \Spec(k) \ar[d] \\
\Spec(k) \ar[r] & X
}
$$
がある。Spaces, Lemma
\ref{spaces-lemma-space-presentation}
により、$X = \Spec(k)/R$ は商層である。
$\Spec(k) \to X$ はエタールなので、射 $s$ と $t$ はエタールである。
したがって $R = \coprod_{i \in I} \Spec(k_i)$ は体のスペクトルの非交和であり、
$s$ と $t$ はともに有限分離体拡大
$s, t : k \subset k_i$ を誘導する。Morphisms, Lemma
\ref{morphisms-lemma-etale-over-field} を参照されたい。
$$
R = \Spec(k) \times_X \Spec(k)
= (\Spec(k) \times_S \Spec(k)) \times_{X \times_S X, \Delta} X
$$
であり、仮定により $\Delta$ は準コンパクトなので、
$R \to \Spec(k) \times_S \Spec(k)$ は準コンパクトである。
$\Spec(k) \times_S \Spec(k)$ はアフィンなので、$R$ は準コンパクトである。
したがって $I$ は有限である。これは $s$ と $t$ が有限局所自由射であることを
意味する。よって Groupoids, Proposition
\ref{groupoids-proposition-finite-flat-equivalence}
により、$\Spec(k)/R$ は $\Spec(k')$ によって表現される。
ここで $k' \subset k$ は有限局所自由で、
$$
k' = \{x \in k \mid s_i(x) = t_i(x)\text{、すべての }i \in I\}
$$
である。$k'$ が体であることは容易に分かる。
\end{proof}

\begin{remark}
\label{spaces-properties-remark-cannot-decide-yet}
Lemma \ref{spaces-properties-lemma-point-like-spaces} は良好な代数空間について成り立つ。
Decent Spaces, Lemma
\ref{decent-spaces-lemma-decent-point-like-spaces}
を参照されたい。実際、一点をもつ良好な代数空間はスキームである。
Decent Spaces, Lemma \ref{decent-spaces-lemma-when-field}
を参照されたい。局所分離的な代数空間は良好なので、
$X$ が局所分離的な場合にもこれは成り立つ。Decent Spaces, Lemma
\ref{decent-spaces-lemma-locally-separated-decent}
を参照されたい。
\end{remark}







\section{代数空間のエタール射}
\label{spaces-properties-section-etale-morphisms}

\noindent
この節は本来、代数空間の射を扱う章に属する。しかし代数空間の
小エタールサイトを定義するために、別の代数空間上エタールな
代数空間という概念が必要である。したがって、スキームから
代数空間へのエタール射、および代数空間の間のエタール射について
いくらか準備をしなければならない。代数空間のエタール射について
さらに詳しくは Morphisms of Spaces, Section
\ref{spaces-morphisms-section-etale} を参照されたい。

\begin{lemma}
\label{spaces-properties-lemma-etale-over-space}
$S$ をスキームとする。$X$ を $S$ 上の代数空間とし、
$U$, $U'$ を $S$ 上のスキームとする。
\begin{enumerate}
\item $U \to U'$ がスキームのエタール射であり、
$U' \to X$ が $U'$ から $X$ へのエタール射ならば、
合成 $U \to X$ は $U$ から $X$ へのエタール射である。
\item $\varphi : U \to X$ と $\varphi' : U' \to X$ が
$X$ へのエタール射であり、$\varphi = \varphi' \circ \chi$
を満たすスキームの射 $\chi : U \to U'$ があるならば、
$\chi$ はスキームのエタール射である。
\item $\chi : U \to U'$ がスキームの全射エタール射であり、
$\varphi' : U' \to X$ が $\varphi = \varphi' \circ \chi$
をエタールにする射ならば、$\varphi'$ はエタールである。
\end{enumerate}
\end{lemma}

\begin{proof}
スキームから代数空間へのエタール射の定義は、
スキームから代数空間への任意の射が表現可能であることを介して、
Spaces, Definition
\ref{spaces-definition-relative-representable-property}
から得られることを思い出そう。

\medskip\noindent
補題の (1) はこの事実、エタール射が合成で保たれること
（Morphisms, Lemma \ref{morphisms-lemma-composition-etale}）、
ならびに形式的な Spaces, Lemmas
\ref{spaces-lemma-composition-representable-transformations-property}
および
\ref{spaces-lemma-morphism-schemes-gives-representable-transformation-property}
から従う。

\medskip\noindent
(2) を証明するため、$S$ 上のスキーム $W$ と全射エタール射
$W \to X$ を選ぶ。$\chi$ の基底変換
$\chi_W : W \times_X U \to W \times_X U'$ を考える。
$W \times_X U$ と $W \times_X U'$ は $W$ 上エタールなので、
Morphisms, Lemma \ref{morphisms-lemma-etale-permanence}
により $\chi_W$ はエタールである。一方、可換図式
$$
\xymatrix{
W \times_X U \ar[r] \ar[d] & W \times_X U' \ar[d] \\
U \ar[r] & U'
}
$$
の二つの垂直矢印はエタールかつ全射である。
したがって Descent, Lemma
\ref{descent-lemma-syntomic-smooth-etale-permanence}
により $U \to U'$ はエタールである。

\medskip\noindent
(3) を証明するため、$W$ を $S$ 上のスキームとし、
$W \to X$ を射とする。上と同様に図式
$$
\xymatrix{
W \times_X U \ar[r] \ar[d] & W \times_X U' \ar[d] \ar[r] & W \ar[d] \\
U \ar[r] & U' \ar[r] & X
}
$$
を考える。$W \times_X U \to W \times_X U'$ は
$U \to U'$ の基底変換として全射エタールであり、
$W \times_X U \to W$ はエタールである。
したがって Descent, Lemma
\ref{descent-lemma-syntomic-smooth-etale-permanence}
により $W \times_X U' \to W$ はエタールである。
定義により、これは $\varphi'$ がエタールであることを意味する。
\end{proof}

\begin{definition}
\label{spaces-properties-definition-etale}
$S$ をスキームとする。$S$ 上の代数空間の射
$f : X \to Y$ が {\it エタール}であるとは、$U$ がスキームである
任意のエタール射 $\varphi : U \to X$ に対して、
合成 $f \circ \varphi$ もエタールであることをいう。
\end{definition}

\noindent
$X$ と $Y$ がスキームならば、これは通常のスキームのエタール射の
概念と一致する。実際、$X \to Y$ が代数空間の表現可能射であるときは常に、
これは Spaces, Definition
\ref{spaces-definition-relative-representable-property}
を介して定義される概念と一致する。
これは、以下の Lemma \ref{spaces-properties-lemma-etale-local} と
Spaces, Lemma \ref{spaces-lemma-representable-morphisms-spaces-property}
を組み合わせれば従う。

\begin{lemma}
\label{spaces-properties-lemma-etale-local}
$S$ をスキームとし、$f : X \to Y$ を $S$ 上の代数空間の射とする。
次は同値である。
\begin{enumerate}
\item $f$ はエタールである。
\item $U$ がスキームであり、合成 $f \circ \varphi$ が
（代数空間の射として）エタールであるような全射エタール射
$\varphi : U \to X$ が存在する。
\item $V$ がスキームであり、基底変換 $V \times_Y X \to V$ が
（代数空間の射として）エタールであるような全射エタール射
$\psi : V \to Y$ が存在する。
\item 可換図式
$$
\xymatrix{
U \ar[d] \ar[r] & V \ar[d] \\
X \ar[r] & Y
}
$$
が存在し、$U$, $V$ はスキーム、垂直矢印はエタール、
左の垂直矢印は全射であり、水平矢印はエタールである。
\end{enumerate}
\end{lemma}

\begin{proof}
(4) が (1) を含意することを示そう。(4) のような図式が与えられたとする。
$W$ をスキームとし、$W \to X$ をエタール射とする。
このとき $W \times_X U \to U$ はエタールである。
したがって $W \times_X U \to V$ はスキームのエタール射
$W \times_X U \to U$ と $U \to V$ の合成としてエタールである。
よって Lemma \ref{spaces-properties-lemma-etale-over-space} (1) により
$W \times_X U \to Y$ はエタールである。
射影 $W \times_X U \to W$ も全射エタールなので、
Lemma \ref{spaces-properties-lemma-etale-over-space} (3) から
$W \to Y$ がエタールであると結論する。

\medskip\noindent
(1) が (4) を含意することを示そう。(1) を仮定する。
可換図式
$$
\xymatrix{
U \ar[d] \ar[r] & V \ar[d] \\
X \ar[r] & Y
}
$$
で、$U \to X$ と $V \to Y$ は全射エタールであるものを選ぶ。
Spaces, Lemma
\ref{spaces-lemma-lift-morphism-presentations}
を参照されたい。仮定により射 $U \to Y$ はエタールであり、
したがって Lemma \ref{spaces-properties-lemma-etale-over-space} (2) により
$U \to V$ はエタールである。

\medskip\noindent
(2) と (3) も (1) と同値であることの証明は省略する。
\end{proof}

\begin{lemma}
\label{spaces-properties-lemma-composition-etale}
代数空間の二つのエタール射の合成はエタールである。
\end{lemma}

\begin{proof}
定義から直ちに従う。
\end{proof}

\begin{lemma}
\label{spaces-properties-lemma-base-change-etale}
代数空間のエタール射を代数空間の任意の射で基底変換したものは
エタールである。
\end{lemma}

\begin{proof}
$X \to Y$ を $S$ 上の代数空間のエタール射とし、
$Z \to Y$ を代数空間の射とする。
スキーム $U$ と全射エタール射 $U \to X$ を選ぶ。
スキーム $W$ と全射エタール射 $W \to Z$ を選ぶ。
すると $U \to Y$ はエタールなので、図式
$$
\xymatrix{
W \times_Y U \ar[d] \ar[r] & W \ar[d] \\
Z \times_Y X \ar[r] & Z
}
$$
の上の水平矢印はエタールである。
さらに左の垂直矢印は全射エタールである
（検証は省略する）。したがって Lemma
\ref{spaces-properties-lemma-etale-local} により、下の水平矢印はエタールである。
\end{proof}

\begin{lemma}
\label{spaces-properties-lemma-etale-permanence}
$S$ をスキームとする。$X, Y, Z$ を代数空間とする。
$g : X \to Z$, $h : Y \to Z$ をエタール射とし、
$h \circ f = g$ を満たす射 $f : X \to Y$ を取る。
このとき $f$ はエタールである。
\end{lemma}

\begin{proof}
可換図式
$$
\xymatrix{
U \ar[d] \ar[r]_\chi & V \ar[d] \\
X \ar[r] & Y
}
$$
で、$U \to X$ と $V \to Y$ は全射エタールであるものを選ぶ。
Spaces, Lemma
\ref{spaces-lemma-lift-morphism-presentations}
を参照されたい。仮定により、射
$\varphi : U \to X \to Z$ と
$\psi : V \to Y \to Z$ はエタールである。
さらに $f, g, h$ に関する仮定により
$\psi \circ \chi = \varphi$ である。
したがって Lemma \ref{spaces-properties-lemma-etale-over-space} (2) により
$U \to V$ はエタールである。
\end{proof}

\begin{lemma}
\label{spaces-properties-lemma-etale-open}
$S$ をスキームとする。$X \to Y$ が $S$ 上の代数空間の
エタール射ならば、付随する位相空間の写像
$|X| \to |Y|$ は開である。
\end{lemma}

\begin{proof}
これは Lemma \ref{spaces-properties-lemma-etale-local} の図式と
Lemma \ref{spaces-properties-lemma-topology-points} から明らかである。
\end{proof}

\noindent
最後に、面白い補題を一つ述べる。スキームへのエタール射をもつ
代数空間がスキームであるとは限らない。Spaces, Example
\ref{spaces-example-non-representable-descent}
を参照されたい。しかし標的が体のスペクトルならば、これは正しい。

\begin{lemma}
\label{spaces-properties-lemma-etale-over-field-scheme}
$S$ をスキームとする。$k$ を体とし、
$X \to \Spec(k)$ を $S$ 上のエタール射とする。
このとき $X$ はスキームである。
\end{lemma}

\begin{proof}
$U$ をアフィンスキームとし、$U \to X$ をエタール射とする。
Definition \ref{spaces-properties-definition-etale} により、
$U \to \Spec(k)$ はエタール射である。
したがって $U = \coprod_{i = 1, \ldots, n} \Spec(k_i)$ は、
$k$ の有限分離拡大 $k_i$ のスペクトルの有限非交和である。
Morphisms, Lemma \ref{morphisms-lemma-etale-over-field}
を参照されたい。
$R = U \times_X U \to U \times_{\Spec(k)} U$ は単射であり、
$U \times_{\Spec(k)} U$ も $k$ の有限分離拡大のスペクトルの
有限非交和である。したがって Schemes, Lemma
\ref{schemes-lemma-mono-towards-spec-field}
により、$R$ も同様に $k$ の有限分離拡大のスペクトルの有限非交和である。
この $U$ と $R$ はアフィンであり、二つの射影
$R \to U$ はともに有限局所自由である。
したがって Groupoids, Proposition
\ref{groupoids-proposition-finite-flat-equivalence}
により $U/R$ はスキームである。
さらに Spaces, Lemma \ref{spaces-lemma-finding-opens}
により、これは $X$ の開部分空間でもある。
Lemma \ref{spaces-properties-lemma-subscheme} により $X$ はスキームである。
\end{proof}













\section{空間と fpqc 被覆}
\label{spaces-properties-section-fpqc}

\noindent
$S$ をスキームとする。$S$ 上の代数空間は、追加の性質をもつ
fppf 位相の層として定義される。したがって、それが fpqc 位相についても
層条件を満たすことは直ちには明らかでない
（Topologies, Definition
\ref{topologies-definition-sheaf-property-fpqc}）。
この節では、それが正しいことを示す Gabber の議論を与える。
ただし、代数空間 $X$ が fpqc 位相について層条件を満たすと言うとき、
実際に考えるのは $T, T_i$ が大サイト
$(\Sch/S)_{fppf}$ の対象であるような fpqc 被覆
$\{f_i : T_i \to T\}_{i \in I}$ だけである
（本章の規約どおりである。Section
\ref{spaces-properties-section-conventions} を参照）。

\begin{proposition}[Gabber]
\label{spaces-properties-proposition-sheaf-fpqc}
$S$ をスキームとし、$X$ を $S$ 上の代数空間とする。
このとき $X$ は fpqc 位相について層条件を満たす。
\end{proposition}

\begin{proof}
$X$ は Zariski 位相について層なので、次を示せば十分である。
アフィンの全射平坦射 $f : T' \to T$ が与えられたとき、
$X(T)$ は二つの写像 $X(T') \to X(T' \times_T T')$ の等化子である。
Topologies, Lemma \ref{topologies-lemma-sheaf-property-fpqc}
を参照されたい（引用した補題はすべてのスキームの圏上で定義された
関手について定式化されているため、ここでは短い議論を一つ省略している）。

\medskip\noindent
$a, b : T \to X$ を $a \circ f = b \circ f$ を満たす二つの射とする。
$a = b$ を示さなければならない。ファイバー積
$$
E = X \times_{\Delta_{X/S}, X \times_S X, (a, b)} T.
$$
を考える。Spaces, Lemma
\ref{spaces-lemma-properties-diagonal}
により $\Delta_{X/S}$ は表現可能な単射である。
したがって $E \to T$ はスキームの単射である。
仮定 $a \circ f = b \circ f$ は、$T' \to T$ が
（一意に）$E$ を経由することを含意する。可換図式
$$
\xymatrix{
T' \times_T E \ar[r] \ar[d] & E \ar[d] \\
T' \ar[r] \ar@/^5ex/[u] \ar[ru] & T
}
$$
を考える。射影 $T' \times_T E \to T'$ は切断をもつ単射なので、
同型である。したがって Descent, Lemma
\ref{descent-lemma-descending-property-isomorphism}
により $E \to T$ は同型である。これは望みどおり $a = b$ を意味する。

\medskip\noindent
次に、二つの合成 $T' \times_T T' \to T' \to X$ が一致する射
$c : T' \to X$ を取る。$T' \to T$ との合成が $c$ である射
$a : T \to X$ を見つけなければならない。
$|U| \to |X|$ の像が $|c| : |T'| \to |X|$ の像を含むような、
アフィンスキーム $U$ とエタール射 $U \to X$ を選ぶ。
これは Lemmas \ref{spaces-properties-lemma-topology-points} および
\ref{spaces-properties-lemma-cover-by-union-affines}、アフィンの有限非交和がアフィンであること、
ならびに $|T'|$ の準コンパクト性によって可能である
（短い議論は省略する）。$U \to X$ は分離的なので
（Lemma \ref{spaces-properties-lemma-separated-cover}）、
$$
V = U \times_{X, c} T' \longrightarrow T'
$$
はスキームの全射、エタール、分離射である
（全射性には Lemma \ref{spaces-properties-lemma-points-cartesian}
と $U \to X$ の選択を用いる）。
$c \circ \text{pr}_0 = c \circ \text{pr}_1$ であることは、
次の等式により $V/T'/T$ 上の降下データが得られることを意味する
（Descent, Definition \ref{descent-definition-descent-datum}）。
\begin{align*}
V \times_{T'} (T' \times_T T')
& =
U \times_{X, c \circ \text{pr}_0} (T' \times_T T') \\
& =
(T' \times_T T') \times_{c \circ \text{pr}_1, X} U \\
& =
(T' \times_T T') \times_{T'} V
\end{align*}
More on Morphisms, Lemma
\ref{more-morphisms-lemma-etale-separated-ind-quasi-affine}
により射 $V \to T'$ は ind-準アフィンである
（エタール射は局所準有限である。Morphisms, Lemma
\ref{morphisms-lemma-etale-locally-quasi-finite} を参照）。
More on Groupoids, Lemma
\ref{more-groupoids-lemma-ind-quasi-affine}
により降下データは有効である。
$W \to T$ を、$V$ 上の与えられた降下データおよび
$T' \times_T W$ 上の標準降下データと両立する同型
$\alpha : T' \times_T W \to V$ が存在するような射とする。
このとき $W \to T$ は全射エタールである
（Descent, Lemmas
\ref{descent-lemma-descending-property-surjective} および
\ref{descent-lemma-descending-property-etale}）。
合成
$$
b' : T' \times_T W \longrightarrow V = U \times_{X, c} T' \longrightarrow U
$$
を考える。二つの合成
$b' \circ (\text{pr}_0, 1), 
b' \circ (\text{pr}_1, 1) :
(T' \times_T T') \times_T W \to T' \times_T W \to U$
は、$\alpha$ の選択と $c$ の対応する性質により一致する
（計算は省略する）。したがって Descent, Lemma
\ref{descent-lemma-fpqc-universal-effective-epimorphisms}
により $b'$ は射 $b : W \to U$ へ降下する。図式
$$
\xymatrix{
T' \times_T W \ar[r] \ar[d] & W \ar[r]_b & U \ar[d] \\
T' \ar[rr]^c &  & X
}
$$
は可換である。これは $T$ 上エタール局所的に $a$ の存在を証明した、
すなわち $a' : W \to X$ を得たということである。
しかし最初の段落で一意性を証明したので、このエタール局所解は
貼り合わせ条件を満たす。すなわち
$X(W \times_T W)$ の元として
$\text{pr}_0^*a' = \text{pr}_1^*a'$ である。
$X$ はエタール層なので、$W$ 上で $a'$ に制限される一意な
$a \in X(T)$ が存在する。
\end{proof}


\section{代数空間のエタールサイト}
\label{spaces-properties-section-etale-site}

\noindent
この節では、代数空間の小エタールサイトを定義する。
これはスキームの小エタールサイト $S_\etale$ の類似物である。
Lemma \ref{spaces-properties-lemma-etale-over-space} により、以下の定義において
$X$ のエタールサイトの対象間の任意の射はエタールであり、また
$X_\etale$ の対象上エタールな任意のスキームも $X_\etale$ の対象である。

\begin{definition}
\label{spaces-properties-definition-etale-site}
$S$ をスキームとする。
$S$ を含む大 fppf サイトを $\Sch_{fppf}$ とし、
対応する大エタールサイト（すなわち、台となる圏が同じもの）を
$\Sch_\etale$ とする。
$X$ を $S$ 上の代数空間とする。
$X$ の {\it 小エタールサイト $X_\etale$} を次のように定義する。
\begin{enumerate}
\item $X_\etale$ の対象とは射 $\varphi : U \to X$ であって、
$U \in \Ob((\Sch/S)_\etale)$ がスキームであり、
$\varphi$ がエタール射であるものをいう。
\item 射 $(\varphi : U \to X) \to (\varphi' : U' \to X)$ は、
$\varphi = \varphi' \circ \chi$ を満たすスキームの射
$\chi : U \to U'$ によって与えられる。
\item $X_\etale$ の射の族
$\{(U_i \to X) \to (U \to X)\}_{i \in I}$ が被覆であるとは、
$\{U_i \to U\}_{i \in I}$ が $(\Sch/S)_\etale$ の被覆であることをいう。
\end{enumerate}
\end{definition}

\noindent
この選択の帰結として、一般に代数空間のエタールサイトは終対象をもたない！
一方、$X$ がたまたまスキームである場合、上の定義は
Topologies, Definition \ref{topologies-definition-big-small-etale}
と一致する。

\medskip\noindent
以上が標準的に用いるサイトであるが、さらに二つの変種も用いる。
すなわち、$X$ 上エタールなすべての {\it 代数空間} $U$ を考えると、
以下で定義するサイト $X_{spaces, \etale}$ が得られる。また、
$X$ 上エタールなすべての {\it アフィンスキーム} $U$ を考えると、
以下で定義するサイト $X_{affine, \etale}$ が得られる。
前者は小エタールサイトの関手性を論じる際に用いる。
Lemma \ref{spaces-properties-lemma-functoriality-etale-site} を参照せよ。

\begin{definition}
\label{spaces-properties-definition-spaces-etale-site}
$S$ をスキームとする。
$S$ を含む大 fppf サイトを $\Sch_{fppf}$ とし、
対応する大エタールサイト（すなわち、台となる圏が同じもの）を
$\Sch_\etale$ とする。
$X$ を $S$ 上の代数空間とする。
$X$ のサイト {\it $X_{spaces, \etale}$} を次のように定義する。
\begin{enumerate}
\item $X_{spaces, \etale}$ の対象とは射 $\varphi : U \to X$ であって、
$U$ が $S$ 上の代数空間であり、$\varphi$ が $S$ 上の代数空間の
エタール射であるものをいう。
\item $X_{spaces, \etale}$ の射
$(\varphi : U \to X) \to (\varphi' : U' \to X)$ は、
$\varphi = \varphi' \circ \chi$ を満たす代数空間の射
$\chi : U \to U'$ によって与えられる。
\item $X_{spaces, \etale}$ の射の族
$\{\varphi_i : (U_i \to X) \to (U \to X)\}_{i \in I}$ が被覆であるとは、
$|U| = \bigcup \varphi_i(|U_i|)$ であることをいう。
\end{enumerate}
通常どおり、Sets, Lemma \ref{sets-lemma-coverings-site} に従って、
少なくとも $X_\etale$ の被覆を含むこの型の被覆の集合を選び、
$X_{spaces, \etale}$ をサイトとする。
\end{definition}

\noindent
$X$ の恒等射はエタールなので、$X_{spaces, \etale}$ が
終対象をもつことは明らかである。
対応するトポスが $X$ の小エタールトポスに等しいことを直ちに示そう。

\begin{lemma}
\label{spaces-properties-lemma-compare-etale-sites}
関手
$$
X_\etale \longrightarrow X_{spaces, \etale}, \quad
U/X \longmapsto U/X
$$
は特殊余連続関手
（Sites, Definition \ref{sites-definition-special-cocontinuous-functor}）
であり、したがってトポスの同値
$\Sh(X_\etale) \to \Sh(X_{spaces, \etale})$ を誘導する。
\end{lemma}

\begin{proof}
この関手が Sites, Lemma \ref{sites-lemma-equivalence} の仮定
(1) -- (5) を満たすことを示せばよい。
この関手が連続かつ余連続であることは明らかであり、
これにより仮定 (1) と (2) が従う。
この関手は充満忠実なので、仮定 (3) と (4) が成り立つ。
代数空間は定義によりスキームによる被覆をもつので、仮定 (5) も成り立つ。
\end{proof}

\begin{remark}
\label{spaces-properties-remark-explain-equivalence}
Lemma \ref{spaces-properties-lemma-compare-etale-sites} の意味を説明しよう。
$S$ をスキーム、$X$ を $S$ 上の代数空間とする。
$\mathcal{F}$ を $X$ の小エタールサイト $X_\etale$ 上の層とする。
この補題は、部分圏 $X_\etale$ への制限が $\mathcal{F}$ に戻るような
$X_{spaces, \etale}$ 上の一意な層 $\mathcal{F}'$ が存在すると述べている。
$U \to X$ が代数空間のエタール射であるとき、
$\mathcal{F}'(U)$ はどのように計算されるだろうか。
代数空間の定義により、スキーム $U'$ と全射エタール射
$U' \to U$ が存在する。すると $\{U' \to U\}$ は
$X_{spaces, \etale}$ の被覆なので、等化子図式
$$
\xymatrix{
\mathcal{F}'(U) \ar[r] &
\mathcal{F}(U') \ar@<1ex>[r] \ar@<-1ex>[r] &
\mathcal{F}(U' \times_U U').
}
$$
を得る。$U' \times_U U'$ はスキームであるから、
$\mathcal{F}'$ ではなく $\mathcal{F}$ と書けることに注意せよ。
このように、層 $\mathcal{F}$ が与えられたとき
$\mathcal{F}'$ を計算する方法が分かる。
\end{remark}

\begin{definition}
\label{spaces-properties-definition-affine-etale-site}
$S$ をスキームとする。
$S$ を含む大 fppf サイトを $\Sch_{fppf}$ とし、
対応する大エタールサイト（すなわち、台となる圏が同じもの）を
$\Sch_\etale$ とする。
$X$ を $S$ 上の代数空間とする。
$X$ のサイト {\it $X_{affine, \etale}$} を次のように定義する。
\begin{enumerate}
\item $X_{affine, \etale}$ の対象とは射 $\varphi : U \to X$ であって、
$U \in \Ob((\Sch/S)_\etale)$ がアフィンスキームであり、
$\varphi$ がエタール射であるものをいう。
\item $X_{affine, \etale}$ の射
$(\varphi : U \to X) \to (\varphi' : U' \to X)$ は、
$\varphi = \varphi' \circ \chi$ を満たすスキームの射
$\chi : U \to U'$ によって与えられる。
\item $X_{affine, \etale}$ の射の族
$\{\varphi_i : (U_i \to X) \to (U \to X)\}_{i \in I}$ が被覆であるとは、
$\{U_i \to U\}$ が標準エタール被覆であることをいう。
Topologies, Definition \ref{topologies-definition-standard-etale} を参照せよ。
\end{enumerate}
通常どおり、Sets, Lemma \ref{sets-lemma-coverings-site} に従って、
この型の被覆の集合を選び $X_{affine, \etale}$ をサイトとする。
\end{definition}

\begin{lemma}
\label{spaces-properties-lemma-alternative}
$S$ をスキーム、$X$ を $S$ 上の代数空間とする。
関手 $X_{affine, \etale} \to X_\etale$ は特殊余連続であり、
$\Sh(X_{affine, \etale})$ から $\Sh(X_\etale)$ への
トポスの同値を誘導する。
\end{lemma}

\begin{proof}
省略する。ヒント：
Topologies, Lemma \ref{topologies-lemma-affine-big-site-etale}
の証明と比較せよ。
\end{proof}

\begin{definition}
\label{spaces-properties-definition-etale-topos}
$S$ をスキーム、$X$ を $S$ 上の代数空間とする。
$X$ の {\it エタールトポス}、より正確には
$X$ の {\it 小エタールトポス} とは、$X_\etale$ 上の集合の層の圏
$\Sh(X_\etale)$ をいう。
\end{definition}

\noindent
Lemma \ref{spaces-properties-lemma-compare-etale-sites} により
$\Sh(X_\etale) = \Sh(X_{spaces, \etale})$ なので、
これを $X_{spaces, \etale}$ 上の集合の層の圏と考えることもできる。
同様に Lemma \ref{spaces-properties-lemma-alternative} により
$\Sh(X_\etale) = \Sh(X_{affine, \etale})$ である。
このトポスは代数空間の射に関して関手的であることが分かる。
正確な主張は次のとおりである。

\begin{lemma}
\label{spaces-properties-lemma-functoriality-etale-site}
$S$ をスキームとする。
$f : X \to Y$ を $S$ 上の代数空間の射とする。
\begin{enumerate}
\item 連続関手
$$
Y_{spaces, \etale} \longrightarrow X_{spaces, \etale}, \quad
V \longmapsto X \times_Y V
$$
はサイトの射
$$
f_{spaces, \etale} :
X_{spaces, \etale}
\to
Y_{spaces, \etale}.
$$
を誘導する。
\item 対応 $f \mapsto f_{spaces, \etale}$ は合成と両立する。
言い換えると $(f \circ g)_{spaces, \etale}
= f_{spaces, \etale} \circ g_{spaces, \etale}$ である
（Sites, Definition \ref{sites-definition-composition-morphisms-sites} を参照）。
\item $f_{spaces, \etale}$ に付随するトポスの射は、
Lemma \ref{spaces-properties-lemma-compare-etale-sites} を介してトポスの射
$f_{small} : \Sh(X_\etale) \to \Sh(Y_\etale)$ を誘導し、
その構成は合成と両立する。
\item $f$ が代数空間の表現可能射ならば、$f_{small}$ は
連続関手 $V \mapsto X \times_Y V$ に対応するサイトの射
$X_\etale \to Y_\etale$ から得られる。
\end{enumerate}
\end{lemma}

\begin{proof}
(1) で述べた関手が Sites, Proposition
\ref{sites-proposition-get-morphism} の仮定を満たすことを示そう。
すなわち、$Y_{spaces, \etale}$ が終対象（$Y$）をもち、
この関手がそれを $X_{spaces, \etale}$ の終対象（$X$）へ移すことを
示さなければならない。任意の圏で $X \times_Y Y = X$ なので明らかである。
次に、$Y_{spaces, \etale}$ がファイバー積をもつことを示す必要がある。
代数空間の圏はファイバー積をもち、$V$ と $V'$ が $Y$ 上エタールならば
$V \times_Y V'$ も $Y$ 上エタールなので、これは正しい
（上の Lemmas \ref{spaces-properties-lemma-composition-etale} および
\ref{spaces-properties-lemma-base-change-etale} を参照）。
したがって命題を適用でき、(1) で述べたサイトの射を得る。

\medskip\noindent
(2) は定義を展開すれば従う。
(3) は上の Lemma \ref{spaces-properties-lemma-compare-etale-sites} による
$X$ と $Y$ の同値を用いれば明らかである。
$f$ が表現可能ならば、上の関手が圏の可換図式
$$
\xymatrix{
X_\etale \ar[r] &
X_{spaces, \etale} \\
Y_\etale \ar[r] \ar[u] &
Y_{spaces, \etale} \ar[u]
}
$$
に収まるので、(4) が従う。
\end{proof}

\noindent
上の補題よりも少し詳しく、$X$ 上の層と $Y$ 上の層の関係を
記述できる。すなわち、Sheaves, Definition
\ref{sheaves-definition-f-map} と比較して、これを $f$-写像によって
次のように定式化できる。

\begin{definition}
\label{spaces-properties-definition-f-map}
$S$ をスキームとする。
$f : X \to Y$ を $S$ 上の代数空間の射とする。
$\mathcal{F}$ を $X_\etale$ 上の集合の層、
$\mathcal{G}$ を $Y_\etale$ 上の集合の層とする。
{\it $f$-写像 $\varphi : \mathcal{G} \to \mathcal{F}$} とは、
可換図式
$$
\xymatrix{
U \ar[d]_g \ar[r] & X \ar[d]^f \\
V \ar[r] & Y
}
$$
を添字とする写像の族
$\varphi_{(U, V, g)} : \mathcal{G}(V) \to \mathcal{F}(U)$ であって、
$U \in X_\etale$, $V \in Y_\etale$ であり、かつスキームのエタール射
$V' \to V$ と $U' \to U$ をもつ拡大図式
$$
\xymatrix{
U' \ar[r] \ar[d]_{g'} & U \ar[d]_g \ar[r] & X \ar[d]^f \\
V' \ar[r] & V \ar[r] & Y
}
$$
が与えられるたびに、図式
$$
\xymatrix{
\mathcal{G}(V)
\ar[rr]_{\varphi_{(U, V, g)}}
\ar[d]_{\text{制限：}\mathcal{G}} & &
\mathcal{F}(U)
\ar[d]^{\text{制限：}\mathcal{F}} \\
\mathcal{G}(V')
\ar[rr]^{\varphi_{(U', V', g')}} & &
\mathcal{F}(U')
}
$$
が可換となるものをいう。
\end{definition}

\begin{lemma}
\label{spaces-properties-lemma-f-map}
$S$ をスキームとする。
$f : X \to Y$ を $S$ 上の代数空間の射とする。
$\mathcal{F}$ を $X_\etale$ 上の集合の層、
$\mathcal{G}$ を $Y_\etale$ 上の集合の層とする。
次の三つの集合の間には標準的な全単射が存在する。
\begin{enumerate}
\item 写像 $\mathcal{G} \to f_{small, *}\mathcal{F}$ の集合。
\item 写像 $f_{small}^{-1}\mathcal{G} \to \mathcal{F}$ の集合。
\item $f$-写像 $\varphi : \mathcal{G} \to \mathcal{F}$ の集合。
\end{enumerate}
\end{lemma}

\begin{proof}
$f_{small, *}$ と $f_{small}^{-1}$ は随伴関手の対なので、
(1) と (2) は同じものであることに注意する。
$\alpha : f_{small}^{-1}\mathcal{G} \to \mathcal{F}$ が
$Y_\etale$ 上の層の写像であると仮定する。
Definition \ref{spaces-properties-definition-f-map} に現れる図式
$$
\xymatrix{
U \ar[d]_g \ar[r]_{j_U} & X \ar[d]^f \\
V \ar[r]^{j_V} & Y
}
$$
が与えられたとする。この図式の可換性により写像
$g_{small}^{-1}(j_V)^{-1}\mathcal{G} \to (j_U)^{-1}\mathcal{F}$
も得られる（局所化関手の記述については Sites, Section
\ref{sites-section-localize} と比較せよ）。したがって写像
$\varphi_{(V, U, g)} :
\mathcal{G}(V) = (j_V)^{-1}\mathcal{G}(V)
\to
(j_U)^{-1}\mathcal{F}(U) = \mathcal{F}(U)$
が得られる。この規則がさらなる制限と両立し、
$\mathcal{G}$ から $\mathcal{F}$ への $f$-写像を定めることの確認は省略する。

\medskip\noindent
逆に、$f$-写像 $\varphi = (\varphi_{(U, V, g)})$ が与えられたとする。
$\mathcal{G}$（それぞれ $\mathcal{F}$）の
$Y_{spaces, \etale}$（それぞれ $X_{spaces, \etale}$）への拡張を
$\mathcal{G}'$（それぞれ $\mathcal{F}'$）と書く。
Lemma \ref{spaces-properties-lemma-compare-etale-sites} を参照せよ。
層の写像
$$
\mathcal{G}' \longrightarrow (f_{spaces, \etale})_*\mathcal{F}'
$$
を構成しなければならない。そのため、$V \to Y$ を代数空間の
エタール射とする。集合の写像
$$
\mathcal{G}'(V) \to \mathcal{F}'(X \times_Y V)
$$
を構成する必要がある。
$V'$ がスキームであるような全射エタール射 $V' \to V$ を選び、
次いで $U'$ がスキームであるような全射エタール射
$U' \to X \times_U V'$ を選ぶ。スキームの射 $g' : U' \to V'$ と、
さらにスキームの射
$$
g'' : U' \times_{X \times_Y V} U' \longrightarrow V' \times_V V'
$$
を得る。次の図式を考える。
$$
\xymatrix{
\mathcal{F}'(X \times_Y V) \ar[r] &
\mathcal{F}(U') \ar@<1ex>[r] \ar@<-1ex>[r] &
\mathcal{F}(U' \times_{X \times_Y V} U') \\
\mathcal{G}'(X \times_Y V) \ar[r] \ar@{..>}[u] &
\mathcal{G}(V') \ar@<1ex>[r] \ar@<-1ex>[r] \ar[u]_{\varphi_{(U', V', g')}} &
\mathcal{G}(V' \times_V V') \ar[u]_{\varphi_{(U'', V'', g'')}}
}
$$
写像 $\varphi_{...}$ と制限との両立性により、右側の二つの正方形は可換である。
$X_{spaces, \etale}$ における被覆の定義により、横の二行は等化子図式である。
したがって点線の矢印が得られる。これらの矢印が制限写像と両立することの
確認は読者に委ねる。
\end{proof}

\noindent
代数空間の射 $X \to Y$ がエタールならば、トポスの射
$\Sh(X_\etale) \to \Sh(Y_\etale)$ は局所化である。
次がその主張である。

\begin{lemma}
\label{spaces-properties-lemma-etale-morphism-topoi}
$S$ をスキームとし、$f : X \to Y$ を $S$ 上の代数空間の射とする。
$f$ がエタールであると仮定する。この場合、余連続な関手
$$
j : X_\etale \to Y_\etale, \quad
(\varphi : U \to X) \mapsto (f \circ \varphi : U \to Y)
$$
が存在する。トポスの射 $f_{small}$ は $j$ に付随するトポスの射である。
Sites, Lemma \ref{sites-lemma-cocontinuous-morphism-topoi} を参照せよ。
さらに $j$ は連続でもあるので、Sites, Lemma
\ref{sites-lemma-when-shriek} が適用される。特に、
$Y_\etale$ 上のすべての層 $\mathcal{G}$ に対して
$f_{small}^{-1}\mathcal{G}(U) = \mathcal{G}(jU)$ である。
\end{lemma}

\begin{proof}
代数空間のエタール射の定義（Definition \ref{spaces-properties-definition-etale}）そのものにより、
与えられた規則は実際に表示された関手 $j$ を定める。
$j(\varphi : U \to X)$ の $Y_\etale$ における被覆
$\{U_i \to U\}$ は $(\varphi : U \to X)$ の $X_\etale$ における
被覆と同じものなので、$j$ が余連続かつ連続であることは明らかである。
$j$ が $f_{small}$ と同じトポスの射を誘導することを示すことだけが残る。
そのため圏の図式
$$
\xymatrix{
X_\etale \ar[r] \ar[d]^j &
X_{spaces, \etale} \ar@/_/[d]_{j_{spaces}} \\
Y_\etale \ar[r] &
Y_{spaces, \etale} \ar@/_/[u]_{v : V \mapsto X \times_Y V}
}
$$
を考える。ここで関手 $j_{spaces}$ は $j$ の圏
$X_{spaces, \etale}$ への明らかな拡張である。
したがって内側の正方形は可換である。実際、$j_{spaces}$ は
Sites, Section \ref{sites-section-localize} で論じた局所化関手
$j_X : Y_{spaces, \etale}/X \to Y_{spaces, \etale}$ と同一視できる。
ゆえに Sites, Lemma \ref{sites-lemma-localize-given-products} により、
余連続関手 $j_{spaces}$ と図式の関手 $v$ は同じトポスの射を誘導する。
Sites, Lemma \ref{sites-lemma-composition-cocontinuous} により、
内側の正方形（サイト間の余連続関手からなる）の可換性は、
付随するトポスの射の可換図式を与える。
したがって Lemma \ref{spaces-properties-lemma-functoriality-etale-site} における
$f_{small}$ の構成により結論を得る。
\end{proof}

\noindent
上の補題は、代数空間のエタール射 $f : X \to Y$ による
$\mathcal{G}$ の引き戻しが、単に $\mathcal{G}$ の圏 $X_\etale$ への
制限であることを述べている。これを表すため、しばしば略記
\begin{equation}
\label{spaces-properties-equation-restrict}
\mathcal{G}|_{X_\etale} = f_{small}^{-1}\mathcal{G}
\end{equation}
を用いる。この状況で補題の関手 $j : X_\etale \to Y_\etale$ は
忠実であるが、一般には充満忠実ではないことに注意せよ。
これについては Section \ref{spaces-properties-section-localize} でより技術的に論じる。

\begin{lemma}
\label{spaces-properties-lemma-pushforward-etale-base-change}
$S$ をスキームとする。$S$ 上の代数空間のカルテジアン図式
$$
\xymatrix{
X' \ar[r] \ar[d]_{f'} & X \ar[d]^f \\
Y' \ar[r]^g & Y
}
$$
を考える。$\mathcal{F}$ を $X_\etale$ 上の層とする。
$g$ がエタールならば、
\begin{enumerate}
\item $\Sh(Y'_\etale)$ において
$f'_{small, *}(\mathcal{F}|_{X'}) = (f_{small, *}\mathcal{F})|_{Y'}$
である\footnote{図式の可換性と (\ref{spaces-properties-equation-restrict}) により、
$(f')_{small}^{-1}(\mathcal{G}|_{Y'}) =
(f_{small}^{-1}\mathcal{G})|_{X'}$ でもある。}。
\item $\mathcal{F}$ がアーベル層ならば、
$R^if'_{small, *}(\mathcal{F}|_{X'}) = (R^if_{small, *}\mathcal{F})|_{Y'}$ である。
\end{enumerate}
\end{lemma}

\begin{proof}
次の関手の図式を考える。
$$
\xymatrix{
X'_{spaces, \etale} \ar[r]_j &
X_{spaces, \etale} \\
Y'_{spaces, \etale} \ar[r]^j \ar[u]^{V' \mapsto V' \times_{Y'} X'} &
Y_{spaces, \etale} \ar[u]_{V \mapsto V \times_Y X}
}
$$
横の矢印は局所化であり、縦の矢印はサイトの射を誘導する。
したがって Sites, Lemma \ref{sites-lemma-localize-morphism} の
最後の主張から (1) が従う。(2) を示すには、(1) を $\mathcal{F}$ の
単射分解に適用し、制限が完全で単射対象を保つことを用いる
（Cohomology on Sites, Lemma
\ref{sites-cohomology-lemma-cohomology-of-open} を参照）。
\end{proof}

\noindent
次の補題は、代数空間の小エタールサイト上の層を、
その空間上エタールなスキームの小エタールサイト上の層からなる
両立する族と考えられることを述べている。
補題の比較写像 $c_f$ はすべて同型であることに注意せよ。
これは Topologies, Lemma
\ref{topologies-lemma-characterize-sheaf-big-etale} および、
$X_\etale$ の対象間のすべての射がエタールであることと両立する。

\begin{lemma}
\label{spaces-properties-lemma-characterize-sheaf-small-etale}
$S$ をスキーム、$X$ を $S$ 上の代数空間とする。
$X_\etale$ 上の層 $\mathcal{F}$ は次のデータによって与えられる。
\begin{enumerate}
\item 各 $U \in \Ob(X_\etale)$ に対する $U_\etale$ 上の層
$\mathcal{F}_U$。
\item $X_\etale$ の各 $f : U' \to U$ に対する同型
$c_f : f_{small}^{-1}\mathcal{F}_U \to \mathcal{F}_{U'}$。
\end{enumerate}
これらのデータには、$X_\etale$ の任意の $f : U' \to U$ と
$g : U'' \to U'$ に対して、合成 $c_g \circ g_{small}^{-1} c_f$ が
$c_{f \circ g}$ に等しいという条件を課す。
\end{lemma}

\begin{proof}
$g_{small}^{-1}$ は Lemma \ref{spaces-properties-lemma-etale-morphism-topoi} のように
解釈できる。このとき主張はサイトに関する一般的事実
Sites, Lemma \ref{sites-lemma-glue-sheaves-absolute} から従う。
\end{proof}

\noindent
$S$ をスキーム、$X$ を $S$ 上の代数空間とする。
$\varphi : U \to X$ を任意の全射エタール射とし、
$X = U/R$ をそれから得られる $X$ の表示とする。
Spaces, Definition \ref{spaces-definition-presentation} を参照せよ。
特に、$j = (t, s) : R \to U \times_S U$ となる亜群
$(U, R, s, t, c, e, i)$ を得る。
Groupoids, Lemma \ref{groupoids-lemma-equivalence-groupoid} を参照せよ。

\begin{lemma}
\label{spaces-properties-lemma-descent-sheaf}
$S$, $\varphi : U \to X$, $(U, R, s, t, c, e, i)$ は上のとおりとする。
$X_\etale$ 上の任意の層 $\mathcal{F}$ に対して、層\footnote{この補題と
その証明では、$\varphi_{small}^{-1}$ の代わりに単に $\varphi^{-1}$ と書き、
他のすべての引き戻しについても同様にする。}
$\mathcal{G} = \varphi^{-1}\mathcal{F}$ には標準的な同型
$$
\alpha :
t^{-1}\mathcal{G}
\longrightarrow
s^{-1}\mathcal{G}
$$
が備わり、図式
$$
\xymatrix{
& \text{pr}_1^{-1}t^{-1}\mathcal{G} \ar[r]_-{\text{pr}_1^{-1}\alpha} &
\text{pr}_1^{-1}s^{-1}\mathcal{G} \ar@{=}[rd] & \\
\text{pr}_0^{-1}s^{-1}\mathcal{G} \ar@{=}[ru] & & &
c^{-1}s^{-1}\mathcal{G} \\
&
\text{pr}_0^{-1}t^{-1}\mathcal{G} \ar[lu]^{\text{pr}_0^{-1}\alpha} \ar@{=}[r] &
c^{-1}t^{-1}\mathcal{G} \ar[ru]_{c^{-1}\alpha}
}
$$
は可換である。関手 $\mathcal{F} \mapsto (\mathcal{G}, \alpha)$ は、
$X_\etale$ 上の層の圏と上のような対 $(\mathcal{G}, \alpha)$ の圏との
同値を定める。
\end{lemma}

\begin{proof}[Lemma \ref{spaces-properties-lemma-descent-sheaf} の第一の証明]
$\mathcal{C} = X_{spaces, \etale}$ とおく。
Lemma \ref{spaces-properties-lemma-etale-morphism-topoi} とその証明により
$U_{spaces, \etale} = \mathcal{C}/U$ であり、引き戻し関手
$\varphi^{-1}$ は単に制限関手である。
さらに $\{U \to X\}$ はサイト $\mathcal{C}$ の被覆であり、
$R = U \times_X U$ である。同型 $\alpha$ は単に標準的な同一視
$$
\left(\mathcal{F}|_{\mathcal{C}/U}\right)|_{\mathcal{C}/U \times_X U}
=
\left(\mathcal{F}|_{\mathcal{C}/U}\right)|_{\mathcal{C}/U \times_X U}
$$
であり、図式の可換性は貼り合わせデータのコサイクル条件である。
したがってこの補題は層の貼り合わせの特殊な場合である。
Sites, Section \ref{sites-section-glueing-sheaves} を参照せよ。
\end{proof}

\begin{proof}[Lemma \ref{spaces-properties-lemma-descent-sheaf} の第二の証明]
$\alpha$ の存在は $\varphi \circ t = \varphi \circ s$ と、
引き戻しが射に関して関手的であることから従う。
Lemma \ref{spaces-properties-lemma-functoriality-etale-site} を参照せよ。
まったく同様に、すなわち引き戻しの関手性により、
同型 $\alpha$ が可換図式に収まることが分かる。
構成 $\mathcal{F} \mapsto (\varphi^{-1}\mathcal{F}, \alpha)$ は
層 $\mathcal{F}$ に関して明らかに関手的である。よって関手を得る。

\medskip\noindent
逆に、対 $(\mathcal{G}, \alpha)$ が与えられたとする。
$V \to X$ を $X_\etale$ の対象とする。
この場合、射 $V' = U \times_X V \to V$ はスキームの全射エタール射なので、
$\{V' \to V\}$ は $V$ のエタール被覆である。
$\mathcal{G}' = (V' \to V)^{-1}\mathcal{G}$ とおく。
$R = U \times_X U$ かつ $t = \text{pr}_0$, $s = \text{pr}_0$ なので、
$V' \times_V V' = R \times_X V$ であり、射影
$s', t' : V' \times_V V' \to V'$ は $t$ と $s$ の引き戻しに等しい。
したがって $\alpha$ は同型
$\alpha' : (t')^{-1}\mathcal{G}' \to (s')^{-1}\mathcal{G}'$ に引き戻される。
以上を踏まえて、単に
$$
\xymatrix{
\mathcal{F}(V) \ar@{=}[r] &
\text{Equalizer}(\mathcal{G}(V') \ar@<1ex>[r] \ar@<-1ex>[r] &
\mathcal{G}(V' \times_V V').
}
$$
と定義する。これが層を定めることの確認は省略する。
射 $V \to U$ が存在するとき $\mathcal{G}(V) = \mathcal{F}(V)$ であることは、
この場合、等化子が
$H^0(\{V' \to V\}, \mathcal{G}) = \mathcal{G}(V)$ であることから分かる。
\end{proof}

\section{小エタールサイトの点}
\label{spaces-properties-section-points-small-etale-site}

\noindent
この節は 'Etale Cohomology, Section
\ref{etale-cohomology-section-stalks} の類似物である。

\begin{definition}
\label{spaces-properties-definition-geometric-point}
$S$ をスキーム、$X$ を $S$ 上の代数空間とする。
\begin{enumerate}
\item $X$ の {\it 幾何学的点} とは射
$\overline{x} : \Spec(k) \to X$ であって、$k$ が代数閉体であるものをいう。
しばしば記号を濫用して $\overline{x} = \Spec(k)$ と書く。
\item 各幾何学的点 $\overline{x}$ には、対応する「像」の点
$x \in |X|$ がある。$\overline{x}$ を
{\it $x$ の上にある幾何学的点} と呼ぶ。
\end{enumerate}
\end{definition}

\noindent
スキームの小エタールサイトの場合とまったく同じ方法で、
$X_\etale$ 上の層の幾何学的点における茎を取れることが分かる。
そのため、エタール近傍の概念を次のように定義する。

\begin{definition}
\label{spaces-properties-definition-etale-neighbourhood}
$S$ をスキーム、$X$ を $S$ 上の代数空間とする。
$\overline{x}$ を $X$ の幾何学的点とする。
\begin{enumerate}
\item $X$ における $\overline{x}$ の {\it エタール近傍} とは可換図式
$$
\xymatrix{
& U \ar[d]^\varphi \\
{\bar x} \ar[r]^{\bar x} \ar[ur]^{\bar u} & X
}
$$
であって、$\varphi$ が $S$ 上の代数空間のエタール射であるものをいう。
この状況を表すため、記法
$\varphi : (U, \overline{u}) \to (X, \overline{x})$ を用いる。
\item {\it エタール近傍の射}
$(U, \overline{u}) \to (U', \overline{u}')$ とは、
$\overline{u}' = h \circ \overline{u}$ を満たす $X$-射
$h : U \to U'$ をいう。
\end{enumerate}
\end{definition}

\noindent
ここでは $U$ が代数空間であることを許している。
$X_\etale$ 上の層の茎を取るときは $X_\etale$ に属する $U$ に制限する
必要があるため、その場合は $U$ がスキームの場合だけを考える。
あるいはサイト $X_{space, \etale}$ を用いて、すべてのエタール近傍を
考えることもできる。次の補題の最後の主張により違いは生じない。

\begin{lemma}
\label{spaces-properties-lemma-cofinal-etale}
$S$ をスキーム、$X$ を $S$ 上の代数空間とする。
$\overline{x}$ を $X$ の幾何学的点とする。
エタール近傍の圏は余フィルター圏である。より正確には次が成り立つ。
\begin{enumerate}
\item $(U_i, \overline{u}_i)_{i = 1, 2}$ を $X$ における
$\overline{x}$ の二つのエタール近傍とする。このとき第三のエタール近傍
$(U, \overline{u})$ と射
$(U, \overline{u}) \to (U_i, \overline{u}_i)$, $i = 1, 2$ が存在する。
\item $h_1, h_2: (U, \overline{u}) \to (U', \overline{u}')$ を
$\overline{s}$ のエタール近傍間の二つの射とする。このときエタール近傍
$(U'', \overline{u}'')$ と射
$h : (U'', \overline{u}'') \to (U, \overline{u})$ であって、
$h_1$ と $h_2$ を等化するもの、すなわち
$h_1 \circ h = h_2 \circ h$ を満たすものが存在する。
\end{enumerate}
さらに、任意のエタール近傍
$(U, \overline{u}) \to (X, \overline{x})$ に対して、
$U'$ がスキームであるようなエタール近傍の射
$(U', \overline{u}') \to (U, \overline{u})$ が存在する。
\end{lemma}

\begin{proof}
(1) について、ファイバー積 $U = U_1 \times_X U_2$ を考える。
エタール射は基底変換と合成で保たれるので、これは $U_1$ と $U_2$ の
いずれの上でもエタールである。Lemmas \ref{spaces-properties-lemma-base-change-etale}
および \ref{spaces-properties-lemma-composition-etale} を参照せよ。
$(\overline{u}_1, \overline{u}_2)$ によって定まる写像
$\overline{u} \to U$ は、それに $U_1$ と $U_2$ の両方へ写る
エタール近傍の構造を与える。

\medskip\noindent
(2) について、$U''$ をファイバー積
$$
\xymatrix{
U'' \ar[r] \ar[d] & U \ar[d]^{(h_1, h_2)} \\
U' \ar[r]^-\Delta & U' \times_X U'.
}
$$
として定義する。$\overline{u}$ と $\overline{u}'$ は $X$ 上で
$\overline{x}$ と一致するので、$\overline{u}'' = (\overline{u}, \overline{u}')$
は $U''$ の幾何学的点である。特に $U'' \not = \emptyset$ である。
さらに $U'$ は $X$ 上エタールなので、ファイバー積 $U'\times_X U'$ も
同様にエタールである（上で $U_1 \times_X U_2$ の場合に見たとおり）。
したがって Lemma \ref{spaces-properties-lemma-etale-permanence} により縦の矢印
$(h_1, h_2)$ はエタールである。ゆえに $U''$ は基底変換により
$U'$ 上エタールであり、したがって $X$ 上でもエタールである
（エタール射の合成はエタールである）。
よって $(U'', \overline{u}'')$ は (2) で課された問題の解である。

\medskip\noindent
最後の主張を示すため、$U'$ がスキームであるような任意の全射エタール射
$U' \to U$ を選ぶ。このとき $U' \times_U \overline{u}$ はスキームであり、
$k$ が代数閉体である $\overline{u} = \Spec(k)$ 上全射かつエタールである。
したがって（Morphisms, Lemma
\ref{morphisms-lemma-etale-over-field} を参照）
$U' \times_U \overline{u} \to \overline{u}$ は切断をもち、
これが所望の $\overline{u}'$ を与える。
\end{proof}

\begin{lemma}
\label{spaces-properties-lemma-geometric-lift-to-usual}
$S$ をスキーム、$X$ を $S$ 上の代数空間とする。
$\overline{x} : \Spec(k) \to X$ を $x \in |X|$ の上にある
$X$ の幾何学的点とする。$\varphi : U \to X$ を代数空間のエタール射とし、
$u \in |U|$ は $\varphi(u) = x$ を満たすとする。
このとき $u$ の上にある幾何学的点
$\overline{u} : \Spec(k) \to U$ であって、
$\overline{x} = \varphi \circ \overline{u}$ を満たすものが存在する。
\end{lemma}

\begin{proof}
$u' \in U'$ をもつアフィンスキーム $U'$ と、$u'$ を $u$ へ写す
エタール射 $U' \to U$ を選ぶ。
$(U', u') \to (X, x)$ に対して補題を証明できれば結論が従う。
したがって $U$ がスキームであり、特に $U \to X$ が表現可能であると
仮定してよい。カルテジアン図式
$$
\xymatrix{
\Spec(k) \times_{\overline{x}, X, \varphi} U
\ar[d]_{\text{pr}_1} \ar[r]_-{\text{pr}_2} & U
\ar[d]^\varphi \\
\Spec(k) \ar[r]^-{\overline{x}} & X
}
$$
を考える。射影 $\text{pr}_1$ はエタール射の基底変換なのでエタールである。
Lemma \ref{spaces-properties-lemma-base-change-etale} を参照せよ。
したがってスキーム $\Spec(k) \times_{\overline{x}, X, \varphi} U$ は、
$k$ の有限分離拡大の非交和である。
Morphisms, Lemma \ref{morphisms-lemma-etale-over-field} を参照せよ。
しかし $k$ は代数閉なので、これらの拡大はすべて自明である。
よって $\Spec(k) \times_{\overline{x}, X, \varphi} U$ は
$\Spec(k)$ のコピーの非交和であり、その各々は
$\varphi \circ \overline{u} = \overline{x}$ を満たす幾何学的点
$\overline{u}$ に対応する。Lemma \ref{spaces-properties-lemma-points-cartesian} により写像
$$
|\Spec(k) \times_{\overline{x}, X, \varphi} U|
\longrightarrow
|\Spec(k)| \times_{|X|} |U|
$$
は全射なので、$u$ の上にある $\overline{u}$ を選べる。
\end{proof}

\begin{lemma}
\label{spaces-properties-lemma-geometric-lift-to-cover}
$S$ をスキーム、$X$ を $S$ 上の代数空間とする。
$\overline{x}$ を $X$ の幾何学的点とする。
$(U, \overline{u})$ を $\overline{x}$ のエタール近傍とする。
$\{\varphi_i : U_i \to U\}_{i \in I}$ を
$X_{spaces, \etale}$ におけるエタール被覆とする。
このとき $i \in I$ と $\overline{u}_i : \overline{x} \to U_i$ であって、
$\varphi_i : (U_i, \overline{u}_i) \to (U, \overline{u})$ が
エタール近傍の射となるものが存在する。
\end{lemma}

\begin{proof}
$u \in |U|$ を $\overline{u}$ の像とする。
$|U| = \bigcup_{i \in I} \varphi_i(|U_i|)$ なので、$x$ へ写る点
$u_i \in U_i$ をもつ $i$ が存在する。
$(U_i, u_i) \to (U, u)$ と $\overline{u}$ に
Lemma \ref{spaces-properties-lemma-geometric-lift-to-usual} を適用すれば、
所望の幾何学的点を得る。
\end{proof}

\begin{definition}
\label{spaces-properties-definition-stalk}
$S$ をスキーム、$X$ を $S$ 上の代数空間とする。
$\mathcal{F}$ を $X_\etale$ 上の前層とする。
$\overline{x}$ を $X$ の幾何学的点とする。
$\overline{x}$ における $\mathcal{F}$ の {\it 茎} とは
$$
\mathcal{F}_{\bar x}
=
\colim_{(U, \overline{u})} \mathcal{F}(U)
$$
をいう。ここで $(U, \overline{u})$ は、
$U \in \Ob(X_\etale)$ である $X$ における
$\overline{x}$ のすべてのエタール近傍を走る。
\end{definition}

\noindent
Lemma \ref{spaces-properties-lemma-cofinal-etale} により、この余極限はフィルター付き添字圏、
すなわち $X_\etale$ におけるエタール近傍の圏の反対圏を走る。
より正確には、Lemma \ref{spaces-properties-lemma-cofinal-etale} はすべてのエタール近傍の
圏の反対圏がフィルター圏であることを述べ、$X_\etale$ に属するものからなる
充満部分圏は共終部分圏であるから、その反対圏もフィルター圏である。

\medskip\noindent
したがって $\mathcal{F}_{\overline{x}}$ の元は三つ組
$(U, \overline{u}, \sigma)$ と考えられる。ここで
$U \in \Ob(X_\etale)$ かつ $\sigma \in \mathcal{F}(U)$ である。
二つの三つ組 $(U, \overline{u}, \sigma)$,
$(U', \overline{u}', \sigma')$ が同じ茎の元を定めるのは、
$U'' \in \Ob(X_\etale)$ である第三のエタール近傍
$(U'', \overline{u}'')$ とエタール近傍の射
$h : (U'', \overline{u}'') \to (U, \overline{u})$,
$h' : (U'', \overline{u}'') \to (U', \overline{u}')$ が存在し、
$\mathcal{F}(U'')$ において $h^*\sigma = (h')^*\sigma'$ が成り立つ場合である。
Categories, Section \ref{categories-section-directed-colimits} を参照せよ。

\medskip\noindent
また、$\mathcal{F}'$ が $X_\etale$ 上の $\mathcal{F}$ に対応する
$X_{spaces, \etale}$ 上の層ならば、
\begin{equation}
\label{spaces-properties-equation-stalk-spaces-etale}
\mathcal{F}_{\overline{x}} = \colim_{(U, \overline{u})} \mathcal{F}'(U)
\end{equation}
である。ここでは余極限は $\overline{x}$ のすべてのエタール近傍を走る。
今後、特に断らず $X_\etale$ を用いる観点と
$X_{spaces, \etale}$ を用いる観点との間を行き来する。

\medskip\noindent
特に、$\mathcal{F}$ がアーベル群、環などの前層ならば、
$\mathcal{F}_{\overline{x}}$ は、アーベル群、環などの有向余極限に
群構造を入れる通常の方法によって、アーベル群、環などになる。

\begin{lemma}
\label{spaces-properties-lemma-stalk-gives-point}
\begin{slogan}
代数空間の幾何学的点は、そのエタールトポスの点を与える。
\end{slogan}
$S$ をスキームとする。
$X$ を $S$ 上の代数空間とする。
$\overline{x}$ を $X$ の幾何学的点とする。
関手
$$
u : X_\etale \longrightarrow \textit{Sets}, \quad
U \longmapsto |U_{\overline{x}}|
$$
を考える。このとき $u$ はサイト $X_\etale$ の点 $p$ を定め
（Sites, Definition \ref{sites-definition-point}）、それに付随する茎関手
$\mathcal{F} \mapsto \mathcal{F}_p$
（Sites, Equation \ref{sites-equation-stalk}）は上で定義した関手
$\mathcal{F} \mapsto \mathcal{F}_{\overline{x}}$ である。
\end{lemma}

\begin{proof}
Lemma \ref{spaces-properties-lemma-geometric-lift-to-usual} の証明で、スキーム
$U_{\overline{x}} = \overline{x} \times_X U$ は
$\overline{x}$ と同型なスキームの非交和であることを見た。
したがって $|U_{\overline{x}}|$ は $\overline{x}$ の上にある
$U$ の幾何学的点の集合、すなわち Definition
\ref{spaces-properties-definition-etale-neighbourhood} の図式に収まる射
$\overline{u} : \overline{x} \to U$ の集まりとも考えられる。
これより $u(X)$ は一元集合であり、$U \to V$ と $W \to V$ が
$X_\etale$ の射であるとき
$u(U \times_V W) = u(U) \times_{u(V)} u(W)$ である。
さらに $X_\etale$ の被覆 $\{U_i \to U\}_{i \in I}$ が与えられれば、
Lemma \ref{spaces-properties-lemma-geometric-lift-to-cover} により
$\coprod u(U_i) \to u(U)$ は全射である。
したがって Sites, Proposition \ref{sites-proposition-point-limits} を
適用でき、$p$ はサイト $X_\etale$ の点である。
最後に、関手 $\mathcal{F} \mapsto \mathcal{F}_{\overline{s}}$ は、
Sites, Equation \ref{sites-equation-stalk} において $p$ に付随する関手
$\mathcal{F} \mapsto \mathcal{F}_p$ とまったく同じ余極限によって
与えられる。これが最後の主張を証明する。
\end{proof}

\begin{lemma}
\label{spaces-properties-lemma-stalk-exact}
$S$ をスキームとする。
$X$ を $S$ 上の代数空間とする。
$\overline{x}$ を $X$ の幾何学的点とする。
\begin{enumerate}
\item 茎関手
$\textit{PAb}(X_\etale) \to \textit{Ab}$,
$\mathcal{F}  \mapsto  \mathcal{F}_{\overline{x}}$
は完全である。
\item $X_\etale$ 上の任意の集合の前層 $\mathcal{F}$ に対して
$(\mathcal{F}^\#)_{\overline{x}} = \mathcal{F}_{\overline{x}}$ である。
\item 関手
$\textit{Ab}(X_\etale) \to \textit{Ab}$,
$\mathcal{F} \mapsto \mathcal{F}_{\overline{x}}$ は完全である。
\item 同様に、茎関手 $\mathcal{F} \mapsto \mathcal{F}_{\overline{x}}$ による関手
$\textit{PSh}(X_\etale) \to \textit{Sets}$ および
$\Sh(X_\etale) \to \textit{Sets}$ は完全であり
（Categories, Definition \ref{categories-definition-exact} を参照）、
任意の余極限と可換である。
\end{enumerate}
\end{lemma}

\begin{proof}
この結果は Modules on Sites, Section
\ref{sites-modules-section-stalks} の一般論から従う。
実際、$\mathcal{F} \mapsto \mathcal{F}_{\overline{x}}$ は $X$ の
小エタールサイトの点から得られる。
Lemma \ref{spaces-properties-lemma-stalk-gives-point} を参照せよ。
スキームの小エタールサイトの場合におけるこれらの主張の一部の
直接的な証明については、'Etale Cohomology, Lemma
\ref{etale-cohomology-lemma-stalk-exact} の証明を参照せよ。
\end{proof}

\noindent
後で、茎関手 $\mathcal{F} \mapsto \mathcal{F}_{\overline{x}}$ が
実際には射 $\overline{x}$ に沿う引き戻しであることを見る。
その意味で、次の補題は上の補題の一般化である。

\begin{lemma}
\label{spaces-properties-lemma-stalk-pullback}
$S$ をスキームとする。
$f : X \to Y$ を $S$ 上の代数空間の射とする。
\begin{enumerate}
\item 関手
$f_{small}^{-1} :
\textit{Ab}(Y_\etale)
\to
\textit{Ab}(X_\etale)$
は完全である。
\item 関手
$f_{small}^{-1} :
\Sh(Y_\etale)
\to
\Sh(X_\etale)$
は完全、すなわち有限極限および有限余極限と可換である。
Categories, Definition \ref{categories-definition-exact} を参照せよ。
\item 代数空間の任意のエタール射 $V \to Y$ に対して
$f_{small}^{-1}h_V = h_{X \times_Y V}$ である。
\item $\overline{x} \to X$ を幾何学的点とする。
$\mathcal{G}$ を $Y_\etale$ 上の層とする。
このとき標準的な同一視
$$
(f_{small}^{-1}\mathcal{G})_{\overline{x}} = \mathcal{G}_{\overline{y}}.
$$
が存在する。ここで $\overline{y} = f \circ \overline{x}$ である。
\end{enumerate}
\end{lemma}

\begin{proof}
Lemma \ref{spaces-properties-lemma-functoriality-etale-site} において $f_{small}$ は
$f_{spaces, small}$ を介して定義されたことを思い出す。
(1), (2), (3) は、
$f_{spaces, \etale} :
X_{spaces, \etale}
\to
Y_{spaces, \etale}$
がサイトの射であることの一般的な帰結である。
(2) については Sites, Definition
\ref{sites-definition-morphism-sites}、(1) については
Modules on Sites, Lemma
\ref{sites-modules-lemma-flat-pullback-exact}、(3) については
Sites, Lemma \ref{sites-lemma-pullback-representable-sheaf} を参照せよ。

\medskip\noindent
(4) の証明。これは Sites, Lemma
\ref{sites-lemma-point-morphism-sites} の、
Lemma \ref{spaces-properties-lemma-stalk-gives-point} を介した特殊な場合である。
直接的な証明も与える。Lemma \ref{spaces-properties-lemma-stalk-exact} により、
茎を取ることは層化と可換であることに注意する。
$Y_\etale$ への制限が $\mathcal{G}$ であるような
$Y_{spaces, \etale}$ 上の層を $\mathcal{G}'$ とする。
$f_{spaces, \etale}^{-1}\mathcal{G}'$ は前層
$$
U \longrightarrow \colim_{U \to X \times_Y V} \mathcal{G}'(V),
$$
に付随する層であることを思い出す。
Sites, Sections \ref{sites-section-continuous-functors} および
\ref{sites-section-functoriality-PSh} を参照せよ。
したがって
\begin{align*}
(f_{spaces, \etale}^{-1}\mathcal{G}')_{\overline{x}}
& =
\colim_{(U, \overline{u})} f_{spaces, \etale}^{-1}\mathcal{G}'(U)
\\
& = \colim_{(U, \overline{u})}
\colim_{a : U \to X \times_Y V} \mathcal{G}'(V) \\
& = \colim_{(V, \overline{v})} \mathcal{G}'(V) \\
& = \mathcal{G}'_{\overline{y}}
\end{align*}
である。第三の等式において、対 $(U, \overline{u})$ と写像
$a : U \to X \times_Y V$ は対 $(V, a \circ \overline{u})$ に対応する。
$\mathcal{G}'$（それぞれ $f_{spaces, \etale}^{-1}\mathcal{G}'$）の茎は、
$\mathcal{G}$（それぞれ $f_{small}^{-1}\mathcal{G}$）の茎と一致する。
Equation (\ref{spaces-properties-equation-stalk-spaces-etale}) を参照せよ。これで結論が従う。
\end{proof}

\begin{remark}
\label{spaces-properties-remark-stalk-pullback}
この注意は 'Etale Cohomology, Remark
\ref{etale-cohomology-remark-stalk-pullback} の類似物である。
$S$ をスキームとする。
$X$ を $S$ 上の代数空間とする。
$\overline{x} : \Spec(k) \to X$ を $X$ の幾何学的点とする。
'Etale Cohomology, Theorem
\ref{etale-cohomology-theorem-equivalence-sheaves-point} により、
$\Spec(k)_\etale$ 上の層の圏は集合の圏と同値である
（層をその大域切断へ写すことによる）。
したがって、射 $\overline{x}$ に適用した
Lemma \ref{spaces-properties-lemma-stalk-pullback} の (4) から、関手
$$
\Sh(X_\etale) \longrightarrow \textit{Sets}, \quad
\mathcal{F} \longmapsto \mathcal{F}_{\overline{x}}
$$
は関手
$$
\Sh(X_\etale)
\longrightarrow
\Sh(\Spec(k)_\etale) = \textit{Sets},
\quad
\mathcal{F} \longmapsto \overline{x}^*\mathcal{F}
$$
と同型である。したがって茎関手は、
（Lemma \ref{spaces-properties-lemma-stalk-gives-point} の結果における単なる抽象的な
トポスの射ではなく）幾何学的射に沿う引き戻し関手とみなせる。
\end{remark}

\begin{remark}
\label{spaces-properties-remark-map-stalks}
$S$ をスキームとする。
$X$ を $S$ 上の代数空間とする。
$x \in |X|$ とする。
$x$ の上にある任意の二つの幾何学的点 $\overline{x}$ と
$\overline{x}'$ に対して、茎関手は同型であると主張する。
$|X|$ の定義により、可換図式
$$
\xymatrix{
\overline{x}'' \ar[r] \ar[d] \ar[rd]^{\overline{x}''} &
\overline{x}' \ar[d]^{\overline{x}'} \\
\overline{x} \ar[r]^{\overline{x}} & X.
}
$$
が存在するような第三の幾何学的点 $\overline{x}''$ を見つけられる。
茎関手 $\mathcal{F} \mapsto \mathcal{F}_{\overline{x}}$ は
射 $\overline{x}$ に沿う引き戻しによって与えられ
（他のものについても同様）、引き戻しの関手性から結論を得る。
\end{remark}

\noindent
次の定理は、代数空間の小エタールサイトが十分多くの点をもつことを述べる。

\begin{theorem}
\label{spaces-properties-theorem-exactness-stalks}
$S$ をスキーム、$X$ を $S$ 上の代数空間とする。
集合の層の写像 $a : \mathcal{F} \to \mathcal{G}$ が単射
（それぞれ全射）であるための必要十分条件は、$X$ のすべての幾何学的点に
対して茎上の写像
$a_{\overline{x}} : \mathcal{F}_{\overline{x}} \to \mathcal{G}_{\overline{x}}$
が単射（それぞれ全射）であることである。
$X_\etale$ 上のアーベル層の列が完全であるための必要十分条件は、
$S$ の幾何学的点におけるすべての茎上で完全であることである。
\end{theorem}

\begin{proof}
$X$ がスキームならば定理が正しいことは既知である。
'Etale Cohomology, Theorem
\ref{etale-cohomology-theorem-exactness-stalks} を参照せよ。
$U$ がスキームであるような全射エタール射 $f : U \to X$ を選ぶ。
$\{U \to X\}$ は（$X_{spaces, \etale}$ における）被覆なので、
層の写像が単射または全射かどうかは $U$ への制限によって確認できる。
ここで $\overline{u} : \Spec(k) \to U$ が $U$ の幾何学的点ならば、
$(\mathcal{F}|_U)_{\overline{u}} = \mathcal{F}_{\overline{x}}$
である。ここで $\overline{x} = f \circ \overline{u}$ である。
（これは $\overline{u}$ と $\overline{x}$ における茎を定義する余極限から
明らかであるが、Lemma \ref{spaces-properties-lemma-stalk-pullback} からも従う。）
したがって $U$ に対する結果から $X$ に対する結果が従い、結論を得る。
\end{proof}

\noindent
次の補題は初読時には飛ばすべきである。

\begin{lemma}
\label{spaces-properties-lemma-points-small-etale-site}
$S$ をスキームとする。
$X$ を $S$ 上の代数空間とする。
$p : \Sh(pt) \to \Sh(X_\etale)$ を $X$ の小エタールトポスの点とする。
このとき $X$ の幾何学的点 $\overline{x}$ であって、茎関手
$\mathcal{F} \mapsto \mathcal{F}_p$ が茎関手
$\mathcal{F} \mapsto \mathcal{F}_{\overline{x}}$ と同型になるものが存在する。
\end{lemma}

\begin{proof}
Sites, Lemma \ref{sites-lemma-point-site-topos} により、サイトの点と
付随するトポスの点の間には一対一対応がある。
したがって $p$ はサイト $X_\etale$ の点を定める関手
$u : X_\etale \to \textit{Sets}$ によって与えられると仮定してよい。
構造射 $j : U \to X$ が全射であるような対象
$U \in \Ob(X_\etale)$ を取る。$h_U$ は終層へ全射する層であることに
注意せよ。茎を取ることは完全なので、
$(h_U)_p = u(U)$ は空でない
（Sites, Lemma \ref{sites-lemma-points-recover} を用いよ）。
$x \in u(U)$ を選ぶ。Sites, Lemma
\ref{sites-lemma-point-localize} により、
$p = j_{small} \circ q$ となる点
$q : \Sh(pt) \to \Sh(U_\etale)$ を得る。
したがって関手的に $\mathcal{F}_p = (\mathcal{F}|_U)_q$ である。
'Etale Cohomology, Lemma
\ref{etale-cohomology-lemma-points-small-etale-site} により、
$U$ の幾何学的点 $\overline{u}$ と、
$\mathcal{G} \in \Sh(U_\etale)$ に対する関手的な同型
$\mathcal{G}_q = \mathcal{G}_{\overline{u}}$ が存在する。
$\overline{x} = j \circ \overline{u}$ とおく。
すると Lemma \ref{spaces-properties-lemma-stalk-pullback} により、
$X_\etale$ 上の $\mathcal{F}$ に関して関手的に
$\mathcal{F}_{\overline{x}} \cong (\mathcal{F}|_U)_{\overline{u}}$
であることが分かり、結論を得る。
\end{proof}

\section{アーベル層の台}
\label{spaces-properties-section-support}

\noindent
まず局所切断の台について論じる。

\begin{lemma}
\label{spaces-properties-lemma-support-subsheaf-final}
$S$ をスキーム、$X$ を $S$ 上の代数空間とする。
$\mathcal{F}$ を $X$ のエタールトポスの終対象の部分層とする
（Sites, Example \ref{sites-example-singleton-sheaf} を参照）。
このとき $\mathcal{F} = h_W$ を満たす一意な開部分空間
$W \subset X$ が存在する。
\end{lemma}

\begin{proof}
この条件は、$\Ob(X_{spaces, \etale})$ のすべての
$\varphi : U \to X$ に対して $\mathcal{F}(U)$ が一元集合または
空集合であることを意味する。特に、局所切断は常に貼り合わさる。
$\mathcal{F}(U) \not = \emptyset$ ならば、
$\varphi(U) \subset X$ は開部分空間であり
（Lemma \ref{spaces-properties-lemma-etale-open}）、
$\{\varphi : U \to \varphi(U)\}$ は $X_{spaces, \etale}$ の被覆なので、
$\mathcal{F}(\varphi(U)) \not = \emptyset$ である。
$W = \bigcup_{\varphi : U \to S, \mathcal{F}(U) \not = \emptyset} \varphi(U)$
と取れば結論が従う。
\end{proof}

\begin{lemma}
\label{spaces-properties-lemma-zero-over-image}
$S$ をスキームとする。
$X$ を $S$ 上の代数空間とする。
$\mathcal{F}$ を $X_{spaces, \etale}$ 上のアーベル層とする。
$\sigma \in \mathcal{F}(U)$ を局所切断とする。
次を満たす開部分空間 $W \subset U$ が存在する。
\begin{enumerate}
\item $W \subset U$ は $\sigma|_W = 0$ を満たす $U$ の最大の開部分空間である。
\item $X_{spaces, \etale}$ のすべての $\varphi : V \to U$ に対して
$$
\sigma|_V = 0 \Leftrightarrow \varphi(V) \subset W,
$$
である。
\item $U$ のすべての幾何学的点 $\overline{u}$ に対して
$$
(U, \overline{u}, \sigma) = 0\text{ が次の茎で成り立つ：}\mathcal{F}_{\overline{x}}
\Leftrightarrow
\overline{u} \in W
$$
である。ここで $\overline{x} = (U \to X) \circ \overline{u}$ である。
\end{enumerate}
\end{lemma}

\begin{proof}
$\mathcal{F}$ はエタール位相における層なので、$U_{Zar}$ への
$\mathcal{F}$ の制限は Zariski 位相における $U$ 上の層である。
したがって性質 (1) をもつ Zariski 開部分 $W$ が存在する。
Modules, Lemma \ref{modules-lemma-support-section-closed} を参照せよ。
$\varphi : V \to U$ を $X_{spaces, \etale}$ の矢印とする。
$\varphi(V) \subset U$ は開部分空間であり
（Lemma \ref{spaces-properties-lemma-etale-open}）、
$\{V \to \varphi(V)\}$ はエタール被覆であることに注意する。
したがって $\sigma|_V = 0$ ならば、$\mathcal{F}$ の層条件により
$\sigma|_{\varphi(V)} = 0$ である。これで (2) が従う。
(3) を証明するには、$(U, \overline{u}, \sigma)$ が
$\mathcal{F}_{\overline{x}}$ の零元を定めるならば
$\overline{u} \in W$ であることを示せばよい。
この仮定は、$\sigma|_V = 0$ となるエタール近傍の射
$(V, \overline{v}) \to (U, \overline{u})$ が存在することを意味する。
したがって (2) により $V \to U$ は $W$ の中へ写り、
$\overline{u} \in W$ である。
\end{proof}

\noindent
$S$ をスキームとする。
$X$ を $S$ 上の代数空間とする。
$x \in |X|$ とする。
$\mathcal{F}$ を $X_\etale$ 上の層とする。
Remark \ref{spaces-properties-remark-map-stalks} により、$x$ の上にある幾何学的点における
層 $\mathcal{F}$ の茎の同型類はよく定まる。

\begin{definition}
\label{spaces-properties-definition-support}
$S$ をスキームとする。
$X$ を $S$ 上の代数空間とする。
$\mathcal{F}$ を $X_\etale$ 上のアーベル層とする。
\begin{enumerate}
\item {\it $\mathcal{F}$ の台} とは、$x$ の上にある任意の（ある）
幾何学的点 $\overline{x}$ に対して
$\mathcal{F}_{\overline{x}} \not = 0$ となる点 $x \in |X|$ の集合をいう。
\item $\sigma \in \mathcal{F}(U)$ を切断とする。
{\it $\sigma$ の台} とは閉部分集合 $U \setminus W$ をいう。
ここで $W \subset U$ は、$\sigma$ の制限が零となる $U$ の最大の開部分集合である
（Lemma \ref{spaces-properties-lemma-zero-over-image} を参照）。
\end{enumerate}
\end{definition}

\begin{lemma}
\label{spaces-properties-lemma-support-section-closed}
$S$ をスキームとする。
$X$ を $S$ 上の代数空間とする。
$\mathcal{F}$ を $X_\etale$ 上のアーベル層とする。
$U \in \Ob(X_\etale)$ および $\sigma \in \mathcal{F}(U)$ とする。
\begin{enumerate}
\item $\sigma$ の台は $|X|$ において閉である。
\item $\sigma + \sigma'$ の台は、
$\sigma, \sigma' \in \mathcal{F}(X)$ の台の和集合に含まれる。
\item $\varphi : \mathcal{F} \to \mathcal{G}$ が
$X_\etale$ 上のアーベル層の写像ならば、$\varphi(\sigma)$ の台は
$\sigma \in \mathcal{F}(U)$ の台に含まれる。
\item $\mathcal{F}$ の台は、$\mathcal{F}$ のすべての局所切断の台の像の
和集合である。
\item $\mathcal{F} \to \mathcal{G}$ が全射ならば、
$\mathcal{G}$ の台は $\mathcal{F}$ の台の部分集合である。
\item $\mathcal{F} \to \mathcal{G}$ が単射ならば、
$\mathcal{F}$ の台は $\mathcal{G}$ の台の部分集合である。
\end{enumerate}
\end{lemma}

\begin{proof}
(1) は定義により成り立つ。
(2) と (3) は、$\mathcal{F}$ と $\mathcal{G}$ の $U_{Zar}$ への制限に
対して成り立つので従う。Modules, Lemma
\ref{modules-lemma-support-section-closed} を参照せよ。
(4) は Lemma \ref{spaces-properties-lemma-zero-over-image} の (3) の直接の帰結である。
(5) と (6) は他の部分から従う。
\end{proof}

\begin{lemma}
\label{spaces-properties-lemma-support-sheaf-rings-closed}
代数空間の小エタールサイト上の環の層の台は閉である。
\end{lemma}

\begin{proof}
（本書の規約によれば）環が $0$ であるための必要十分条件は
$1 = 0$ であるから、環の層の台は単位切断の台である。
\end{proof}

\section{代数空間の構造層}
\label{spaces-properties-section-structure-sheaf}

\noindent
代数空間の構造層は、次の補題の環の層である。

\begin{lemma}
\label{spaces-properties-lemma-sheaf-condition-holds}
$S$ をスキーム、$X$ を $S$ 上の代数空間とする。
規則 $U \mapsto \Gamma(U, \mathcal{O}_U)$ は
$X_\etale$ 上の環の層を定める。
\end{lemma}

\begin{proof}
被覆の定義と Descent, Lemma
\ref{descent-lemma-sheaf-condition-holds} から直ちに従う。
\end{proof}

\begin{definition}
\label{spaces-properties-definition-structure-sheaf}
$S$ をスキームとする。
$X$ を $S$ 上の代数空間とする。
$X$ の {\it 構造層} とは、Lemma
\ref{spaces-properties-lemma-sheaf-condition-holds} で記述した小エタールサイト
$X_\etale$ 上の環の層 $\mathcal{O}_X$ をいう。
\end{definition}

\noindent
Lemma \ref{spaces-properties-lemma-characterize-sheaf-small-etale} によれば、層
$\mathcal{O}_X$ は、$X_\etale$ の対象 $U$ を走るエタール層
$(\mathcal{O}_X)_U$ の系に対応する。
その補題の証明と定義から、単に
$(\mathcal{O}_X)_U = \mathcal{O}_U$ であることは明らかである。
ここで $\mathcal{O}_U$ は Descent, Definition
\ref{descent-definition-structure-sheaf} で導入した
$U_\etale$ の構造層である。特に $X$ がスキームならば、
$X$ の小エタールサイト上の層 $\mathcal{O}_X$ が再び得られる。

\medskip\noindent
Lemma \ref{spaces-properties-lemma-compare-etale-sites} の同値
$\Sh(X_\etale) = \Sh(X_{spaces, \etale})$ を介して、
$\mathcal{O}_X$ を $X_{spaces, \etale}$ 上の環の層と考えることもできる。
$Y \to X$ が $X_{spaces, \etale}$ の対象であるときに
$\mathcal{O}_X(Y)$、特に $\mathcal{O}_X(X)$ を計算する方法は
Remark \ref{spaces-properties-remark-explain-equivalence} で説明されている。

\begin{lemma}
\label{spaces-properties-lemma-morphism-ringed-topoi}
$S$ をスキームとする。
$f : X \to Y$ を $S$ 上の代数空間の射とする。
このとき標準的な写像
$f^\sharp : f_{small}^{-1}\mathcal{O}_Y \to \mathcal{O}_X$ が存在し、
$$
(f_{small}, f^\sharp) :
(\Sh(X_\etale), \mathcal{O}_X)
\longrightarrow
(\Sh(Y_\etale), \mathcal{O}_Y)
$$
は環付きトポスの射である。さらに次が成り立つ。
\begin{enumerate}
\item 構成 $f \mapsto (f_{small}, f^\sharp)$ は合成と両立する。
\item $f$ がスキームの射ならば、$f^\sharp$ は Descent, Remark
\ref{descent-remark-change-topologies-ringed} で記述された写像である。
\end{enumerate}
\end{lemma}

\begin{proof}
Lemma \ref{spaces-properties-lemma-f-map} により、$\mathcal{O}_Y$ から
$\mathcal{O}_X$ への $f$-写像を与えれば十分である。
言い換えると、$U \in X_\etale$, $V \in Y_\etale$ である
各可換図式
$$
\xymatrix{
U \ar[d]_g \ar[r] & X \ar[d]^f \\
V \ar[r] & Y
}
$$
に対して、環の写像
$
(f^\sharp)_{(U, V, g)} :
\Gamma(V, \mathcal{O}_V)
\to
\Gamma(U, \mathcal{O}_U).
$
を与える必要がある。もちろん
$(f^\sharp)_{(U, V, g)} = g^\sharp$ と取る。
これが制限写像と両立し、したがって実際に $f$-写像を与えることは明らかである。
合成との両立性、および Descent, Remark
\ref{descent-remark-change-topologies-ringed} の構成との一致の確認は省略する。
\end{proof}

\begin{lemma}
\label{spaces-properties-lemma-reduced-space}
$S$ をスキーム、$X$ を $S$ 上の代数空間とする。
次は同値である。
\begin{enumerate}
\item $X$ は被約である。
\item すべての $x \in |X|$ に対して $x$ における $X$ の局所環は被約である
（Remark \ref{spaces-properties-remark-list-properties-local-ring-local-etale-topology}）。
\end{enumerate}
この場合、$\Gamma(X, \mathcal{O}_X)$ は被約環であり、
$f \in \Gamma(X, \mathcal{O}_X)$ が $X = V(f)$ を満たすならば $f = 0$ である。
\end{lemma}

\begin{proof}
(1) と (2) の同値性は、$X$ 上エタールなアフィンスキームに
Properties, Lemma \ref{properties-lemma-characterize-reduced} を
適用すれば従う。最後の主張は、引用した補題と、
$\Gamma(X, \mathcal{O}_X)$ が $X$ 上エタールなある被約スキーム $U$ の
$\Gamma(U, \mathcal{O}_U)$ の部分環であるという事実から従う。
\end{proof}

\section{構造層の茎}
\label{spaces-properties-section-stalks-structure-sheaf}

\noindent
この節は 'Etale Cohomology, Section
\ref{etale-cohomology-section-stalks-structure-sheaf} の類似物である。

\begin{lemma}
\label{spaces-properties-lemma-describe-etale-local-ring}
$S$ をスキームとする。
$X$ を $S$ 上の代数空間とする。
$\overline{x}$ を $X$ の幾何学的点とする。
$(U, \overline{u})$ を、$U$ がスキームであるような
$\overline{x}$ のエタール近傍とする。このとき
$$
\mathcal{O}_{X, \overline{x}} =
\mathcal{O}_{U, \overline{u}} =
\mathcal{O}_{U, u}^{sh}
$$
である。左辺は $X$ の構造層の茎であり、右辺は
$\overline{u}$ の中心である点 $u$ における $U$ の局所環の厳密 Hensel 化である。
\end{lemma}

\begin{proof}
$U_\etale$ 上の構造層 $\mathcal{O}_U$ は $X$ の構造層の制限である。
したがって第一の等式は Lemma \ref{spaces-properties-lemma-stalk-pullback} の (4) から従う。
第二の等式は 'Etale Cohomology, Lemma
\ref{etale-cohomology-lemma-describe-etale-local-ring} で説明されている。
\end{proof}

\begin{definition}
\label{spaces-properties-definition-etale-local-rings}
$S$ をスキームとする。
$X$ を $S$ 上の代数空間とする。
$\overline{x}$ を点 $x \in |X|$ の上にある $X$ の幾何学的点とする。
\begin{enumerate}
\item {\it $\overline{x}$ における $X$ のエタール局所環} とは、
$X_\etale$ 上の構造層 $\mathcal{O}_X$ の $\overline{x}$ における茎をいう。
記法：$\mathcal{O}_{X, \overline{x}}$。
\item {\it $\overline{x}$ における $X$ の厳密 Hensel 化} とは
スキーム $\Spec(\mathcal{O}_{X, \overline{x}})$ をいう。
\end{enumerate}
\end{definition}

\noindent
$\overline{x}$ における $X$ の厳密 Hensel 化の（$X$ 上のスキームとしての）
同型型は、点 $x \in |X|$ のみに依存し、$x$ の上にある幾何学的点の
選択には依存しない。Remark \ref{spaces-properties-remark-map-stalks} を参照せよ。

\begin{lemma}
\label{spaces-properties-lemma-etale-site-locally-ringed}
$S$ をスキームとする。
$X$ を $S$ 上の代数空間とする。
構造層 $\mathcal{O}_X$ を備えた小エタールサイト $X_\etale$ は
局所環付きサイトである。Modules on Sites, Definition
\ref{sites-modules-definition-locally-ringed} を参照せよ。
\end{lemma}

\begin{proof}
茎 $\mathcal{O}_{X, \overline{x}}$ が局所的であり、
$S_\etale$ が十分多くの点をもつことから従う。
Lemma \ref{spaces-properties-lemma-describe-etale-local-ring} および
Theorem \ref{spaces-properties-theorem-exactness-stalks} を参照せよ。
これにより小エタールサイトが局所環付きになることについては、
Modules on Sites, Lemmas
\ref{sites-modules-lemma-locally-ringed-stalk} および
\ref{sites-modules-lemma-ringed-stalk-not-zero} を参照せよ。
\end{proof}

\begin{lemma}
\label{spaces-properties-lemma-dimension-local-ring}
$S$ をスキーム、$X$ を $S$ 上の代数空間とする。
$x \in |X|$ を点とし、$d \in \{0, 1, 2, \ldots, \infty\}$ とする。
次は同値である。
\begin{enumerate}
\item $x$ における $X$ の局所環の次元
（Definition \ref{spaces-properties-definition-dimension-local-ring}）は $d$ である。
\item $x$ の上にあるある幾何学的点 $\overline{x}$ に対して
$\dim(\mathcal{O}_{X, \overline{x}}) = d$ である。
\item $x$ の上にある任意の幾何学的点 $\overline{x}$ に対して
$\dim(\mathcal{O}_{X, \overline{x}}) = d$ である。
\end{enumerate}
\end{lemma}

\begin{proof}
(2) と (3) の同値性は、$\mathcal{O}_{X, \overline{x}}$ の同型型が
$x \in |X|$ のみに依存することから従う。
Remark \ref{spaces-properties-remark-map-stalks} を参照せよ。
Lemma \ref{spaces-properties-lemma-describe-etale-local-ring} を用いると、
(1) と (2)$+$(3) の同値性は、任意の局所環 $R$ に対して
$\dim(R) = \dim(R^{sh})$ であるという主張に帰着する。
これは More on Algebra, Lemma
\ref{more-algebra-lemma-henselization-dimension} である。
\end{proof}

\begin{lemma}
\label{spaces-properties-lemma-dimension-decent-invariant-under-etale}
$S$ をスキームとする。$f : X \to Y$ を $S$ 上の代数空間のエタール射とし、
$x \in X$ とする。このとき (1) $\dim_x(X) = \dim_{f(x)}(Y)$ であり、
(2) $x$ における $X$ の局所環の次元は $f(x)$ における $Y$ の局所環の
次元に等しい。$f$ が全射ならば、(3) $\dim(X) = \dim(Y)$ である。
\end{lemma}

\begin{proof}
スキーム $U$、点 $u \in U$、および $u$ を $x$ へ写すエタール射
$U \to X$ を選ぶ。このとき合成 $U \to Y$ もエタールであり、
$u$ を $f(x)$ へ写す。したがって、該当する整数は点 $u$ における
スキーム $U$ の挙動によって定義されるので、(1) と (2) が従う。
(1) については Definition \ref{spaces-properties-definition-dimension-at-point} を参照せよ。
(3) は (1) の直接の帰結である。
Definition \ref{spaces-properties-definition-dimension} を参照せよ。
\end{proof}

\begin{lemma}
\label{spaces-properties-lemma-reduced-local-ring}
$S$ をスキーム、$X$ を $S$ 上の代数空間とする。
$x \in |X|$ を点とする。次は同値である。
\begin{enumerate}
\item $x$ における $X$ の局所環は被約である
（Remark \ref{spaces-properties-remark-list-properties-local-ring-local-etale-topology}）。
\item $x$ の上にあるある幾何学的点 $\overline{x}$ に対して
$\mathcal{O}_{X, \overline{x}}$ は被約である。
\item $x$ の上にある任意の幾何学的点 $\overline{x}$ に対して
$\mathcal{O}_{X, \overline{x}}$ は被約である。
\end{enumerate}
\end{lemma}

\begin{proof}
(2) と (3) の同値性は、$\mathcal{O}_{X, \overline{x}}$ の同型型が
$x \in |X|$ のみに依存することから従う。
Remark \ref{spaces-properties-remark-map-stalks} を参照せよ。
Lemma \ref{spaces-properties-lemma-describe-etale-local-ring} を用いると、
(1) と (2)$+$(3) の同値性は、局所環が被約であるための必要十分条件は
その厳密 Hensel 化が被約であることだという主張に帰着する。
これは More on Algebra, Lemma
\ref{more-algebra-lemma-henselization-reduced} である。
\end{proof}

\section{局所的既約性}
\label{spaces-properties-section-irreducible-local-ring}

\noindent
代数空間の点には、よく定まったエタール局所環がある。
スキームの場合、これは局所環の厳密 Hensel 化に対応する。
一般には、与えられた点を通るスキームまたは代数空間の既約成分の個数を
エタール局所環から知ることはできず、幾何学的枝の個数だけを数えられる。

\begin{lemma}
\label{spaces-properties-lemma-irreducible-local-ring}
$S$ をスキームとする。
$X$ を $S$ 上の代数空間とする。
$x \in |X|$ を点とする。次は同値である。
\begin{enumerate}
\item 任意のスキーム $U$、エタール射 $a : U \to X$、および
$a(u) = x$ を満たす $u \in U$ に対して、局所環
$\mathcal{O}_{U, u}$ は一意な極小素イデアルをもつ。
\item 任意のスキーム $U$、エタール射 $a : U \to X$、および
$a(u) = x$ を満たす $u \in U$ に対して、$u$ を通る $U$ の
既約成分が一意に存在する。
\item 任意のスキーム $U$、エタール射 $a : U \to X$、および
$a(u) = x$ を満たす $u \in U$ に対して、局所環
$\mathcal{O}_{U, u}$ は単枝である。
\item 任意のスキーム $U$、エタール射 $a : U \to X$、および
$a(u) = x$ を満たす $u \in U$ に対して、局所環
$\mathcal{O}_{U, u}$ は幾何学的に単枝である。
\item $x$ の上にある任意の幾何学的点 $\overline{x}$ に対して、
$\mathcal{O}_{X, \overline{x}}$ は一意な極小素イデアルをもつ。
\end{enumerate}
\end{lemma}

\begin{proof}
(1) と (2) の同値性は、$u$ を通る $U$ の既約成分と
$u$ における $U$ の局所環の極小素イデアルとの間に $1$-$1$ 対応が
あることから従う。$a : U \to X$ と $u \in U$ は (1) のとおりとする。
Lemma \ref{spaces-properties-lemma-describe-etale-local-ring} により、
$\mathcal{O}_{X, \overline{x}}$ は $\mathcal{O}_{U, u}$ の厳密 Hensel 化である。
特に、More on Algebra, Lemma
\ref{more-algebra-lemma-geometrically-unibranch} により (4) と (5) は同値である。
(2), (3), (4) の同値性は More on Morphisms, Lemma
\ref{more-morphisms-lemma-nr-branches} から従う。
\end{proof}

\begin{definition}
\label{spaces-properties-definition-unibranch}
$S$ をスキーム、$X$ を $S$ 上の代数空間とする。
$x \in |X|$ とする。Lemma \ref{spaces-properties-lemma-irreducible-local-ring} の
同値な条件が成り立つとき、$X$ は {\it $x$ において幾何学的に単枝}
であるという。すべての $x \in |X|$ において $X$ が
幾何学的に単枝であるとき、$X$ は {\it 幾何学的に単枝} であるという。
\end{definition}

\noindent
これはスキームに対する定義
（Properties, Definition \ref{properties-definition-unibranch}）と整合する。

\begin{lemma}
\label{spaces-properties-lemma-nr-branches-local-ring}
$S$ をスキーム、$X$ を $S$ 上の代数空間とする。
$x \in |X|$ を点とし、$n \in \{1, 2, \ldots\}$ を整数とする。
次は同値である。
\begin{enumerate}
\item 任意のスキーム $U$、エタール射 $a : U \to X$、および
$a(u) = x$ を満たす $u \in U$ に対して、局所環 $\mathcal{O}_{U, u}$ の
極小素イデアルの個数は $\leq n$ であり、少なくとも一つの
$U, a, u$ の選択に対してその個数は $n$ である。
\item 任意のスキーム $U$、エタール射 $a : U \to X$、および
$a(u) = x$ を満たす $u \in U$ に対して、$u$ を通る $U$ の既約成分の
個数は $\leq n$ であり、少なくとも一つの $U, a, u$ の選択に対して
その個数は $n$ である。
\item 任意のスキーム $U$、エタール射 $a : U \to X$、および
$a(u) = x$ を満たす $u \in U$ に対して、$u$ における $U$ の枝の個数は
$\leq n$ であり、少なくとも一つの $U, a, u$ の選択に対して
その個数は $n$ である。
\item 任意のスキーム $U$、エタール射 $a : U \to X$、および
$a(u) = x$ を満たす $u \in U$ に対して、$u$ における $U$ の
幾何学的枝の個数は $n$ である。
\item $\mathcal{O}_{X, \overline{x}}$ の極小素イデアルの個数は $n$ である。
\end{enumerate}
\end{lemma}

\begin{proof}
(1) と (2) の同値性は、$u$ を通る $U$ の既約成分と
$u$ における $U$ の局所環の極小素イデアルとの間に $1$-$1$ 対応が
あることから従う。$a : U \to X$ と $u \in U$ は (1) のとおりとする。
Lemma \ref{spaces-properties-lemma-describe-etale-local-ring} により、
$\mathcal{O}_{X, \overline{x}}$ は $\mathcal{O}_{U, u}$ の厳密 Hensel 化である。
$u$ における $U$ の（幾何学的）枝の個数とは、
$\mathcal{O}_{U, u}$ の（厳密）Hensel 化の極小素イデアルの個数であることを
思い出す。特に (4) と (5) は同値である。
(2), (3), (4) の同値性は More on Morphisms, Lemma
\ref{more-morphisms-lemma-nr-branches} から従う。
\end{proof}

\begin{definition}
\label{spaces-properties-definition-number-of-geometric-branches}
$S$ をスキームとする。
$X$ を $S$ 上の代数空間とする。
$x \in |X|$ とする。{\it $x$ における $X$ の幾何学的枝の個数} とは、
Lemma \ref{spaces-properties-lemma-nr-branches-local-ring} の同値な条件が成り立つ場合には
$n \in \mathbf{N}$、そうでなければ $\infty$ をいう。
\end{definition}

\section{Noether 代数空間}
\label{spaces-properties-section-noetherian}

\noindent
Section \ref{spaces-properties-section-types-properties} ですでに局所 Noether 代数空間を定義した。

\begin{definition}
\label{spaces-properties-definition-noetherian}
$S$ をスキーム、$X$ を $S$ 上の代数空間とする。
$X$ が準コンパクト、準分離、かつ局所 Noether であるとき、
$X$ は {\it Noether} であるという。
\end{definition}

\noindent
Noether 代数空間 $X$ は単に準コンパクトかつ局所 Noether であるだけでなく、
準分離でもあることに注意する。局所 Noether スキームは準分離であるため、
これは Noether スキームの定義と矛盾しない。
Properties, Lemma
\ref{properties-lemma-locally-Noetherian-quasi-separated} を参照せよ。
これは代数空間には当てはまらない。実際、
$X = \mathbf{A}^1_k/\mathbf{Z}$（Spaces, Example
\ref{spaces-example-affine-line-translation} を参照）は
局所 Noether かつ準コンパクトであるが準分離ではない
（したがって本書の定義では Noether ではない）。

\medskip\noindent
上の選択の帰結として、Noether 代数空間上有限型の代数空間は
自動的には Noether ではない。すなわち、Morphisms, Lemma
\ref{morphisms-lemma-finite-type-noetherian} の類似は成り立たない。
正しい主張は、Noether 代数空間上有限表示の代数空間は Noether である
というものである（Morphisms of Spaces, Lemma
\ref{spaces-morphisms-lemma-finite-presentation-noetherian} を参照）。

\medskip\noindent
Noether 代数空間 $X$ はスキームに非常に近い。
この節の残りでは、それを示すいくつかの補題を集める。

\begin{lemma}
\label{spaces-properties-lemma-Noetherian-topology}
$S$ をスキーム、$X$ を $S$ 上の代数空間とする。
\begin{enumerate}
\item $X$ が局所 Noether ならば、$|X|$ は局所 Noether 位相空間である。
\item $X$ が準コンパクトかつ局所 Noether ならば、
$|X|$ は Noether 位相空間である。
\end{enumerate}
\end{lemma}

\begin{proof}
$X$ が局所 Noether であると仮定する。
スキーム $U$ と全射エタール射 $U \to X$ を選ぶ。
$X$ が局所 Noether なので $U$ も局所 Noether である。
Properties, Lemma \ref{properties-lemma-Noetherian-topology} により、
$|U|$ は局所 Noether 位相空間である。
$|U| \to |X|$ は開かつ全射なので、Topology, Lemma
\ref{topology-lemma-image-Noetherian} により $|X|$ は局所 Noether である。
これで (1) が証明された。$X$ が準コンパクトかつ局所 Noether ならば、
$|X|$ は準コンパクトかつ局所 Noether である。
したがって Topology, Lemma
\ref{topology-lemma-quasi-compact-locally-Noetherian-Noetherian} により
$|X|$ は Noether である。
\end{proof}

\begin{lemma}
\label{spaces-properties-lemma-Noetherian-sober}
$S$ をスキーム、$X$ を $S$ 上の代数空間とする。
$X$ が Noether ならば、$|X|$ は sober Noether 位相空間である。
\end{lemma}

\begin{proof}
準分離代数空間の台位相空間は sober である。
Lemma \ref{spaces-properties-lemma-quasi-separated-sober} を参照せよ。
Lemma \ref{spaces-properties-lemma-Noetherian-topology} により、それは Noether である。
\end{proof}

\begin{lemma}
\label{spaces-properties-lemma-Noetherian-local-ring-Noetherian}
$S$ をスキームとし、$X$ を $S$ 上の Noether 代数空間とする。
$\overline{x}$ を $X$ の幾何学的点とする。このとき
$\mathcal{O}_{X, \overline{x}}$ は Noether 局所環である。
\end{lemma}

\begin{proof}
$U$ がスキームであるような $\overline{x}$ のエタール近傍
$(U, \overline{u})$ を選ぶ。Lemma \ref{spaces-properties-lemma-describe-etale-local-ring} により、
$\mathcal{O}_{X, \overline{x}}$ は $u$ における $U$ の局所環の
厳密 Hensel 化である。Noether 空間の定義により、スキーム $U$ は
局所 Noether である。したがって More on Algebra, Lemma
\ref{more-algebra-lemma-henselization-noetherian} により結論を得る。
\end{proof}

\section{正則代数空間}
\label{spaces-properties-section-regular}

\noindent
Section \ref{spaces-properties-section-types-properties} ですでに正則代数空間を定義した。

\begin{lemma}
\label{spaces-properties-lemma-regular}
$S$ をスキームとする。
$X$ を $S$ 上の局所 Noether 代数空間とする。
次は同値である。
\begin{enumerate}
\item $X$ は正則である。
\item すべてのエタール局所環 $\mathcal{O}_{X, \overline{x}}$ は正則である。
\end{enumerate}
\end{lemma}

\begin{proof}
$U$ をスキームとし、$U \to X$ を全射エタール射とする。
仮定により $U$ は局所 Noether である。さらに、すべてのエタール局所環
$\mathcal{O}_{X, \overline{x}}$ は $U$ 上のある局所環の厳密 Hensel 化であり、
逆も成り立つ。Lemma \ref{spaces-properties-lemma-describe-etale-local-ring} を参照せよ。
したがって More on Algebra, Lemma
\ref{more-algebra-lemma-henselization-regular} により、(2) は
$U$ のすべての局所環が正則であること、すなわち $U$ が正則スキームであることと
同値である（Properties, Lemma
\ref{properties-lemma-characterize-regular} を参照）。
Definition \ref{spaces-properties-definition-type-property} により、これは (1) と同値である。
\end{proof}

\noindent
Descent, Lemma \ref{descent-lemma-regular-local-ring-local} を用いて、
代数空間 $X$ が点 $x$ において正則であることの意味を定義できる。

\begin{definition}
\label{spaces-properties-definition-regular-at-point}
$S$ をスキーム、$X$ を $S$ 上の代数空間とする。
$x \in |X|$ を点とする。スキームから $X$ へのエタール射
$a : U \to X$ と、$a(u) = x$ を満たす点 $u \in U$ からなる
任意の（同値なことに、ある）対 $(a : U \to X, u)$ に対して
$\mathcal{O}_{U, u}$ が正則局所環であるとき、
{\it $X$ は $x$ において正則} であるという。
\end{definition}

\noindent
Definition \ref{spaces-properties-definition-property-at-point},
Lemma \ref{spaces-properties-lemma-local-source-target-at-point}, および
Descent, Lemma \ref{descent-lemma-regular-local-ring-local} を参照せよ。

\begin{lemma}
\label{spaces-properties-lemma-regular-at-x}
$S$ をスキーム、$X$ を $S$ 上の代数空間とする。
$x \in |X|$ を点とする。次は同値である。
\begin{enumerate}
\item $X$ は $x$ において正則である。
\item $x$ の上にある任意の（同値なことに、ある）幾何学的点
$\overline{x}$ に対して、エタール局所環
$\mathcal{O}_{X, \overline{x}}$ は正則である。
\end{enumerate}
\end{lemma}

\begin{proof}
$U$ をスキーム、$u \in U$ を点とし、$a : U \to X$ を $u$ を $x$ へ写す
エタール射とする。$x$ の上にある $X$ の任意の幾何学的点
$\overline{x}$ に対して、エタール局所環
$\mathcal{O}_{X, \overline{x}}$ は $u$ における $U$ の局所環の
厳密 Hensel 化である。Lemma \ref{spaces-properties-lemma-describe-etale-local-ring} を参照せよ。
したがって More on Algebra, Lemma
\ref{more-algebra-lemma-henselization-regular} により結論を得る。
\end{proof}

\begin{lemma}
\label{spaces-properties-lemma-regular-normal}
正則代数空間は正規である。
\end{lemma}

\begin{proof}
これは定義とスキームの場合から従う。
Properties, Lemma \ref{properties-lemma-regular-normal} を参照せよ。
\end{proof}

\section{代数空間上の加群の層}
\label{spaces-properties-section-modules}

\noindent
$X$ が代数空間ならば、$X$ 上の加群の層とは、$X$ の小エタールサイト上の
$\mathcal{O}_X$-加群の層をいう。ここで $\mathcal{O}_X$ は $X$ の構造層である。
加群の層の圏を $\textit{Mod}(\mathcal{O}_X)$ と書く。

\medskip\noindent
代数空間の射 $f : X \to Y$ が与えられると、
Lemma \ref{spaces-properties-lemma-morphism-ringed-topoi} により環付きトポスの射を得る。
したがって Modules on Sites, Definition
\ref{sites-modules-definition-pushforward} により、よく定まった引き戻し関手と
直像関手
\begin{equation}
\label{spaces-properties-equation-push-pull}
f^* :
\textit{Mod}(\mathcal{O}_Y)
\longrightarrow
\textit{Mod}(\mathcal{O}_X), \quad
f_* :
\textit{Mod}(\mathcal{O}_X)
\longrightarrow
\textit{Mod}(\mathcal{O}_Y)
\end{equation}
を得る。これらは通常の仕方で随伴する。
$g : Y \to Z$ が $S$ 上の代数空間の別の射ならば、
環付きトポスの射が対応する仕方で合成されるので（補題による）、単に
$(g \circ f)^* = f^* \circ g^*$ および
$(g \circ f)_* = g_* \circ f_*$ である。

\begin{lemma}
\label{spaces-properties-lemma-etale-exact-pullback}
$S$ をスキームとする。
$f : X \to Y$ を $S$ 上の代数空間のエタール射とする。
このとき $f^{-1}\mathcal{O}_Y = \mathcal{O}_X$ であり、
$\mathcal{O}_Y$-加群の任意の層 $\mathcal{G}$ に対して
$f^*\mathcal{G} = f_{small}^{-1}\mathcal{G}$ である。
特に $f^* : \textit{Mod}(\mathcal{O}_Y) \to
\textit{Mod}(\mathcal{O}_X)$ は完全である。
\end{lemma}

\begin{proof}
Lemma \ref{spaces-properties-lemma-etale-morphism-topoi} における逆像の記述と構造層の定義により、
$f_{small}^{-1}\mathcal{O}_Y = \mathcal{O}_X$ であることは明らかである。
定義により引き戻しは
$$
f^*\mathcal{G} =
f_{small}^{-1}\mathcal{G} \otimes_{f_{small}^{-1}\mathcal{O}_Y}
\mathcal{O}_X
$$
なので、$f^*\mathcal{G} = f_{small}^{-1}\mathcal{G}$ と結論する。
$f_{small}$ はトポスの射であるから $f_{small}^{-1}$ は完全であり、
完全性は明らかである。
\end{proof}

\noindent
Equation (\ref{spaces-properties-equation-restrict}) で導入した記号の濫用を続け、
上の補題の状況では
\begin{equation}
\label{spaces-properties-equation-restrict-modules}
\mathcal{G}|_{X_\etale}
= f^*\mathcal{G}
= f_{small}^{-1}\mathcal{G}
\end{equation}
と書く。これについては Section \ref{spaces-properties-section-localize} でより技術的に論じる。

\begin{lemma}
\label{spaces-properties-lemma-pushforward-etale-base-change-modules}
$S$ をスキームとする。$S$ 上の代数空間のカルテジアン図式
$$
\xymatrix{
X' \ar[r] \ar[d]_{f'} & X \ar[d]^f \\
Y' \ar[r]^g & Y
}
$$
を考える。$\mathcal{F} \in \textit{Mod}(\mathcal{O}_X)$ とする。
$g$ がエタールならば、
$f'_*(\mathcal{F}|_{X'}) = (f_*\mathcal{F})|_{Y'}$\footnote{図式の可換性と
(\ref{spaces-properties-equation-restrict-modules}) により
$(f')^*(\mathcal{G}|_{Y'}) = (f^*\mathcal{G})|_{X'}$ でもある。}であり、
$\textit{Mod}(\mathcal{O}_{Y'})$ において
$R^if'_*(\mathcal{F}|_{X'}) = (R^if_*\mathcal{F})|_{Y'}$ である。
\end{lemma}

\begin{proof}
これは加群の場合における
Lemma \ref{spaces-properties-lemma-pushforward-etale-base-change} の言い換えである。
\end{proof}

\begin{lemma}
\label{spaces-properties-lemma-characterize-module-small-etale}
$S$ をスキーム、$X$ を $S$ 上の代数空間とする。
$\mathcal{O}_X$-加群の層 $\mathcal{F}$ は次のデータによって与えられる。
\begin{enumerate}
\item 各 $U \in \Ob(X_\etale)$ に対する、$U_\etale$ 上の
$\mathcal{O}_U$-加群の層 $\mathcal{F}_U$。
\item $X_\etale$ の各 $f : U' \to U$ に対する同型
$c_f : f_{small}^*\mathcal{F}_U \to \mathcal{F}_{U'}$。
\end{enumerate}
これらのデータには、$X_\etale$ の任意の $f : U' \to U$ と
$g : U'' \to U'$ に対して、合成 $c_g \circ g_{small}^*c_f$ が
$c_{f \circ g}$ に等しいという条件を課す。
\end{lemma}

\begin{proof}
Lemmas \ref{spaces-properties-lemma-etale-exact-pullback} および
\ref{spaces-properties-lemma-characterize-sheaf-small-etale} を組み合わせ、
$X_\etale$ の対象間の任意の射がスキームのエタール射であることを用いる。
\end{proof}

\section{エタール局所化}
\label{spaces-properties-section-localize}

\noindent
この節を読むことは何としても避けるべきである。

\medskip\noindent
$X \to Y$ を代数空間のエタール射とする。
このとき $X$ は $Y_{spaces, \etale}$ の対象であり、定義から直ちに
（Lemma \ref{spaces-properties-lemma-etale-morphism-topoi} の証明も参照）
\begin{equation}
\label{spaces-properties-equation-localize}
X_{spaces, \etale} = Y_{spaces, \etale}/X
\end{equation}
である。右辺は対象 $X$ におけるサイト $Y_{spaces, \etale}$ の局所化である。
Sites, Definition \ref{sites-definition-localize} を参照せよ。
さらに Lemma \ref{spaces-properties-lemma-etale-exact-pullback} により、
この同一視は構造層と両立する。
したがって環付きサイト $(X_{spaces, \etale}, \mathcal{O}_X)$ は、
対象 $X$ における環付きサイト $(Y_{spaces, \etale}, \mathcal{O}_Y)$ の
局所化と同一視される：
\begin{equation}
\label{spaces-properties-equation-localize-ringed}
(X_{spaces, \etale}, \mathcal{O}_X) =
(Y_{spaces, \etale}/X, \mathcal{O}_Y|_{Y_{spaces, \etale}/X})
\end{equation}
右辺で用いた環付きサイトの局所化は Modules on Sites, Definition
\ref{sites-modules-definition-localize-ringed-site} で定義されている。

\medskip\noindent
ここで $X \to Y$ を代数空間のエタール射とし、$X$ がスキームであると仮定する。
すると $X$ は $Y_\etale$ の対象であり、
\begin{equation}
\label{spaces-properties-equation-localize-at-scheme}
X_\etale = Y_\etale/X
\end{equation}
および上と同様に
\begin{equation}
\label{spaces-properties-equation-localize-at-scheme-ringed}
(X_\etale, \mathcal{O}_X) =
(Y_\etale/X, \mathcal{O}_Y|_{Y_\etale/X})
\end{equation}
が従う。

\medskip\noindent
最後に、$X \to Y$ が代数空間のエタール射で、$X$ がアフィンスキームならば、
$X$ は $Y_{affine, \etale}$ の対象であり、
\begin{equation}
\label{spaces-properties-equation-localize-at-affine}
X_{affine, \etale} = Y_{affine, \etale}/X
\end{equation}
および上と同様に
\begin{equation}
\label{spaces-properties-equation-localize-at-affine-ringed}
(X_{affine, \etale}, \mathcal{O}_X) =
(Y_{affine, \etale}/X, \mathcal{O}_Y|_{Y_{affine, \etale}/X})
\end{equation}
である。

\medskip\noindent
次に、これらの局所化が射と両立することを示す。

\begin{lemma}
\label{spaces-properties-lemma-relocalize-morphism}
$S$ をスキームとする。$p$ と $q$ がエタールであるような
$S$ 上の代数空間の可換図式
$$
\xymatrix{
U \ar[d]_p \ar[r]_g & V \ar[d]^q \\
X \ar[r]^f & Y
}
$$
を考える。$U \to X$ と $V \to Y$ に対する同一視
(\ref{spaces-properties-equation-localize-ringed}) を介して、環付きトポスの射
$$
(g_{spaces, \etale}, g^\sharp) :
(\Sh(U_{spaces, \etale}), \mathcal{O}_U)
\longrightarrow
(\Sh(V_{spaces, \etale}), \mathcal{O}_V)
$$
は、環付きサイトの射 $(f_{spaces, \etale}, f^\sharp)$ と
$g$ に対応する写像 $c : U \to V \times_Y X$ から Modules on Sites,
Lemma \ref{sites-modules-lemma-relocalize-morphism-ringed-sites} で構成される射
$(f_{spaces, \etale, c}, f_c^\sharp)$ と $2$-同型である。
\end{lemma}

\begin{proof}
射 $(f_{spaces, \etale, c}, f_c^\sharp)$ は、局所化と基底変換写像の合成
$f' \circ j$ として定義される。同様に $g$ は合成
$U \to V \times_Y X \to V$ である。
したがって次の二つの場合に補題を証明すれば十分である：
(1) $f = \text{id}$、(2) $U = X \times_Y V$。
(1) の場合、Lemma \ref{spaces-properties-lemma-etale-permanence} により射
$g : U \to V$ はエタールである。
したがって Equations (\ref{spaces-properties-equation-localize}) および
(\ref{spaces-properties-equation-localize-ringed}) の周辺の議論により、
$(g_{spaces, \etale}, g^\sharp)$ は局所化射である。
これはこの場合の補題の内容そのものである。
(2) の場合、射 $g_{spaces, \etale}$ は環付きサイトの射を与える関手
$V_{spaces, \etale} \to U_{spaces, \etale}$,
$V'/V \mapsto V' \times_V U/U$ から得られる。
これは射 $f'$ を定義する関手でもある。
Sites, Lemma \ref{sites-lemma-localize-morphism} を参照せよ。
この場合に $(f')^\sharp = g^\sharp$ であることの確認は省略する
（いずれも $f^\sharp$ の $U_{spaces, \etale}$ への制限である）。
\end{proof}

\begin{lemma}
\label{spaces-properties-lemma-relocalize-morphism-at-schemes}
Lemma \ref{spaces-properties-lemma-relocalize-morphism} と同じ記法と仮定を用い、
さらに $U$ と $V$ がスキームであると仮定する。
$U \to X$ と $V \to Y$ に対する同一視
(\ref{spaces-properties-equation-localize-at-scheme-ringed}) を介して、環付きトポスの射
$$
(g_{small}, g^\sharp) :
(\Sh(U_\etale), \mathcal{O}_U)
\longrightarrow
(\Sh(V_\etale), \mathcal{O}_V)
$$
は、$(f_{small}, f^\sharp)$ と $g$ に対応する写像
$s : h_U \to f_{small}^{-1}h_V$ から Modules on Sites, Lemma
\ref{sites-modules-lemma-relocalize-morphism-ringed-topoi} で構成される射
$(f_{small, s}, f_s^\sharp)$ と $2$-同型である。
\end{lemma}

\begin{proof}
環付きトポスの射 $(g_{small}, g^\sharp)$ は、環付きサイトの射
$(g_{spaces, \etale}, g^\sharp)$ に付随する環付きトポスの射と
$2$-同型であることに注意する。
したがって Lemma \ref{spaces-properties-lemma-relocalize-morphism} および
Modules on Sites, Lemma
\ref{sites-modules-lemma-relocalize-morphism-compare} により結論を得る。
\end{proof}

\noindent
最後に、代数空間 $Y$ の小エタールサイト $Y_\etale$ 上の集合の層と、
$Y$ 上エタールな代数空間との関係を論じる。
$S$ をスキーム、$Y$ を $S$ 上の代数空間とする。
$\mathcal{F}$ を $\Sh(Y_\etale)$ の対象とする。規則
$$
X : (\Sch/S)_{fppf}^{opp} \longrightarrow \textit{Sets}
$$
および
$$
X(T) = \{(y, s) \mid y : T \to Y\text{ は }S\text{ 上の射であり、}
s \in \Gamma(T, y_{small}^{-1}\mathcal{F})\}
$$
によって定義される関手を考える。射 $g : T' \to T$ が与えられると、
制限写像は $(y, s)$ を $(y \circ g, g_{small}^{-1}s)$ へ写す。
Lemma \ref{spaces-properties-lemma-functoriality-etale-site} により
$y_{small} \circ g_{small} = (y \circ g)_{small}$ なので、これは意味をもつ。
対 $(y, s)$ を $y$ へ写す標準的な写像 $X \to Y$ がある。

\begin{lemma}
\label{spaces-properties-lemma-sheaf-gives-space}
$S$ をスキーム、$Y$ を $S$ 上の代数空間とする。
$\mathcal{F}$ を $Y_\etale$ 上の集合の層とする。
集合論的な条件が満たされるならば（証明を参照）、次が成り立つ。
\begin{enumerate}
\item 上で $\mathcal{F}$ に付随させた関手 $X$ は代数空間である。
\item 写像 $X \to Y$ は代数空間のエタール射である。
\item 同一視 $\Sh(Y_\etale) = \Sh(Y_{spaces, \etale})$ を介して
$\mathcal{F} \cong h_X$ である。
\item $\mathcal{F} \cong f_{small, !}*$ である。ここで $*$ は圏
$\Sh(X_\etale)$ の終対象であり、Lemma
\ref{spaces-properties-lemma-etale-morphism-topoi} により $f_{small, !}$ が存在する。
\end{enumerate}
\end{lemma}

\begin{proof}
$X$ が fppf 位相に対する層であることを証明しよう。
実際、$\{g_i : T_i \to T\}$ を $(\Sch/S)_{fppf}$ の被覆とし、
$(y_i, s_i) \in X(T_i)$ は貼り合わせ条件、すなわち
$(y_i, s_i)$ と $(y_j, s_j)$ の $T_i \times_T T_j$ への制限が
一致するという条件を満たすとする。
$Y$ は fppf 位相に対する層なので、$y_i = y \circ g_i$ を満たす一意な射
$y : T \to Y$ が $y_i$ から得られる。このとき
$y_{i, small}^{-1}\mathcal{F} =
g_{i, small}^{-1}y_{small}^{-1}\mathcal{F}$ である。
したがって 'Etale Cohomology, Lemma
\ref{etale-cohomology-lemma-describe-pullback} により、切断 $s_i$ は
$y_{small}^{-1}\mathcal{F}$ の切断へ一意に貼り合わさる。

\medskip\noindent
$\mathcal{F} \in \Ob(\Sh(Y_\etale))$ を
$X \in \Ob((\Sch/S)_{fppf})$ へ写す構成は有限極限とすべての余極限を保つ。
実際、各関手 $y_{small}^{-1}$ がこの性質をもつ。
もちろん、$V \in \Ob(Y_\etale)$ ならば、この構成は
$Y_\etale$ 上の表現可能層 $h_V$ を、$V$ によって表現される
表現可能関手へ写す。

\medskip\noindent
Sites, Lemma \ref{sites-lemma-sheaf-coequalizer-representable} により、
集合 $I$、各 $i \in I$ に対する $Y_\etale$ の対象 $V_i$、および
$Y_\etale$ 上の層の全射
$$
\coprod h_{V_i} \longrightarrow \mathcal{F}
$$
を見つけられる。必要な集合論的条件は、添字集合 $I$ が大きすぎないことである
\footnote{$Y$ の幾何学的点における $\mathcal{F}$ の茎の濃度の上限が、
$(\Sch/S)_{fppf}$ のある対象の大きさで抑えられれば十分である。}。
このとき $V = \coprod V_i$ は $(\Sch/S)_{fppf}$ の対象であり、
したがって $Y_\etale$ の対象である。また全射 $h_V \to \mathcal{F}$ を得る。

\medskip\noindent
$\Sh(Y_\etale)$ における $h_V$ とそれ自身との積は
$h_{V \times_Y V}$ であることに注意する。ファイバー積
$$
h_V \times_\mathcal{F} h_V \subset h_{V \times_Y V}
$$
を考える。$V \times_Y V$ の開部分スキーム $R$ であって
$h_V \times_\mathcal{F} h_V = h_R$ となるものが存在する。
Lemma \ref{spaces-properties-lemma-support-subsheaf-final} を参照せよ（細部を省略する）。
Yoneda の補題により $Y_\etale$ の二つの射 $s, t : R \to V$ を得て、
$\Sh(Y_\etale)$ における余等化子図式
$$
\xymatrix{
h_R \ar@<1ex>[r] \ar@<-1ex>[r] &
h_V \ar[r] &
\mathcal{F}
}
$$
を得る。もちろん射 $s, t$ はエタールであり、エタール同値関係
$(t, s) : R \to V \times_S V$ を定める。

\medskip\noindent
直前の二段落の議論により、$(\Sch/S)_{fppf}$ における余等化子図式
$$
\xymatrix{
R \ar@<1ex>[r] \ar@<-1ex>[r] &
V \ar[r] &
X
}
$$
を得る。したがって Spaces, Theorem
\ref{spaces-theorem-presentation} により $X = V/R$ は代数空間である。
これで (1) が証明された。(2) は $V \to Y$ がエタールであることから従う。
(3) は $X$ と $h_X$ の定義から直ちに従う。(4) の証明は省略する。
これは Lemma \ref{spaces-properties-lemma-etale-morphism-topoi} の余連続関手 $j$ に付随する射を、
上で論じた $Y_{spaces, \etale}$ の $X$ における局所化としての
$X_{spaces, \etale}$ の記述、および Sites, Lemma
\ref{sites-lemma-describe-j-shriek-representable} と照合すれば従う。
\end{proof}

\section{射の復元}
\label{spaces-properties-section-morphisms}

\noindent
この節では、代数空間にその局所環付き小エタールトポスを対応させる規則が、
適切な意味で充満忠実であることを証明する。
Theorem \ref{spaces-properties-theorem-fully-faithful} を参照せよ。

\begin{lemma}
\label{spaces-properties-lemma-morphism-locally-ringed}
$S$ をスキームとする。
$f : X \to Y$ を $S$ 上の代数空間の射とする。
$f$ に付随する環付きトポスの射 $(f_{small}, f^\sharp)$ は
局所環付きトポスの射である。Modules on Sites, Definition
\ref{sites-modules-definition-morphism-locally-ringed-topoi} を参照せよ。
\end{lemma}

\begin{proof}
$(X_\etale, \mathcal{O}_{X_\etale})$ と
$(Y_\etale, \mathcal{O}_{Y_\etale})$ が局所環付きサイトであることを
すでに見たので、この主張は意味をもつ。
Lemma \ref{spaces-properties-lemma-etale-site-locally-ringed} を参照せよ。
さらに $X_\etale$ が十分多くの点をもつことも既知である。
Theorem \ref{spaces-properties-theorem-exactness-stalks} を参照せよ。
したがって $(f_{small}, f^\sharp)$ が Modules on Sites, Lemma
\ref{sites-modules-lemma-locally-ringed-morphism} の条件 (3) を
満たすことを証明すれば十分である。
そのため $X_\etale$ の点 $p$ を取る。
Lemma \ref{spaces-properties-lemma-points-small-etale-site} により、$p$ は $X$ の
幾何学的点 $\overline{x}$ に対応する。
Lemma \ref{spaces-properties-lemma-stalk-pullback} により、点
$q = f_{small} \circ p$ は $Y$ の幾何学的点
$\overline{y} = f \circ \overline{x}$ に対応する。
したがって証明すべき主張は、誘導されるエタール局所環の写像
$$
\mathcal{O}_{Y, \overline{y}} \longrightarrow \mathcal{O}_{X, \overline{x}}
$$
が局所環準同型であることだ。これは直接証明できるが、代わりに
スキームに対する対応する結果から導こう。そのため可換図式
$$
\xymatrix{
U \ar[d] \ar[r]_\psi & V \ar[d] \\
X \ar[r] & Y
}
$$
であって、$U$ と $V$ がスキーム、縦の矢印が全射エタールであるものを選ぶ
（Spaces, Lemma \ref{spaces-lemma-lift-morphism-presentations} を参照）。
持ち上げ $\overline{u} : \overline{x} \to U$ を選ぶ
（Lemma \ref{spaces-properties-lemma-geometric-lift-to-cover} により可能）。
$\overline{v} = \psi \circ \overline{u}$ とおく。
エタール局所環の可換図式
$$
\xymatrix{
\mathcal{O}_{U, \overline{u}} &
\mathcal{O}_{V, \overline{v}} \ar[l] \\
\mathcal{O}_{X, \overline{x}} \ar[u] &
\mathcal{O}_{Y, \overline{y}}. \ar[l] \ar[u]
}
$$
を得る。'Etale Cohomology, Lemma
\ref{etale-cohomology-lemma-morphism-locally-ringed} により、
上の横矢印は局所環準同型である。
最後に Lemma \ref{spaces-properties-lemma-describe-etale-local-ring} により縦矢印は同型である。
したがって結論を得る。
\end{proof}

\begin{lemma}
\label{spaces-properties-lemma-2-morphism}
$S$ をスキームとする。
$X$, $Y$ を $S$ 上の代数空間とする。
$f : X \to Y$ を $S$ 上の代数空間の射とする。
$t$ を $(f_{small}, f^\sharp)$ からそれ自身への $2$-射とする。
Modules on Sites, Definition
\ref{sites-modules-definition-2-morphism-ringed-topoi} を参照せよ。
このとき $t = \text{id}$ である。
\end{lemma}

\begin{proof}
$X'$（それぞれ $Y'$）を、$\Spec(\mathbf{Z})$ 上の代数空間とみなした
$X$ とする。Spaces, Definition \ref{spaces-definition-base-change} を参照せよ。
構成から $(X_{small}, \mathcal{O})$ は
$(X'_{small}, \mathcal{O})$ に等しく、$Y$ についても同様であることは明らかである。
したがって $X'$ と $Y'$ を用いてよい。言い換えると、
$S = \Spec(\mathbf{Z})$ と仮定してよい。

\medskip\noindent
$S = \Spec(\mathbf{Z})$ とし、$f : X \to Y$ と $t$ は補題のとおりとする。
これは $t : f^{-1}_{small} \to f^{-1}_{small}$ が関手の変換であり、図式
$$
\xymatrix{
f_{small}^{-1}\mathcal{O}_Y
\ar[rd]_{f^\sharp}  & &
f_{small}^{-1}\mathcal{O}_Y \ar[ll]^t \ar[ld]^{f^\sharp} \\
& \mathcal{O}_X
}
$$
が可換であることを意味する。
$V \to Y$ がエタールで、$V$ がアフィンであるとする。
$V = \Spec(B)$ と書く。$\mathbf{Z}$-代数としての $B$ の生成元
$b_j \in B$, $j \in J$ を選ぶ。
$T = \Spec(\mathbf{Z}[\{x_j\}_{j \in J}])$ とおく。
以下では、任意のスキーム $U$ に対する等式
$\Mor_{\Sch}(U, T) = \prod_{j \in J} \Gamma(U, \mathcal{O}_U)$ を
特に断らず用いる。
全射環準同型 $\mathbf{Z}[x_j] \to B$, $x_j \mapsto b_j$ は
閉埋め込み $V \to T$ に対応する。
$Y$ 上の代数空間の単射
$$
i : V \longrightarrow T_Y = T \times Y
$$
を得る。$Y_\etale$ 上の層の言葉では、射 $i$ は層の単射
$h_i : h_V \to \prod_{j \in J} \mathcal{O}_Y$ を誘導する。
$i$ の $X$ への基底変換
$i' : X \times_Y V \to T_X$ も単射である
（Spaces, Lemma
\ref{spaces-lemma-base-change-representable-transformations-property}）。
したがって $i' : X \times_Y V \to T_X$ は単射であり、これはさらに
$h_{i'} : h_{X \times_Y V} \to \prod_{j \in J} \mathcal{O}_X$ が
層の単射であることを意味する。
Lemma \ref{spaces-properties-lemma-stalk-pullback} の同一視
$f_{small}^{-1}h_V = h_{X \times_Y V}$ を介して、写像 $h_{i'}$ は
$$
\xymatrix{
f_{small}^{-1}h_V \ar[r]^-{f^{-1}h_i} &
\prod_{j \in J} f_{small}^{-1}\mathcal{O}_Y
\ar[r]^{\prod f^\sharp} &
\prod_{j \in J} \mathcal{O}_X
}
$$
に等しい（確認は省略する）。これは写像
$t : f_{small}^{-1}h_V \to f_{small}^{-1}h_V$ が可換図式
$$
\xymatrix{
f_{small}^{-1}h_V \ar[r]^-{f^{-1}h_i} \ar[d]^t &
\prod_{j \in J} f_{small}^{-1}\mathcal{O}_Y
\ar[r]^-{\prod f^\sharp} \ar[d]^{\prod t} &
\prod_{j \in J} \mathcal{O}_X \ar[d]^{\text{id}}\\
f_{small}^{-1}h_V \ar[r]^-{f^{-1}h_i} &
\prod_{j \in J} f_{small}^{-1}\mathcal{O}_Y
\ar[r]^-{\prod f^\sharp} &
\prod_{j \in J} \mathcal{O}_X
}
$$
に収まることを意味する。右の正方形の可換性は、上で説明した
$t$ に関する仮定から成り立つ。
上の議論により横矢印の合成は単射なので、左の縦矢印も恒等写像である。
$Y_\etale$ 上の任意の集合の層は、$V$ がアフィンであるような型 $h_V$ の
層の（巨大な）余積からの全射をもつ
（Lemma \ref{spaces-properties-lemma-alternative} と Sites, Lemma
\ref{sites-lemma-sheaf-coequalizer-representable} を組み合わせよ）。
したがって所望のとおり
$t : f_{small}^{-1} \to f_{small}^{-1}$ は恒等変換である。
\end{proof}

\begin{lemma}
\label{spaces-properties-lemma-faithful}
$S$ をスキームとする。
$X$, $Y$ を $S$ 上の代数空間とする。
$S$ 上の代数空間の二つの射 $a, b : X \to Y$ について、
環付きトポスの $2$-圏における $2$-同型
$(a_{small}, a^\sharp) \cong (b_{small}, b^\sharp)$ が存在するならば、
両者は等しい。
\end{lemma}

\begin{proof}
$t : a_{small}^{-1} \to b_{small}^{-1}$ をこの $2$-同型とする。
$X_\etale$ 上の層と $X_{spaces, \etale}$ 上の層の間に違いはないので、
同値なことに $t$ を変換
$t : a_{spaces, \etale}^{-1} \to b_{spaces, \etale}^{-1}$ と考えられる。
$U$ と $V$ がスキームで、$p$ と $q$ が全射エタールであるような可換図式
$$
\xymatrix{
U \ar[d]_p \ar[r]_\alpha  & V \ar[d]^q \\
X \ar[r]^a & Y
}
$$
を選ぶ。図式
$$
\xymatrix{
h_U \ar[r]_-\alpha \ar@{=}[d] & a_{spaces, \etale}^{-1}h_V \ar[d]^t \\
h_U \ar@{..>}[r] & b_{spaces, \etale}^{-1}h_V
}
$$
を考える。層 $b_{spaces, \etale}^{-1}h_V$ は
$h_{V \times_{Y, b} X}$ と同型なので、点線の矢印は可換図式
$$
\xymatrix{
U \ar[d]_p \ar[r]_\beta  & V \ar[d]^q \\
X \ar[r]^b & Y
}
$$
に収まるスキームの射 $\beta : U \to V$ から得られる。
次の $2$-同型の列が存在すると主張する。
\begin{align*}
(\alpha_{small}, \alpha^\sharp)
& \cong
(\alpha_{spaces, \etale}, \alpha^\sharp) \\
& \cong
(a_{spaces, \etale, c}, a_c^\sharp) \\
& \cong
(b_{spaces, \etale, d}, b_d^\sharp) \\
& \cong
(\beta_{spaces, \etale}, \beta^\sharp) \\
& \cong
(\beta_{small}, \beta^\sharp)
\end{align*}
最初と最後の $2$-同型は、$U_{spaces, \etale}$ 上の層と
$U_\etale$ 上の層との同一視、および $V$ に対する同様の同一視から得られる。
第二と第四の $2$-同型は、$\alpha$ によって誘導される
$c : U \to X \times_{a, Y} V$ と、$\beta$ によって誘導される
$d : U \to X \times_{b, Y} V$ に対する
Lemma \ref{spaces-properties-lemma-relocalize-morphism} の同型である。
中央の $2$-同型は変換 $t$ から得られる。
実際、関手 $a_{spaces, \etale, c}^{-1}$ は関手
$$
(\mathcal{H} \to h_V) \longmapsto
(a_{spaces, \etale}^{-1}\mathcal{H}
\times_{a_{spaces, \etale}^{-1}h_V, \alpha}
h_U \to
h_U)
$$
に対応し、$b_{spaces, \etale, d}^{-1}$ についても同様である。
Sites, Lemma \ref{sites-lemma-relocalize-morphism} を参照せよ。
ここでは $Y_{spaces, \etale}/V$ 上の層を
$\Sh(Y_{spaces, \etale})$ の矢印 $(\mathcal{H} \to h_V)$ として同一視し、
$U/X$ についても同様にしている。Sites, Lemma
\ref{sites-lemma-essential-image-j-shriek} を参照せよ。
この同一視を介して、構造層 $\mathcal{O}_V$ は対
$(\mathcal{O}_Y \times h_V \to h_V)$ に対応し、
$\mathcal{O}_U$ についても同様である。
Modules on Sites, Lemma
\ref{sites-modules-lemma-localize-compare} を参照せよ。
$t$ は $\alpha$ と $\beta$ を入れ替えるので、$t$ は同型
$$
t :
a_{spaces, \etale}^{-1}\mathcal{H}
\times_{a_{spaces, \etale}^{-1}h_V, \alpha}
h_U
\longrightarrow
b_{spaces, \etale}^{-1}\mathcal{H}
\times_{b_{spaces, \etale}^{-1}h_V, \beta}
h_U
$$
を $h_U$ 上、$(\mathcal{H} \to h_V)$ に関して関手的に誘導する。
また、上の構造層 $\mathcal{O}_U$ と $\mathcal{O}_V$ の記述により、
$t$ は $a^\sharp$ と $b^\sharp$ と両立するので、
$t$ は $a_c^\sharp$ と $b_d^\sharp$ とも両立する。
したがって環付きトポスの射
$(\alpha_{small}, \alpha^\sharp)$ と
$(\beta_{small}, \beta^\sharp)$ は $2$-同型である。
'Etale Cohomology, Lemma \ref{etale-cohomology-lemma-faithful} により
$\alpha = \beta$ と結論する！
$p : U \to X$ は層の全射なので、$a = b$ が従う。
\end{proof}

\noindent
この節の主結果は次のとおりである。

\begin{theorem}
\label{spaces-properties-theorem-fully-faithful}
$X$, $Y$ を $\Spec(\mathbf{Z})$ 上の代数空間とする。
$$
(g, g^\sharp) :
(\Sh(X_\etale), \mathcal{O}_X)
\longrightarrow
(\Sh(Y_\etale), \mathcal{O}_Y)
$$
を局所環付きトポスの射とする。このとき、
$(g, g^\sharp)$ が $(f_{small}, f^\sharp)$ と同型になるような
代数空間の一意な射 $f : X \to Y$ が存在する。
言い換えると、構成
$$
\textit{Spaces}/\Spec(\mathbf{Z})
\longrightarrow \textit{Locally ringed topoi},
\quad
X \longrightarrow (X_\etale, \mathcal{O}_X)
$$
は充満忠実である（右辺の射は $2$-同型を除いて考える）。
\end{theorem}

\begin{proof}
一意性は Lemma \ref{spaces-properties-lemma-faithful} で見た。
したがって存在を証明すれば十分である。
この証明では Equation (\ref{spaces-properties-equation-localize-at-scheme-ringed}) の同一視と、
Lemma \ref{spaces-properties-lemma-relocalize-morphism-at-schemes} の結果を自由に用いる。

\medskip\noindent
$U \in \Ob(X_\etale)$, $V \in \Ob(Y_\etale)$ とし、
$s \in g^{-1}h_V(U)$ を切断とする。
$s$ を層の写像 $s : h_U \to g^{-1}h_V$ と考えられる。
Modules on Sites, Lemma
\ref{sites-modules-lemma-relocalize-morphism-ringed-topoi} により、
環付きトポスの射の可換図式
$$
\xymatrix{
(\Sh(X_\etale/U), \mathcal{O}_U)
\ar[rr]_-{(j, j^\sharp)} \ar[d]_{(g_s, g_s^\sharp)} & &
(\Sh(X_\etale), \mathcal{O}_X) \ar[d]^{(g, g^\sharp)} \\
(\Sh(V_\etale), \mathcal{O}_V) \ar[rr] & &
(\Sh(Y_\etale), \mathcal{O}_Y).
}
$$
を得る。'Etale Cohomology, Theorem
\ref{etale-cohomology-theorem-fully-faithful} により、
$(g_s, g_s^\sharp)$ が $(f_{s, small}, f_s^\sharp)$ と $2$-同型になるような
スキームの一意な射 $f_s : U \to V$ を得る。
以上で説明した構成 $(U, V, s) \leadsto f_s$ は次の関手性を満たす。
$X_\etale$ の射 $a : U' \to U$、$Y_\etale$ の射 $b : V' \to V$、および
図式
$$
\xymatrix{
h_{U'} \ar[d]_a \ar[r]_{s'} & g^{-1}h_{V'} \ar[d]^{g^{-1}b} \\
h_U \ar[r]^s & g^{-1}h_V
}
$$
が可換となる写像 $s' : h_{U'} \to g^{-1}h_{V'}$ が与えられたとする。
このときスキームの図式
$$
\xymatrix{
U' \ar[r]_-{f_{s'}} \ar[d]_a & u(V') \ar[d]^{u(b)} \\
U \ar[r]^-{f_s} & u(V)
}
$$
は可換である。その理由は、Modules on Sites, Lemma
\ref{sites-modules-lemma-relocalize-morphism-ringed-topoi} で構成された射
$(g_s, g_s^\sharp)$ に同じ条件が成り立つことと、'Etale Cohomology,
Theorem \ref{etale-cohomology-theorem-fully-faithful} の一意性である。

\medskip\noindent
問題は射 $f_s$ を代数空間の射へ貼り合わせることである。
そのため、まずスキーム $V$ と全射エタール射 $V \to Y$ を選ぶ。
これは $h_V \to *$ が全射であることを意味するので、
$g^{-1}h_V \to *$ も全射である。
したがってスキーム $U$、全射エタール射 $U \to X$、および射
$s : h_U \to g^{-1}h_V$ が存在する。
次に $R = V \times_Y V$, $R' = U \times_X U$ とおく。
$g^{-1}$ は完全なので
$g^{-1}h_R = g^{-1}h_V \times g^{-1}h_V$ を得る。
したがって $s$ は射 $s \times s : h_{R'} \to g^{-1}h_R$ を誘導する。
上の構成を適用すると、スキームの射の可換図式
$$
\xymatrix{
R' \ar@<1ex>[d] \ar@<-1ex>[d] \ar[rr]_{f_{s \times s}} & &
R \ar@<1ex>[d] \ar@<-1ex>[d] \\
U \ar[rr]^{f_s} & &
V
}
$$
を得る。$X = U/R'$ および $Y = V/R$ なので
（Spaces, Lemma \ref{spaces-lemma-space-presentation} を参照）、
この図式は明らかな可換図式に収まる代数空間の射
$f : X \to Y$ を定める。
なお $(f_{small}, f^\sharp)$ が $(g, g^\sharp)$ と $2$-同型であることを
示さなければならない。上の構成によって与えられた $2$-同型を
$t_V : f_{s, small}^{-1} \to g_s^{-1}$ および
$t_R : f_{s \times s, small}^{-1} \to g_{s \times s}^{-1}$ とする。
$\mathcal{G}$ を $Y_\etale$ 上の層とする。このとき $t_V$ は同型
$$
f_{small}^{-1}\mathcal{G}|_{U_\etale}
=
f_{s, small}^{-1}\mathcal{G}|_{V_\etale}
\xrightarrow{t_V}
g_s^{-1}\mathcal{G}|_{V_\etale}
=
g^{-1}\mathcal{G}|_{U_\etale}.
$$
を定める。さらに、この同型をいずれかの射影 $R' \to U$ によって
$R'$ へ引き戻したものは同型
$$
f_{small}^{-1}\mathcal{G}|_{R'_\etale}
=
f_{s \times s, small}^{-1}\mathcal{G}|_{R_\etale}
\xrightarrow{t_R}
g_{s \times s}^{-1}\mathcal{G}|_{R_\etale}
=
g^{-1}\mathcal{G}|_{R'_\etale}.
$$
である。$\{U \to X\}$ はサイト $X_{spaces, \etale}$ の被覆なので、
これは最初の表示された同型が層の同型
$t : f_{small}^{-1}\mathcal{G} \to g^{-1}\mathcal{G}$ へ降下することを
意味する（細部を省略する）。$t_V$ と $t_R$ は関手の変換なので、
この同型は $\mathcal{G}$ に関して関手的である。
最後に、$t_V$ と $t_R$ がそれぞれ両立するので、$t$ は
$f^\sharp$ と $g^\sharp$ と両立する（細部をいくつか省略する）。
これで定理の証明は完了する。
\end{proof}

\begin{lemma}
\label{spaces-properties-lemma-isomorphism-ringed-topoi}
$X$, $Y$ を $\mathbf{Z}$ 上の代数空間とする。
$$
(g, g^\sharp) :
(\Sh(X_\etale), \mathcal{O}_X)
\longrightarrow
(\Sh(Y_\etale), \mathcal{O}_Y)
$$
が環付きトポスの同型ならば、$(g, g^\sharp)$ が
$(f_{small}, f^\sharp)$ と同型になるような代数空間の一意な射
$f : X \to Y$ が存在し、さらに $f$ は代数空間の同型である。
\end{lemma}

\begin{proof}
Theorem \ref{spaces-properties-theorem-fully-faithful} により、$(g, g^\sharp)$ が
局所環付きトポスの射であることを示せば十分である。
Modules on Sites, Lemma
\ref{sites-modules-lemma-locally-ringed-morphism} により
（またサイト $X_\etale$ は十分多くの点をもつので）、
$X_\etale$ の任意の点 $p$ と $q = f \circ p$ に対して、
$g^\sharp$ により誘導される写像
$\mathcal{O}_{Y, q} \to \mathcal{O}_{X, p}$ が局所環準同型であることを
確認すれば十分である。これは同型なので明らかである。
\end{proof}

\section{代数空間上の準連接層}
\label{spaces-properties-section-quasi-coherent}

\noindent
Descent, Sections \ref{descent-section-quasi-coherent-sheaves},
\ref{descent-section-quasi-coherent-cohomology}, および
\ref{descent-section-quasi-coherent-sheaves-bis} で、スキーム $U$ に対して、
$U$ 上の準連接 $\mathcal{O}_U$-加群と $U$ の小エタールサイト上の
準連接 $\mathcal{O}$-加群との間に違いがないことを見た。
したがって次の定義を表現可能代数空間に適用すると、
スキーム上の準連接層という元来の概念
（Schemes, Section \ref{schemes-section-quasi-coherent}）と両立する。

\begin{definition}
\label{spaces-properties-definition-quasi-coherent}
$S$ をスキーム、$X$ を $S$ 上の代数空間とする。
{\it 準連接} $\mathcal{O}_X$-加群とは、Modules on Sites, Definition
\ref{sites-modules-definition-site-local} の意味で環付きサイト
$(X_\etale, \mathcal{O}_X)$ 上の準連接加群であるものをいう。
$X$ 上の準連接層の圏を $\QCoh(\mathcal{O}_X)$ と書く。
\end{definition}

\noindent
準連接性は内在的な概念なので（Modules on Sites, Lemma
\ref{sites-modules-lemma-special-locally-free} を参照）、これは、
対応する $X_{spaces, \etale}$ 上の $\mathcal{O}_X$-加群が準連接であることと
同値である。

\medskip\noindent
通常どおり、準連接層は引き戻しに関してよく振る舞う。

\begin{lemma}
\label{spaces-properties-lemma-pullback-quasi-coherent}
$S$ をスキームとする。
$f : X \to Y$ を $S$ 上の代数空間の射とする。
引き戻し関手
$f^* : \textit{Mod}(\mathcal{O}_Y) \to \textit{Mod}(\mathcal{O}_X)$ は
準連接層を保つ。
\end{lemma}

\begin{proof}
これは一般的な事実である。Modules on Sites, Lemma
\ref{sites-modules-lemma-local-pullback} を参照せよ。
\end{proof}

\noindent
$X$ と $Y$ がたまたまスキームならば、この引き戻し関手は
準連接加群の層の間の通常の引き戻し関手と一致することに注意する。
Descent, Proposition
\ref{descent-proposition-equivalence-quasi-coherent-functorial} を参照せよ。
これと $X$ の小エタールサイトの対象上の準連接層とを比較する
定番の補題は次のとおりである。

\begin{lemma}
\label{spaces-properties-lemma-characterize-quasi-coherent-small-etale}
$S$ をスキーム、$X$ を $S$ 上の代数空間とする。
準連接 $\mathcal{O}_X$-加群 $\mathcal{F}$ は次のデータによって与えられる。
\begin{enumerate}
\item 各 $U \in \Ob(X_\etale)$ に対する、$U_\etale$ 上の準連接
$\mathcal{O}_U$-加群 $\mathcal{F}_U$。
\item $X_\etale$ の各 $f : U' \to U$ に対する同型
$c_f : f_{small}^*\mathcal{F}_U \to \mathcal{F}_{U'}$。
\end{enumerate}
これらのデータには、$X_\etale$ の任意の $f : U' \to U$ と
$g : U'' \to U'$ に対して、合成 $c_g \circ g_{small}^*c_f$ が
$c_{f \circ g}$ に等しいという条件を課す。
\end{lemma}

\begin{proof}
Lemmas \ref{spaces-properties-lemma-pullback-quasi-coherent} および
\ref{spaces-properties-lemma-characterize-module-small-etale} を組み合わせる。
\end{proof}

\begin{lemma}
\label{spaces-properties-lemma-stalk-quasi-coherent}
$S$ をスキーム、$X$ を $S$ 上の代数空間とする。
$\mathcal{F}$ を準連接 $\mathcal{O}_X$-加群とする。
$x \in |X|$ を点とし、$\overline{x}$ を $x$ の上にある幾何学的点とする。
最後に、$\varphi : (U, \overline{u}) \to (X, \overline{x})$ を、
$U$ がスキームであるようなエタール近傍とする。このとき
$$
(\varphi^*\mathcal{F})_u \otimes_{\mathcal{O}_{U, u}}
\mathcal{O}_{X, \overline{x}} =
\mathcal{F}_{\overline{x}}
$$
である。ここで $u \in U$ は $\overline{u}$ の像である。
\end{lemma}

\begin{proof}
Lemma \ref{spaces-properties-lemma-describe-etale-local-ring} により
$\mathcal{O}_{X, \overline{x}} = \mathcal{O}_{U, u}^{sh}$ なので、
テンソル積は意味をもつ。さらに Definition \ref{spaces-properties-definition-stalk} から
$$
\mathcal{F}_{\overline{u}} = \colim (\varphi^*\mathcal{F})_u
$$
であることは明らかである。ここで余極限は補題のような
$\varphi : (U, \overline{u}) \to (X, \overline{x})$ を走る。
したがって補題の主張の左辺から右辺への標準的な写像がある。
$\mathcal{O}_{X, \overline{x}}$ にも同様の余極限による記述があり、
Lemma \ref{spaces-properties-lemma-characterize-quasi-coherent-small-etale} により、
$(U', \overline{u}') \to (U, \overline{u})$ がエタール近傍の射であるたびに
$$
((\varphi')^*\mathcal{F})_{u'} =
(\varphi^*\mathcal{F})_u \otimes_{\mathcal{O}_{U, u}} \mathcal{O}_{U', u'}
$$
である。証明を完了するため、$\otimes$ が余極限と可換することを用いる。
\end{proof}

\begin{lemma}
\label{spaces-properties-lemma-stalk-pullback-quasi-coherent}
$S$ をスキームとする。
$f : X \to Y$ を $S$ 上の代数空間の射とする。
$\mathcal{G}$ を準連接 $\mathcal{O}_Y$-加群とする。
$\overline{x}$ を $X$ の幾何学的点とし、
$\overline{y} = f \circ \overline{x}$ を $Y$ におけるその像とする。
このとき、引き戻しの茎と、茎と $\overline{x}$ における $X$ の局所環との
テンソル積との間に標準的な同型
$$
(f^*\mathcal{G})_{\overline{x}} =
\mathcal{G}_{\overline{y}} \otimes_{\mathcal{O}_{Y, \overline{y}}}
\mathcal{O}_{X, \overline{x}}
$$
が存在する。
\end{lemma}

\begin{proof}
$f^*\mathcal{G} =
f_{small}^{-1}\mathcal{G} \otimes_{f_{small}^{-1}\mathcal{O}_Y} \mathcal{O}_X$
なので、これは Lemma \ref{spaces-properties-lemma-stalk-pullback} の引き戻しの茎の記述と、
茎を取ることがテンソル積と可換するという事実から従う。
より直接的には次のように分かる。$U$ と $V$ がスキームで、
$p$ と $q$ が全射エタールであるような可換図式
$$
\xymatrix{
U \ar[d]_p \ar[r]_\alpha  & V \ar[d]^q \\
X \ar[r]^a & Y
}
$$
を選ぶ。Lemma \ref{spaces-properties-lemma-geometric-lift-to-usual} により、
$\overline{x} = p \circ \overline{u}$ を満たす $U$ の幾何学的点
$\overline{u}$ を選べる。$\overline{v} = \alpha \circ \overline{u}$ とおく。
すると
\begin{align*}
(f^*\mathcal{G})_{\overline{x}} & =
(p^*f^*\mathcal{G})_u \otimes_{\mathcal{O}_{U, u}}
\mathcal{O}_{X, \overline{x}} \\
& = (\alpha^*q^*\mathcal{G})_u \otimes_{\mathcal{O}_{U, u}}
\mathcal{O}_{X, \overline{x}} \\
& = (q^*\mathcal{G})_v \otimes_{\mathcal{O}_{V, v}}
\mathcal{O}_{U, u} \otimes_{\mathcal{O}_{U, u}}
\mathcal{O}_{X, \overline{x}} \\
& = (q^*\mathcal{G})_v \otimes_{\mathcal{O}_{V, v}}
\mathcal{O}_{X, \overline{x}} \\
& = (q^*\mathcal{G})_v \otimes_{\mathcal{O}_{V, v}}
\mathcal{O}_{Y, \overline{y}} \otimes_{\mathcal{O}_{Y, \overline{y}}}
\mathcal{O}_{X, \overline{x}} \\
& = \mathcal{G}_{\overline{y}} \otimes_{\mathcal{O}_{Y, \overline{y}}}
\mathcal{O}_{X, \overline{x}}
\end{align*}
である。ここでは Lemma \ref{spaces-properties-lemma-stalk-quasi-coherent} を（二度）用い、
またスキーム上の準連接層の引き戻しに対する対応する結果
Sheaves, Lemma \ref{sheaves-lemma-stalk-pullback-modules} を用いた。
\end{proof}

\begin{lemma}
\label{spaces-properties-lemma-characterize-quasi-coherent}
$S$ をスキーム、$X$ を $S$ 上の代数空間とする。
$\mathcal{F}$ を $\mathcal{O}_X$-加群の層とする。次は同値である。
\begin{enumerate}
\item $\mathcal{F}$ は準連接 $\mathcal{O}_X$-加群である。
\item $S$ 上の代数空間のエタール射 $f : Y \to X$ であって、
$|f| : |Y| \to |X|$ が全射であり、$f^*\mathcal{F}$ が $Y$ 上
準連接となるものが存在する。
\item スキーム $U$ と全射エタール射 $\varphi : U \to X$ であって、
$\varphi^*\mathcal{F}$ が準連接 $\mathcal{O}_U$-加群となるものが存在する。
\item 任意のアフィンスキーム $U$ とエタール射 $\varphi : U \to X$ に対して、
制限 $\varphi^*\mathcal{F}$ は準連接 $\mathcal{O}_U$-加群である。
\end{enumerate}
\end{lemma}

\begin{proof}
$\text{id}_X$ を考えれば (1) が (2) を含意することは明らかである。
$f : Y \to X$ は (2) のとおりとし、$V \to Y$ をスキームから $Y$ への
全射エタール射とする。このとき合成 $V \to X$ も全射エタールであり、
Lemma \ref{spaces-properties-lemma-pullback-quasi-coherent} により $\mathcal{F}$ の $V$ への
引き戻しも準連接である。したがって (2) は (3) を含意する。

\medskip\noindent
$U \to X$ は (3) のとおりとする。
Equation (\ref{spaces-properties-equation-restrict-modules}) で導入した記号の濫用を用いる。
$\mathcal{F}|_{U_\etale}$ は準連接なので、各
$\mathcal{F}|_{U_{i, \etale}}$ が大域表示をもつようなエタール被覆
$\{U_i \to U\}$ が存在する。Modules on Sites, Definition
\ref{sites-modules-definition-global} および Lemma
\ref{sites-modules-lemma-local-final-object} を参照せよ。
$V \to X$ を $X_\etale$ の対象とする。
$U \to X$ は全射かつエタールなので、写像の族
$\{U_i \times_X V \to V\}$ は $V$ のエタール被覆である。
射 $U_i \times_X V \to U_i$ を介して、
$\mathcal{F}|_{U_{i, \etale}}$ の大域表示を制限し、
$\mathcal{F}|_{(U_i \times_X V)_\etale}$ の大域表示を得られる。
したがって $X_\etale$ 上の層 $\mathcal{F}$ は Modules on Sites,
Definition \ref{sites-modules-definition-site-local} の条件を満たし、
準連接である。

\medskip\noindent
(3) と (4) の同値性は、任意のスキームがアフィン開被覆をもつことから従う。
\end{proof}

\begin{lemma}
\label{spaces-properties-lemma-properties-quasi-coherent}
$S$ をスキーム、$X$ を $S$ 上の代数空間とする。
$X$ 上の準連接層の圏 $\QCoh(\mathcal{O}_X)$ は次の性質をもつ。
\begin{enumerate}
\item 準連接層の任意の直和は準連接である。
\item 準連接層の任意の余極限は準連接である。
\item 準連接層の射の核と余核は準連接である。
\item $\mathcal{O}_X$-加群の短完全列
$0 \to \mathcal{F}_1 \to \mathcal{F}_2 \to \mathcal{F}_3 \to 0$ が
与えられたとき、三つのうち二つが準連接ならば第三も準連接である。
\item 二つの準連接 $\mathcal{O}_X$-加群のテンソル積は準連接である。
\item 二つの準連接 $\mathcal{O}_X$-加群 $\mathcal{F}$, $\mathcal{G}$ が
与えられ、$\mathcal{F}$ が有限表示であるとする
（Section \ref{spaces-properties-section-properties-modules} を参照）。このとき内部 Hom
$\SheafHom_{\mathcal{O}_X}(\mathcal{F}, \mathcal{G})$ は準連接である。
\end{enumerate}
\end{lemma}

\begin{proof}
$X$ がスキームならば、これは Descent, Lemma
\ref{descent-lemma-properties-quasi-coherent} である。
エタール局所化により補題をこの場合へ帰着させる。

\medskip\noindent
スキーム $U$ と全射エタール射 $\varphi : U \to X$ を選ぶ。
記号は $\textit{Mod}(\mathcal{O}_U) =
\textit{Mod}(U_\etale, \mathcal{O}_U)$ および
$\QCoh(\mathcal{O}_U) = \QCoh(U_\etale, \mathcal{O}_U)$ とする。
言い換えると、$U$ はスキームであるが、$U$ 上の準連接加群を
$U$ の小エタールサイト上の加群と考える。
Lemma \ref{spaces-properties-lemma-pullback-quasi-coherent} により可換図式
$$
\xymatrix{
\QCoh(\mathcal{O}_X) \ar[r]_{\varphi^*} \ar[d] &
\QCoh(\mathcal{O}_U) \ar[d] \\
\textit{Mod}(\mathcal{O}_X) \ar[r]^{\varphi^*} &
\textit{Mod}(\mathcal{O}_U)
}
$$
を得る。下の横矢印は制限関手
(\ref{spaces-properties-equation-restrict-modules})
$\mathcal{G} \mapsto \mathcal{G}|_{U_\etale}$ である。
この関手は左随伴と右随伴の両方をもつので、すべての極限および余極限と可換する。
Modules on Sites, Section \ref{sites-modules-section-localize} を参照せよ。
さらに Lemma \ref{spaces-properties-lemma-characterize-quasi-coherent} により、
$\textit{Mod}(\mathcal{O}_X)$ の対象が $\QCoh(\mathcal{O}_X)$ に属するための
必要十分条件は、その $U$ への制限が $\QCoh(\mathcal{O}_U)$ に属することである。
以上の準備のもとで証明を始められる。

\medskip\noindent
(1) の証明。$\mathcal{F}_i$, $i \in I$ を準連接
$\mathcal{O}_X$-加群の族とする。上の議論により
$$
\Big(\bigoplus \mathcal{F}_i\Big)|_{U_\etale} =
\bigoplus \mathcal{F}_i|_{U_\etale}
$$
である。加群 $\mathcal{F}_i|_{U_\etale}$ はいずれも準連接である。
したがってスキームの場合により直和は準連接である。
ゆえに $\bigoplus \mathcal{F}_i$ は $U$ 上の準連接加群へ制限される
加群として準連接である。

\medskip\noindent
(2) の証明。図式 $\mathcal{I} \to \QCoh(\mathcal{O}_X)$,
$i \mapsto \mathcal{F}_i$ を考える。このとき上の議論により
$$
(\colim \mathcal{F}_i)|_{U_\etale} = \colim \mathcal{F}_i|_{U_\etale}
$$
であり、同じ仕方で結論する。

\medskip\noindent
(3) の証明。$a : \mathcal{F} \to \mathcal{F}'$ を
$\QCoh(\mathcal{O}_X)$ の矢印とする。このとき
$\Ker(a)|_{U_\etale} = \Ker(a|_{U_\etale})$ および
$\Coker(a)|_{U_\etale} = \Coker(a|_{U_\etale})$ であり、
同じ仕方で結論する。

\medskip\noindent
(4) の証明。制限
$0 \to \mathcal{F}_1|_{U_\etale} \to \mathcal{F}_2|_{U_\etale}
\to \mathcal{F}_3|_{U_\etale} \to 0$ は短完全である。
したがってこの列に対して三つのうち二つという性質が成り立ち、
前と同様に結論する。

\medskip\noindent
(5) の証明。$\mathcal{F}$ と $\mathcal{G}$ を
$\QCoh(\mathcal{O}_X)$ の対象とする。このとき
$$
(\mathcal{F} \otimes_{\mathcal{O}_X} \mathcal{G})_{U_\etale} =
\mathcal{F}|_{U_\etale} \otimes_{\mathcal{O}_U} \mathcal{G}|_{U_\etale}
$$
であり、前と同様に結論する。

\medskip\noindent
(6) の証明。$\mathcal{F}$ と $\mathcal{G}$ を
$\QCoh(\mathcal{O}_X)$ の対象とし、$\mathcal{F}$ は有限表示であるとする。
$$
\SheafHom_{\mathcal{O}_X}(\mathcal{F}, \mathcal{G})|_{U_\etale} =
\SheafHom_{\mathcal{O}_U}(\mathcal{F}|_{U_\etale}, \mathcal{G}|_{U_\etale})
$$
である。実際、制限は局所化であり
（Section \ref{spaces-properties-section-localize}、特に式
(\ref{spaces-properties-equation-localize-at-scheme-ringed})）を参照）、
内部 Hom の形成は局所化と可換する。
Modules on Sites, Lemma
\ref{sites-modules-lemma-internal-hom-restriction} を参照せよ。
したがって前と同様に結論する。
\end{proof}

\noindent
一般に、代数空間の射に沿う準連接層の直像が準連接であるとは限らない。
この問題には Morphisms of Spaces, Section
\ref{spaces-morphisms-section-pushforward} で戻る。

\section{加群の性質}
\label{spaces-properties-section-properties-modules}

\noindent
Modules on Sites, Sections
\ref{sites-modules-section-global},
\ref{sites-modules-section-local}, および Definition
\ref{sites-modules-definition-flat} で、任意の環付きトポス上の
$\mathcal{O}$-加群の加群について、多くの内在的性質を定義した。
$X$ が代数空間ならば、これらの概念を環付きサイト
$(X_\etale, \mathcal{O}_X)$ 上の加群、または同値なことに環付きサイト
$(X_{spaces, \etale}, \mathcal{O}_X)$ 上の加群に自由に適用する。

\medskip\noindent
大域的性質 $\mathcal{P}$：
\begin{enumerate}
\item[(a)] {\it 自由}、
\item[(b)] {\it 有限自由}、
\item[(c)] {\it 大域切断で生成される}、
\item[(d)] {\it 有限個の大域切断で生成される}、
\item[(e)] {\it 大域表示}をもつ、および
\item[(f)] {\it 大域有限表示}をもつ。
\end{enumerate}
局所的性質 $\mathcal{P}$：
\begin{enumerate}
\item[(g)] {\it 局所自由}、
\item[(f)] {\it 有限局所自由}、
\item[(h)] {\it 局所的に切断で生成される}、
\item[(i)] {\it 局所的に $r$ 個の切断で生成される}、
\item[(j)] {\it 有限型}、
\item[(k)] {\it 準連接}（Section \ref{spaces-properties-section-quasi-coherent} を参照）、
\item[(l)] {\it 有限表示}、
\item[(m)] {\it 連接}、および
\item[(n)] {\it 平坦}。
\end{enumerate}
定義から直ちに従う結果をいくつか挙げる。
\begin{enumerate}
\item $\mathcal{P}=$「連接」の場合を除き、各性質は引き戻しで保たれる。
Modules on Sites, Lemmas \ref{sites-modules-lemma-global-pullback},
\ref{sites-modules-lemma-local-pullback}, および
\ref{sites-modules-lemma-pullback-flat} を参照せよ。
\item 上の各性質（連接を含む）は、代数空間のエタール射による引き戻しで
保たれる（この場合、引き戻しは制限によって与えられる。
Lemma \ref{spaces-properties-lemma-etale-morphism-topoi} を参照）。
\item $f : Y \to X$ を代数空間の全射エタール射とする。
局所的性質 (g) -- (m) の各々について、$f^*\mathcal{F}$ が
$\mathcal{P}$ をもつならば $\mathcal{F}$ も $\mathcal{P}$ をもつ。
これは $\{Y \to X\}$ が $X_{spaces, \etale}$ の被覆であることと、
Modules on Sites, Lemma
\ref{sites-modules-lemma-local-final-object} から従う。
\item $X$ がスキーム、$\mathcal{F}$ が $X_\etale$ 上の準連接加群、
$\mathcal{P}$ が「連接」または「局所自由」以外の任意の性質ならば、
$X_\etale$ 上の $\mathcal{F}$ に対する $\mathcal{P}$ は、
$\mathcal{F}|_{X_{Zar}}$ に対する対応する性質と同値である。
すなわち、$\mathcal{F}$ をスキーム $X$ 上の準連接層とみなしたときの
$\mathcal{P}$ に対応する。Descent, Lemma
\ref{descent-lemma-equivalence-quasi-coherent-properties} を参照せよ。
\item $X$ が局所 Noether スキームで、$\mathcal{F}$ が
$X_\etale$ 上の準連接加群ならば、$\mathcal{F}$ が $X_\etale$ 上連接である
ための必要十分条件は $\mathcal{F}|_{X_{Zar}}$ が連接であることである。
すなわち、これはスキーム $X$ 上の連接層が連接であるという通常の概念に対応する。
Descent, Lemma \ref{descent-lemma-equivalence-quasi-coherent-properties} を参照せよ。
\end{enumerate}

\section{局所射影加群}
\label{spaces-properties-section-locally-projective}

\noindent
Properties, Section \ref{properties-section-locally-projective} で、
局所射影準連接加群の概念を定義したことを思い出す。

\begin{lemma}
\label{spaces-properties-lemma-locally-projective}
$S$ をスキーム、$X$ を $S$ 上の代数空間とする。
$\mathcal{F}$ を準連接 $\mathcal{O}_X$-加群とする。次は同値である。
\begin{enumerate}
\item あるスキーム $U$ と全射エタール射 $U \to X$ に対して、
制限 $\mathcal{F}|_U$ は $U$ 上局所射影である。
\item 任意のスキーム $U$ と任意のエタール射 $U \to X$ に対して、
制限 $\mathcal{F}|_U$ は $U$ 上局所射影である。
\end{enumerate}
\end{lemma}

\begin{proof}
$U \to X$ は (1) のとおりとし、$V$ がスキームであるような
エタール射 $V \to X$ を取る。このとき
$\{U \times_X V \to V\}$ はスキームの fppf 被覆である。
したがって $\mathcal{F}|_U$ が局所射影ならば、
$\mathcal{F}|_{U \times_X V}$ は局所射影であり
（Properties, Lemma
\ref{properties-lemma-locally-projective-pullback} を参照）、
Descent, Lemma \ref{descent-lemma-locally-projective-descends} により
$\mathcal{F}|_V$ も局所射影である。
\end{proof}

\begin{definition}
\label{spaces-properties-definition-locally-projective}
$S$ をスキーム、$X$ を $S$ 上の代数空間とする。
$\mathcal{F}$ を準連接 $\mathcal{O}_X$-加群とする。
Lemma \ref{spaces-properties-lemma-locally-projective} の同値な条件が満たされるとき、
$\mathcal{F}$ は {\it 局所射影} であるという。
\end{definition}

\begin{lemma}
\label{spaces-properties-lemma-locally-projective-pullback}
$S$ をスキームとする。
$f : X \to Y$ を $S$ 上の代数空間の射とする。
$\mathcal{G}$ を準連接 $\mathcal{O}_Y$-加群とする。
$\mathcal{G}$ が $Y$ 上局所射影ならば、$f^*\mathcal{G}$ は
$X$ 上局所射影である。
\end{lemma}

\begin{proof}
$V$ がスキームであるような全射エタール射 $V \to Y$ を選ぶ。
$U$ がスキームであるような全射エタール射
$U \to V \times_Y X$ を選ぶ。誘導される射を $\psi : U \to V$ と書く。
このとき
$$
f^*\mathcal{G}|_U = \psi^*(\mathcal{G}|_V)
$$
である。したがって定義とスキームの場合の結果
Properties, Lemma \ref{properties-lemma-locally-projective-pullback} から
補題が従う。
\end{proof}

\section{準連接層と表示}
\label{spaces-properties-section-quasi-coherent-presentation}

\noindent
$S$ をスキーム、$X$ を $S$ 上の代数空間とする。
$\varphi : U \to X$ を任意の全射エタール射とし、
$X = U/R$ をそれから得られる $X$ の表示とする。
Spaces, Definition \ref{spaces-definition-presentation} を参照せよ。
特に、$j = (t, s) : R \to U \times_S U$ となる亜群
$(U, R, s, t, c)$ を得る。
Groupoids, Lemma \ref{groupoids-lemma-equivalence-groupoid} を参照せよ。
Groupoids, Definition \ref{groupoids-definition-groupoid-module} で、
任意の亜群上の準連接層の概念を定義した。
これらの概念を用いると、次が得られる。

\begin{proposition}
\label{spaces-properties-proposition-quasi-coherent}
$S$, $\varphi : U \to X$, $(U, R, s, t, c)$ は上のとおりとする。
任意の準連接 $\mathcal{O}_X$-加群 $\mathcal{F}$ に対して、層
$\varphi^*\mathcal{F}$ には標準的な同型
$$
\alpha : t^*\varphi^*\mathcal{F} \longrightarrow s^*\varphi^*\mathcal{F}
$$
が備わる。これは Groupoids, Definition
\ref{groupoids-definition-groupoid-module} の条件を満たし、したがって
$(U, R, s, t, c)$ 上の準連接層を定める。関手
$\mathcal{F} \mapsto (\varphi^*\mathcal{F}, \alpha)$ は圏の同値
$$
\begin{matrix}
\text{準連接} \\
\mathcal{O}_X\text{-加群}
\end{matrix}
\longleftrightarrow
\begin{matrix}
\text{準連接加群}\\
\text{上の }(U, R, s, t, c)
\end{matrix}
$$
を定める。
\end{proposition}

\begin{proof}
命題の主張とこの証明では、スキーム上の準連接層を、そのスキームの
小エタールサイト上の準連接層と考える。これは Descent, Sections
\ref{descent-section-quasi-coherent-sheaves},
\ref{descent-section-quasi-coherent-cohomology}, および
\ref{descent-section-quasi-coherent-sheaves-bis} の結果により許される。

\medskip\noindent
$\alpha$ の存在は $\varphi \circ t = \varphi \circ s$ と、引き戻しが射に関して
関手的であることから従う。Equation (\ref{spaces-properties-equation-push-pull}) の周辺の
議論を参照せよ。まったく同様に、すなわち引き戻しの関手性により、
同型 $\alpha$ が Groupoids, Definition
\ref{groupoids-definition-groupoid-module} の条件 (1) を満たすことが分かる。
定義の条件 (2) を見るには $\alpha$ が同型であることを見れば十分であり、
これは明らかである。構成
$\mathcal{F} \mapsto (\varphi^*\mathcal{F}, \alpha)$ は
準連接層 $\mathcal{F}$ に関して明らかに関手的である。
したがって補題の表示式の左から右への関手を得る。

\medskip\noindent
逆に、$(\mathcal{F}, \alpha)$ が $(U, R, s, t, c)$ 上の準連接層であるとする。
$V \to X$ を $X_\etale$ の対象とする。この場合、射
$V' = U \times_X V \to V$ はスキームの全射エタール射なので、
$\{V' \to V\}$ は $V$ のエタール被覆である。
さらに準連接層 $\mathcal{F}$ は $V'$ 上の準連接層
$\mathcal{F}'$ へ引き戻される。
$R = U \times_X U$ かつ $t = \text{pr}_0$, $s = \text{pr}_0$ なので、
$V' \times_V V' = R \times_X V$ であり、射影
$V' \times_V V' \to V'$ は $t$ と $s$ の引き戻しに等しい。
したがって $\alpha$ は同型
$\alpha' : \text{pr}_0^*\mathcal{F}' \to \text{pr}_1^*\mathcal{F}'$ に
引き戻され、対 $(\mathcal{F}', \alpha')$ は被覆 $\{V' \to V\}$ に関する
準連接層の降下データである。
Descent, Proposition
\ref{descent-proposition-fpqc-descent-quasi-coherent} により、この降下データは
有効であり、$V_\etale$ 上の準連接 $\mathcal{O}_V$-加群
$\mathcal{F}_V$ を得る。これが $X_\etale$ 上の準連接層を与えることを示すには、
Lemma \ref{spaces-properties-lemma-characterize-quasi-coherent-small-etale} により、
$X_\etale$ の任意の射 $f : V_1 \to V_2$ に対して、合成と両立する
標準的な同型 $c_f : \mathcal{F}_{V_1} \to \mathcal{F}_{V_2}$ が
存在することを示さなければならない。確認は省略する。
また、これが右の圏から左の圏への関手を定め、上で記述した関手の逆になることの
確認も省略する。
\end{proof}

\begin{proposition}
\label{spaces-properties-proposition-coherator}
$S$ をスキーム、$X$ を $S$ 上の代数空間とする。
\begin{enumerate}
\item 圏 $\QCoh(\mathcal{O}_X)$ は Grothendieck アーベル圏である。
したがって $\QCoh(\mathcal{O}_X)$ は十分多くの単射対象とすべての極限をもつ。
\item 包含関手
$\QCoh(\mathcal{O}_X) \to \textit{Mod}(\mathcal{O}_X)$ は右随伴\footnote{
この関手は {\it coherator} と呼ばれることがある。}
$$
Q : \textit{Mod}(\mathcal{O}_X) \longrightarrow \QCoh(\mathcal{O}_X)
$$
をもち、すべての準連接層 $\mathcal{F}$ に対して随伴写像
$Q(\mathcal{F}) \to \mathcal{F}$ は同型である。
\end{enumerate}
\end{proposition}

\begin{proof}
この証明はスキームの場合の証明の繰り返しである。
Properties, Proposition \ref{properties-proposition-coherator} を参照せよ。
読者にはまずその証明を読むことを勧める。

\medskip\noindent
(1) は、$\QCoh(\mathcal{O}_X)$ が (a) すべての余極限をもち、
(b) フィルター余極限が完全であり、(c) 生成対象をもつことを意味する。
Injectives, Section \ref{injectives-section-grothendieck-conditions} を参照せよ。
Lemma \ref{spaces-properties-lemma-properties-quasi-coherent} により、
$\QCoh(\mathcal{O}_X)$ の余極限は存在し、
$\textit{Mod}(\mathcal{O}_X)$ における余極限と一致する。
Modules on Sites, Lemma \ref{sites-modules-lemma-limits-colimits} により
フィルター余極限は完全である。したがって (a) と (b) が成り立つ。

\medskip\noindent
生成対象を構成するため、$(U, R, s, t, c)$ がエタール亜群スキームとなり、
特に $s$ と $t$ がスキームの平坦射となるような表示 $X = U/R$ を選ぶ。
Groupoids, Lemma \ref{groupoids-lemma-colimit-kappa} のような濃度
$\kappa$ を選ぶ。Groupoids, Lemma
\ref{groupoids-lemma-set-of-iso-classes} のように、
$(U, R, s, t, c)$ 上の $\kappa$-生成準連接加群の族
$(\mathcal{E}_t, \alpha_t)_{t \in T}$ を選ぶ。
Proposition \ref{spaces-properties-proposition-quasi-coherent} の圏同値を介して準連接加群
$(\mathcal{E}_t, \alpha_t)$ に対応する $X$ 上の準連接加群を
$\mathcal{F}_t$ とする。このとき、任意の準連接加群 $\mathcal{H}$ は、
$\mathcal{F}_t$ のいずれかと同型なその準連接部分加群の有向余極限である。
したがって $\bigoplus_t \mathcal{F}_t$ は
$\QCoh(\mathcal{O}_X)$ の生成対象であり、(c) が成り立つと結論する。
極限と単射対象に関する主張は任意の Grothendieck アーベル圏で成り立つ。
Injectives, Theorem
\ref{injectives-theorem-injective-embedding-grothendieck} および
Lemma \ref{injectives-lemma-grothendieck-products} を参照せよ。

\medskip\noindent
(2) の証明。$Q$ を構成するため、次の一般的手順を用いる。
$\textit{Mod}(\mathcal{O}_X)$ の対象 $\mathcal{F}$ が与えられたとき、関手
$$
\QCoh(\mathcal{O}_X)^{opp} \longrightarrow \textit{Sets},\quad
\mathcal{G} \longmapsto \Hom_X(\mathcal{G}, \mathcal{F})
$$
を考える。この関手は余極限を極限へ移すので表現可能である。
Injectives, Lemma \ref{injectives-lemma-grothendieck-brown} を参照せよ。
したがって準連接層 $Q(\mathcal{F})$ と、
$\QCoh(\mathcal{O}_X)$ の $\mathcal{G}$ に対する関手的な同型
$\Hom_X(\mathcal{G}, \mathcal{F}) =
\Hom_X(\mathcal{G}, Q(\mathcal{F}))$ が存在する。
Yoneda の補題（Categories, Lemma \ref{categories-lemma-yoneda}）により、
構成 $\mathcal{F} \leadsto Q(\mathcal{F})$ は $\mathcal{F}$ に関して関手的である。
構成により $Q$ は包含関手の右随伴である。
$\mathcal{F}$ が準連接であるとき $Q(\mathcal{F}) \to \mathcal{F}$ が
同型であることは、包含関手
$\QCoh(\mathcal{O}_X) \to \textit{Mod}(\mathcal{O}_X)$ が
充満忠実であることの形式的な帰結である。
\end{proof}

\section{スキームへの射}
\label{spaces-properties-section-morphisms-to-schemes}

\noindent
次は Schemes, Lemma \ref{schemes-lemma-morphism-into-affine} の類似物である。

\begin{lemma}
\label{spaces-properties-lemma-morphism-to-affine-scheme}
$X$ を $\mathbf{Z}$ 上の代数空間とする。
$T$ をアフィンスキームとする。
$f$ を（大域切断上の）$f^\sharp$ へ写す写像
$$
\Mor(X, T)
\longrightarrow
\Hom(\Gamma(T, \mathcal{O}_T), \Gamma(X, \mathcal{O}_X))
$$
は全単射である。
\end{lemma}

\begin{proof}
写像の逆を構成する。$\varphi : \Gamma(T, \mathcal{O}_T) \to
\Gamma(X, \mathcal{O}_X)$ を環準同型とする。
$X = U/R$ をその表示として選ぶ。
Spaces, Definition \ref{spaces-definition-presentation} を参照せよ。
Schemes, Lemma \ref{schemes-lemma-morphism-into-affine} により、合成
$$
\Gamma(T, \mathcal{O}_T) \to \Gamma(X, \mathcal{O}_X) \to
\Gamma(U, \mathcal{O}_U)
$$
はスキームの一意な射 $g : U \to T$ に対応する。
同じ補題により、二つの合成 $R \to U \to T$ は等しい。
したがって $U \to X \to T$ が $g$ に等しくなるような射
$f : X = U/R \to T$ を得る。構成により図式
$$
\xymatrix{
\Gamma(U, \mathcal{O}_U) & \Gamma(X, \mathcal{O}_X) \ar[l] \\
& \Gamma(T, \mathcal{O}_T) \ar[lu]^{g^\sharp} \ar[u]^{\varphi}_{f^\sharp}
}
$$
は可換である。$U \to X$ はエタール被覆であり
$\mathcal{O}_X$ は $X_\etale$ 上の層なので、$f^\sharp$ は
$\varphi$ に等しい。$f$ の一意性は $g$ の一意性から従う。
\end{proof}

\section{自由作用による商}
\label{spaces-properties-section-quotient-by-free}

\noindent
$S$ をスキームとする。
$X$ を $S$ 上の代数空間とする。
$G$ を抽象群とする。
$a : G \to \text{Aut}(X)$ を準同型、すなわち $a$ を $X$ への $G$ の
{\it 作用} とする。$S$ 上のすべてのスキーム $T$ に対して写像
$$
G \times X(T) \longrightarrow X(T)
$$
が自由であるとき、この作用は {\it 自由} であるという。
（点がよく定まった剰余体をもたないことがあるため、
Spaces, Lemma \ref{spaces-lemma-quotient} のような判定法は使えない。）
作用が自由な場合、商 $X/G$ を代数空間として構成する。
これは後で証明する一般的な Bootstrap, Lemma
\ref{bootstrap-lemma-quotient-free-action} の特殊な場合である。

\begin{lemma}
\label{spaces-properties-lemma-quotient}
$S$ をスキームとする。
$X$ を $S$ 上の代数空間とする。
$G$ を $X$ に自由に作用する抽象群とする。
このとき商層 $X/G$ は代数空間である。
\end{lemma}

\begin{proof}
主張は、前層
$$
T \longmapsto X(T)/G
$$
に付随する層 $F$ が代数空間であることを意味する。
これを見るため表示を構成する。スキーム $U$ と全射エタール射
$\varphi : U \to X$ を選ぶ。$V = \coprod_{g \in G} U$ とおき、
$g \in G$ に対応する成分上で $\psi : V \to X$ を
$a(g) \circ \varphi$ に等しいと定める。
$G$ を成分の置換によって $V$ に作用させる。すなわち、
$g_0 \in G$ は $g$ に対応する成分を、$U$ の恒等射によって
$g_0g$ に対応する成分へ写す。このとき $\psi$ は $G$-同変射である。
したがって次の段落で扱う場合へ帰着される。

\medskip\noindent
$U$ 上に $G$-作用があり、$U \to X$ が全射、エタール、かつ
$G$-同変であると仮定する。この場合、$R = U \times_X U$ 上に、射影
$t, s : R \to U$ と両立する $G$ の誘導作用がある。ここで
$$
X/G = U/\coprod\nolimits_{g \in G} R
$$
であると主張する。写像
$$
j : \coprod\nolimits_{g \in G} R
\longrightarrow
U \times_S U
$$
は $(r, g) \mapsto (t(r), g(s(r)))$ によって与えられる。
$j$ は単射であることに注意する：
$(t(r), g(s(r))) = (t(r'), g'(s(r')))$ ならば
$t(r) = t(r')$ であり、したがって $r$ と $r'$ は $s$ と $t$ の両方で
$X$ に同じ像をもつ。したがって $G$ は $X$ に自由に作用するので
$g = g'$ であり、したがって $s(r) = s(r')$ であり、したがって
$R$ は $U$ 上の同値関係なので $r = r'$ である。
さらに $j$ は同値関係である（細部を省略する）。
$s$ と $t$ はエタールなので、二つの射影
$\coprod\nolimits_{g \in G} R \to U$ はいずれもエタールである。
したがって $j$ はエタール同値関係であり、Spaces, Theorem
\ref{spaces-theorem-presentation} により
$U/\coprod\nolimits_{g \in G} R$ は代数空間である。
写像 $U \to X$ によって誘導される写像
$$
U/\coprod\nolimits_{g \in G} R \longrightarrow X/G
$$
がある。これが層の同型であることの証明は省略する。
\end{proof}
